\documentclass[12pt,oneside]{book}

%============================================================
% PACKAGES
%============================================================

%--------------------
% Page + font
%--------------------
\usepackage[margin=1in]{geometry}
\usepackage[T1]{fontenc}
\usepackage[utf8]{inputenc}
\usepackage{lmodern}
\usepackage{microtype}

%--------------------
% Math
%--------------------
\usepackage{amsmath, amssymb, amsfonts, amsthm, mathtools}
\usepackage{bm}
\usepackage{bbm}
\usepackage{mathrsfs}
\usepackage{dsfont}

%--------------------
% Tables + figures
%--------------------
\usepackage{graphicx}
\usepackage{float}
\usepackage{booktabs}
\usepackage{array}
\usepackage{xltabular}
\usepackage{multirow}
\usepackage{caption}
\usepackage{subcaption}

%--------------------
% Lists
%--------------------
\usepackage{enumitem}
\setlist[itemize]{topsep=2pt,itemsep=2pt}
\setlist[enumerate]{topsep=2pt,itemsep=2pt}

%--------------------
% Colors + hyperlinks
%--------------------
\usepackage[dvipsnames]{xcolor}
\usepackage[
    colorlinks=true,
    linkcolor=MidnightBlue,
    citecolor=MidnightBlue,
    urlcolor=MidnightBlue
]{hyperref}
\usepackage[nameinlink,capitalise]{cleveref}

%--------------------
% Algorithms
%--------------------
\usepackage{algorithm}
\usepackage{algpseudocode}

%--------------------
% Code blocks
%--------------------
\usepackage{listings}

\lstdefinestyle{codestyle}{
    basicstyle=\ttfamily\small,
    breaklines=true,
    frame=single,
    numbers=left,
    numberstyle=\tiny,
    keywordstyle=\color{MidnightBlue},
    commentstyle=\color{ForestGreen},
    stringstyle=\color{BrickRed},
    showstringspaces=false,
    tabsize=4
}

\lstset{style=codestyle}

%--------------------
% Section styling
%--------------------
\usepackage{titlesec}
\usepackage{tocloft}

\setcounter{tocdepth}{2}
\setcounter{secnumdepth}{3}

%============================================================
% THEOREMS
%============================================================

\numberwithin{equation}{section}

\theoremstyle{definition}
\newtheorem{definition}{Definition}[section]
\newtheorem{example}{Example}[section]
\newtheorem{exercise}{Exercise}[section]
\newtheorem{remark}{Remark}[section]

\theoremstyle{plain}
\newtheorem{theorem}{Theorem}[section]
\newtheorem{lemma}{Lemma}[section]
\newtheorem{proposition}{Proposition}[section]
\newtheorem{corollary}{Corollary}[section]

\theoremstyle{remark}
\newtheorem*{note}{Note}

%============================================================
% CUSTOM COMMANDS
%============================================================

% Sets
\newcommand{\R}{\mathbb{R}}
\newcommand{\N}{\mathbb{N}}
\newcommand{\Z}{\mathbb{Z}}
\newcommand{\Q}{\mathbb{Q}}
\newcommand{\C}{\mathbb{C}}

% Probability
\newcommand{\E}{\mathbb{E}}
\newcommand{\Var}{\mathrm{Var}}
\newcommand{\Cov}{\mathrm{Cov}}
\newcommand{\Corr}{\mathrm{Corr}}
\newcommand{\Prob}{\mathbb{P}}

\newcommand{\ind}{\mathbbm{1}}
\newcommand{\iid}{\overset{\text{iid}}{\sim}}
\newcommand{\convd}{\overset{d}{\to}}
\newcommand{\convp}{\overset{p}{\to}}
\newcommand{\as}{\overset{a.s.}{\to}}

% Operators
\newcommand{\given}{\,\middle|\,}
\newcommand{\argmin}{\operatorname*{arg\,min}}
\newcommand{\argmax}{\operatorname*{arg\,max}}

% Linear algebra / analysis
\newcommand{\norm}[1]{\left\lVert #1 \right\rVert}
\newcommand{\abs}[1]{\left\lvert #1 \right\rvert}
\newcommand{\inner}[2]{\left\langle #1, #2 \right\rangle}

\newcommand{\dd}{\,\mathrm{d}}
\newcommand{\grad}{\nabla}
\newcommand{\Hess}{\nabla^2}

\newcommand{\tr}{\operatorname{tr}}
\newcommand{\rank}{\operatorname{rank}}
\newcommand{\diag}{\operatorname{diag}}

% Distributions
\newcommand{\Normal}{\mathcal{N}}
\newcommand{\Bern}{\mathrm{Bernoulli}}
\newcommand{\Bin}{\mathrm{Binomial}}
\newcommand{\Pois}{\mathrm{Poisson}}
\newcommand{\Gam}{\mathrm{Gamma}}
\newcommand{\BetaDist}{\mathrm{Beta}}
\newcommand{\Exp}{\mathrm{Exponential}}
\newcommand{\Dir}{\mathrm{Dirichlet}}
\newcommand{\InvGam}{\mathrm{Inv\text{-}Gamma}}

% Title block
\title{EN.553.747 Mathematics of Data Science}
\author{Jonathan Y. Ma}
\date{June 2026}

\begin{document}

\maketitle

\noindent \emph{This course was taught by Dr. Mateo Diaz in the Fall 2025 Semester. The notes are transcribed by Jonathan Ma and may contain errors and omitted content, so any issues are to be directed towards me.}

\tableofcontents

\chapter{Foundations of High-Dimensional Probability}

\section{Motivation}

The central question in modern data science is:

\begin{quote}
How can one efficiently extract information from high-dimensional data?
\end{quote}

This question is simultaneously:
\begin{itemize}
\item statistical,
\item computational,
\item geometric,
\item and optimization-theoretic.
\end{itemize}

Many classical methods fail in high dimensions due to exponential scaling in ambient dimension. The mathematics of data science studies the probabilistic and geometric structures that permit efficient inference despite this apparent complexity.

Core themes include:
\begin{itemize}
\item concentration inequalities,
\item dimension reduction,
\item randomized linear algebra,
\item compressed sensing,
\item spectral methods,
\item high-dimensional estimation,
\item convex and nonconvex optimization,
\item statistical learning theory.
\end{itemize}

A primary reference for the probabilistic aspects of the course is Roman Vershynin's \emph{High-Dimensional Probability}.  

\section{The Curse of Dimensionality}

\subsection{Numerical Integration in One Dimension}

Suppose we wish to estimate
\[
I = \int_{[0,1]} f(x)\dd x.
\]

A classical deterministic approximation is the Riemann sum:
\[
I \approx \frac{1}{n}\sum_{i=1}^n f(x_i),
\]
where $\{x_i\}$ are grid points.

\begin{definition}
A function $f:[0,1]\to\R$ is called $L$-Lipschitz if
\[
|f(x)-f(y)| \leq L|x-y|
\]
for all $x,y\in[0,1]$.
\end{definition}

\begin{proposition}
If $f$ is $L$-Lipschitz and we partition $[0,1]$ into intervals of width $\varepsilon$, then the Riemann sum approximation error is $O(\varepsilon)$.
\end{proposition}

\begin{proof}
Let the partition intervals be
\[
I_k = [(k-1)\varepsilon,k\varepsilon].
\]

For any $x,y\in I_k$,
\[
|f(x)-f(y)| \leq L\varepsilon.
\]

Thus replacing $f$ by a representative value on each interval introduces at most $L\varepsilon$ local error. Summing over all intervals preserves the same asymptotic scaling:
\[
\left|
\int_0^1 f(x)\dd x
-
\frac{1}{n}\sum_{i=1}^n f(x_i)
\right|
= O(\varepsilon).
\]
\end{proof}

\subsection{High-Dimensional Grids}

Now consider
\[
I_d = \int_{[0,1]^d} f(x)\dd x.
\]

If each coordinate is discretized with mesh width $\varepsilon$, then approximately
\[
\left(\frac{1}{\varepsilon}\right)^d
\]
grid points are required.

Thus the computational cost grows exponentially in dimension.

\begin{definition}
The phenomenon in which computational or statistical complexity grows exponentially with ambient dimension is called the \emph{curse of dimensionality}.
\end{definition}

\begin{remark}
Many deterministic approximation methods become computationally infeasible in high dimensions because even moderate accuracy requires exponentially many samples.
\end{remark}

\section{Monte Carlo Integration}

Randomization provides a fundamentally different approach.

Let
\[
X_1,\dots,X_n \iid \mathrm{Unif}([0,1]^d).
\]

Then
\[
\E[f(X_i)]
=
\int_{[0,1]^d} f(x)\dd x.
\]

Hence we may estimate the integral via
\[
\widehat{I}_n
=
\frac{1}{n}\sum_{i=1}^n f(X_i).
\]

\subsection{Strong Law of Large Numbers}

\begin{theorem}[Strong Law of Large Numbers]
Let $X_1,X_2,\dots$ be iid with $\E|X_1|<\infty$. Then
\[
\frac{1}{n}\sum_{i=1}^n X_i
\as
\E[X_1].
\]
\end{theorem}

Applying this theorem to $f(X_i)$ yields
\[
\widehat{I}_n
\as
\int_{[0,1]^d} f(x)\dd x.
\]

\subsection{Dimension-Independent Error Scaling}

A major advantage of Monte Carlo methods is that convergence rates are independent of ambient dimension.

\begin{theorem}
Assume $f(X)$ has finite variance. Then
\[
\E\left[
\left(
\frac{1}{n}\sum_{i=1}^n f(X_i)-\E[f(X)]
\right)^2
\right]
=
\frac{\Var(f(X))}{n}.
\]
\end{theorem}

\begin{proof}
Define
\[
Y_i = f(X_i)-\E[f(X)].
\]

Then $\E[Y_i]=0$ and
\[
\frac{1}{n}\sum_{i=1}^n f(X_i)-\E[f(X)]
=
\frac{1}{n}\sum_{i=1}^n Y_i.
\]

Therefore
\[
\E\left[
\left(
\frac{1}{n}\sum_{i=1}^n Y_i
\right)^2
\right]
=
\frac{1}{n^2}
\E\left[
\left(
\sum_{i=1}^n Y_i
\right)^2
\right].
\]

Expanding,
\[
=
\frac{1}{n^2}
\sum_{i,j}\E[Y_iY_j].
\]

Since the $Y_i$ are independent and centered,
\[
\E[Y_iY_j]=0
\quad\text{for } i\neq j.
\]

Thus
\[
=
\frac{1}{n^2}
\sum_{i=1}^n \E[Y_i^2]
=
\frac{\Var(f(X))}{n}.
\]
\end{proof}

Taking square roots gives
\[
\mathrm{RMSE}
=
O\left(\frac{1}{\sqrt{n}}\right).
\]

\begin{remark}
The convergence rate
\[
O(n^{-1/2})
\]
does not depend on dimension $d$.
This is one of the fundamental reasons Monte Carlo methods are effective in high-dimensional problems.
\end{remark}

\section{Probability Review}

\subsection{Expectation}

\begin{definition}
Let $X$ be a random vector with density $p(x)$ and let $f:\R^d\to\R$. The expectation of $f(X)$ is
\[
\E[f(X)]
=
\int_{\R^d} f(x)p(x)\dd x.
\]
\end{definition}

\subsection{Variance}

$\mathrm{Var}(X)=\mathbb{E}[(X-\mathbb{E}X)^2]$

\begin{definition}
For a random variable $X$,
\[
\Var(X)
=
\E[(X-\E[X])^2].
\]
\end{definition}

\begin{proposition}
Variance admits the equivalent formula
\[
\Var(X)=\E[X^2]-(\E[X])^2.
\]
\end{proposition}

\begin{proof}
Expand:
\[
(X-\E[X])^2
=
X^2-2X\E[X]+(\E[X])^2.
\]

Taking expectations,
\[
\Var(X)
=
\E[X^2]
-
2(\E[X])^2
+
(\E[X])^2.
\]

Hence
\[
\Var(X)
=
\E[X^2]-(\E[X])^2.
\]
\end{proof}

\subsection{Linearity of Expectation}

\begin{proposition}
For random variables $X_1,\dots,X_n$,
\[
\E\left[\sum_{i=1}^n X_i\right]
=
\sum_{i=1}^n \E[X_i].
\]
\end{proposition}

\begin{proof}
Expectation is an integral operator, and integration is linear.
\end{proof}

\subsection{Variance of Independent Sums}

\begin{proposition}
If $X_1,\dots,X_n$ are independent, then
\[
\Var\left(\sum_{i=1}^n X_i\right)
=
\sum_{i=1}^n \Var(X_i).
\]
\end{proposition}

\begin{proof}
We compute
\[
\Var\left(\sum_i X_i\right)
=
\E\left[
\left(
\sum_i (X_i-\E[X_i])
\right)^2
\right].
\]

Expanding,
\[
=
\sum_i \Var(X_i)
+
\sum_{i\neq j}
\Cov(X_i,X_j).
\]

Independence implies
\[
\Cov(X_i,X_j)=0,
\]
so the cross terms vanish.
\end{proof}

\subsection{Jensen's Inequality}

\begin{definition}
A function $\varphi:\R\to\R$ is convex if
\[
\varphi(\lambda x+(1-\lambda)y)
\leq
\lambda\varphi(x)+(1-\lambda)\varphi(y)
\]
for all $\lambda\in[0,1]$.
\end{definition}

\begin{theorem}[Jensen's Inequality]
If $X$ is integrable and $\varphi$ is convex, then
\[
\varphi(\E[X])
\leq
\E[\varphi(X)].
\]
\end{theorem}

\begin{example}
Taking $\varphi(x)=x^2$ gives
\[
(\E[X])^2 \leq \E[X^2].
\]
\end{example}

\section{The Central Limit Theorem}

Suppose
\[
X_1,\dots,X_n
\iid
(\mu,\sigma^2),
\]
meaning
\[
\E[X_i]=\mu,
\qquad
\Var(X_i)=\sigma^2.
\]

Define
\[
S_n=\sum_{i=1}^n X_i.
\]

The standardized sum is
\[
Z_n
=
\frac{S_n-\E[S_n]}
{\sqrt{\Var(S_n)}}.
\]

Since
\[
\E[S_n]=n\mu,
\qquad
\Var(S_n)=n\sigma^2,
\]
we obtain
\[
Z_n
=
\frac{1}{\sigma\sqrt{n}}
\sum_{i=1}^n (X_i-\mu).
\]

$Z_n=\frac{1}{\sigma\sqrt{n}}\sum_{i=1}^n (X_i-\mu)$

\begin{theorem}[Central Limit Theorem]
Let $X_1,X_2,\dots$ be iid with finite mean $\mu$ and variance $\sigma^2$. Then
\[
Z_n
\convd
\Normal(0,1).
\]
\end{theorem}

\subsection{Gaussian Tail Bounds}

Let
\[
Z\sim \Normal(0,1).
\]

A standard Gaussian tail estimate is
\[
\Prob(Z\geq t)
\leq
e^{-t^2/2}.
\]

\[
\mathbb{P}(Z\ge t)\le e^{-t^2/2}
\]

\begin{proof}
Using Markov's inequality,
\[
\Prob(Z\geq t)
=
\Prob(e^{\lambda Z}\geq e^{\lambda t})
\leq
e^{-\lambda t}\E[e^{\lambda Z}].
\]

Since
\[
\E[e^{\lambda Z}]
=
e^{\lambda^2/2},
\]
we get
\[
\Prob(Z\geq t)
\leq
\exp\left(
\frac{\lambda^2}{2}-\lambda t
\right).
\]

Optimizing over $\lambda$ gives $\lambda=t$, yielding
\[
\Prob(Z\geq t)
\leq
e^{-t^2/2}.
\]
\end{proof}

\section{A Coin-Flipping Example}

Suppose
\[
C_i\in\{-1,+1\}
\]
represents the outcome of the $i$-th coin flip.

For a fair coin,
\[
\E[C_i]=0.
\]

We may test fairness using the empirical mean
\[
\frac{1}{n}\sum_{i=1}^n C_i.
\]

The CLT implies
\[
\sqrt{n}
\left(
\frac{1}{n}\sum_{i=1}^n C_i
\right)
\convd
\Normal(0,1).
\]

Hence,
\[
\Prob\left(
\frac{1}{n}\sum_{i=1}^n C_i
\geq \frac{3}{4}
\right)
=
\Prob\left(
Z_n \geq \frac{\sqrt{n}}{2}
\right).
\]

Using the Gaussian tail bound,
\[
\Prob\left(
Z_n \geq \frac{\sqrt{n}}{2}
\right)
\lesssim
e^{-n/8}.
\]

\begin{remark}
This illustrates a key theme of high-dimensional probability:

\begin{quote}
Random fluctuations often decay exponentially fast.
\end{quote}

Understanding such concentration phenomena is one of the primary goals of the course.
\end{remark}

\section{Course Themes}

Major themes in the mathematics of data science include:
\begin{itemize}
\item concentration inequalities,
\item dimension reduction,
\item clustering,
\item foundations of statistical inference,
\item statistical learning theory,
\item minimax lower bounds,
\item randomized algorithms,
\item high-dimensional estimation,
\item convex and nonconvex optimization.
\end{itemize}

\chapter{Concentration Inequalities and Sub-Gaussian Random Variables}

\section{Tail Bounds and the Chernoff Method}

A central problem in high-dimensional probability is understanding the concentration behavior of random sums.

Given independent random variables
\[
X_1,\dots,X_n,
\]
we seek bounds of the form
\[
\Prob\left(
\left|
\sum_{i=1}^n X_i
\right|
\geq t
\right).
\]

In many situations, the goal is to establish exponentially decaying tail bounds.

\section{Markov's Inequality}

\begin{theorem}[Markov's Inequality]
Let $X$ be a nonnegative random variable and let $t>0$. Then
\[
\Prob(X\geq t)
\leq
\frac{\E[X]}{t}.
\]
\end{theorem}

\begin{proof}
Observe that
\[
X
=
X\ind_{\{X\geq t\}}
+
X\ind_{\{X<t\}}.
\]

Since $X\geq t$ on the event $\{X\geq t\}$,
\[
X\ind_{\{X\geq t\}}
\geq
t\ind_{\{X\geq t\}}.
\]

Taking expectations gives
\[
\E[X]
\geq
t\E[\ind_{\{X\geq t\}}]
=
t\Prob(X\geq t).
\]

Rearranging yields
\[
\Prob(X\geq t)
\leq
\frac{\E[X]}{t}.
\]
\end{proof}

\section{Chebyshev's Inequality}

Applying Markov's inequality to a squared random variable yields a variance-based concentration bound.

\begin{theorem}[Chebyshev's Inequality]
For any random variable $X$ with finite variance,
\[
\Prob(|X-\E[X]|\geq t)
\leq
\frac{\Var(X)}{t^2}.
\]
\end{theorem}

\begin{proof}
Apply Markov's inequality to
\[
Y=(X-\E[X])^2.
\]

Then
\[
\Prob(|X-\E[X]|\geq t)
=
\Prob((X-\E[X])^2\geq t^2)
\leq
\frac{\E[(X-\E[X])^2]}{t^2}.
\]

Since
\[
\E[(X-\E[X])^2]
=
\Var(X),
\]
the result follows.
\end{proof}

\begin{remark}
Chebyshev's inequality gives polynomial tail decay:
\[
\Prob(|X-\E[X]|\geq t)
=
O(t^{-2}).
\]

In high-dimensional settings, this is often too weak. We seek exponential concentration inequalities instead.
\end{remark}

\section{The Chernoff Method}

The key idea behind Chernoff bounds is to transform tail events using monotone functions.

For any increasing function
\[
f:\R\to\R_+,
\]
we have
\[
\Prob(X-\E[X]\geq t)
=
\Prob(f(X-\E[X])\geq f(t)).
\]

A particularly useful choice is the exponential function:
\[
f(x)=e^{\lambda x},
\]
where $\lambda>0$ is a tunable parameter.

Thus,
\[
\Prob(X-\E[X]\geq t)
=
\Prob(e^{\lambda(X-\E[X])}\geq e^{\lambda t}).
\]

Applying Markov's inequality gives
\[
\Prob(X-\E[X]\geq t)
\leq
\frac{\E[e^{\lambda(X-\E[X])}]}{e^{\lambda t}}.
\]

Equivalently,
\[
\Prob(X-\E[X]\geq t)
\leq
\exp\left(
\log\E[e^{\lambda(X-\E[X])}]
-
\lambda t
\right).
\]

Optimizing over $\lambda\geq0$ yields
\[
\log\Prob(X-\E[X]\geq t)
\leq
\inf_{\lambda\geq0}
\left\{
\log\E[e^{\lambda(X-\E[X])}]
-
\lambda t
\right\}.
\]

\section{Fenchel Conjugates}

The optimization expression above naturally leads to convex duality.

\begin{definition}[Fenchel Conjugate]
Let
\[
\psi:\R\to\R\cup\{\infty\}.
\]

The Fenchel conjugate of $\psi$ is
\[
\psi^*(t)
=
\sup_{\lambda\in\R}
\{
\lambda t-\psi(\lambda)
\}.
\]
\end{definition}

Define the log moment generating function
\[
\psi_X(\lambda)
=
\log\E[e^{\lambda(X-\E[X])}].
\]

Then the Chernoff argument becomes
\[
\Prob(X-\E[X]\geq t)
\leq
\exp(-\psi_X^*(t)).
\]

\begin{theorem}[Chernoff Bound]
For any random variable $X$,
\[
\Prob(X-\E[X]\geq t)
\leq
\exp(-\psi_X^*(t)).
\]
\end{theorem}

\begin{proof}
Starting from
\[
\Prob(X-\E[X]\geq t)
\leq
\exp(
\psi_X(\lambda)-\lambda t
),
\]
we optimize over $\lambda\geq0$:
\[
\log\Prob(X-\E[X]\geq t)
\leq
\inf_{\lambda\geq0}
\{
\psi_X(\lambda)-\lambda t
\}.
\]

Since
\[
\inf_{\lambda}
\{
\psi_X(\lambda)-\lambda t
\}
=
-
\sup_{\lambda}
\{
\lambda t-\psi_X(\lambda)
\},
\]
we obtain
\[
\log\Prob(X-\E[X]\geq t)
\leq
-\psi_X^*(t).
\]

Exponentiating yields
\[
\Prob(X-\E[X]\geq t)
\leq
e^{-\psi_X^*(t)}.
\]
\end{proof}

\section{Gaussian Concentration}

Suppose
\[
X\sim \Normal(\mu,\sigma^2).
\]

Its centered moment generating function is
\[
\E[e^{\lambda(X-\mu)}]
=
e^{\sigma^2\lambda^2/2}.
\]

$\mathbb{E}e^{\lambda(X-\mu)}=e^{\sigma^2\lambda^2/2}$

Thus
\[
\psi_X(\lambda)
=
\frac{\sigma^2\lambda^2}{2}.
\]

The Fenchel conjugate becomes
\[
\psi_X^*(t)
=
\sup_{\lambda}
\left\{
\lambda t-\frac{\sigma^2\lambda^2}{2}
\right\}.
\]

This is a quadratic maximization problem.

Differentiating,
\[
t-\sigma^2\lambda=0
\quad\Rightarrow\quad
\lambda^*=\frac{t}{\sigma^2}.
\]

Substituting,
\[
\psi_X^*(t)
=
\frac{t^2}{2\sigma^2}.
\]

Hence,
\[
\Prob(X-\mu\geq t)
\leq
e^{-t^2/(2\sigma^2)}.
\]

$\mathbb{P}(X-\mu\ge t)\le e^{-t^2/(2\sigma^2)}$

\begin{remark}
Gaussian random variables exhibit exponential concentration with quadratic decay in $t$.
\end{remark}

\section{Sub-Gaussian Random Variables}

The Gaussian concentration inequality motivates a broader class of random variables.

\begin{definition}[Sub-Gaussian Random Variable]
A random variable $X$ with mean
\[
\mu=\E[X]
\]
is called $\sigma$-sub-Gaussian if
\[
\log\E[e^{\lambda(X-\mu)}]
\leq
\frac{\sigma^2\lambda^2}{2}
\]
for all $\lambda\in\R$.
\end{definition}

\begin{remark}
Sub-Gaussian random variables have tails no heavier than those of a Gaussian random variable.
\end{remark}

\subsection{Tail Bounds for Sub-Gaussian Variables}

\begin{theorem}
If $X$ is $\sigma$-sub-Gaussian, then
\[
\Prob(X-\mu\geq t)
\leq
e^{-t^2/(2\sigma^2)}.
\]
\end{theorem}

\begin{proof}
By sub-Gaussianity,
\[
\psi_X(\lambda)
\leq
\frac{\sigma^2\lambda^2}{2}.
\]

Therefore,
\[
\psi_X^*(t)
=
\sup_\lambda
\{
\lambda t-\psi_X(\lambda)
\}
\geq
\sup_\lambda
\left\{
\lambda t-\frac{\sigma^2\lambda^2}{2}
\right\}
=
\frac{t^2}{2\sigma^2}.
\]

Applying the Chernoff bound,
\[
\Prob(X-\mu\geq t)
\leq
e^{-t^2/(2\sigma^2)}.
\]
\end{proof}

\begin{proposition}
If $X$ is $\sigma$-sub-Gaussian, then
\[
-X
\]
is also $\sigma$-sub-Gaussian.
\end{proposition}

\begin{proof}
Observe that
\[
\log\E[e^{\lambda(-X+\mu)}]
=
\log\E[e^{-\lambda(X-\mu)}].
\]

Applying sub-Gaussianity with parameter $-\lambda$ gives
\[
\log\E[e^{-\lambda(X-\mu)}]
\leq
\frac{\sigma^2\lambda^2}{2}.
\]
\end{proof}

\begin{corollary}
If $X$ is $\sigma$-sub-Gaussian, then
\[
\Prob(|X-\mu|\geq t)
\leq
2e^{-t^2/(2\sigma^2)}.
\]
\end{corollary}

\begin{proof}
Using the union bound,
\[
\Prob(|X-\mu|\geq t)
=
\Prob(X-\mu\geq t)
+
\Prob(\mu-X\geq t).
\]

Each term is bounded by
\[
e^{-t^2/(2\sigma^2)}.
\]
\end{proof}

\section{Bounded Random Variables are Sub-Gaussian}

\begin{theorem}
Suppose $X$ is supported on $[a,b]$. Then $X$ is
\[
\frac{b-a}{2}
\]
-sub-Gaussian.
\end{theorem}

\begin{proof}
Let
\[
Y=X-\E[X].
\]

Define
\[
\psi_Y(\lambda)
=
\log\E[e^{\lambda Y}].
\]

Differentiating,
\[
\psi_Y'(\lambda)
=
\frac{\E[Ye^{\lambda Y}]}
{\E[e^{\lambda Y}]},
\]
and
\[
\psi_Y''(\lambda)
=
\frac{\E[Y^2e^{\lambda Y}]}
{\E[e^{\lambda Y}]}
-
\left(
\frac{\E[Ye^{\lambda Y}]}
{\E[e^{\lambda Y}]}
\right)^2.
\]

Define the tilted density
\[
f_\lambda(y)
=
\frac{e^{\lambda y}f_Y(y)}
{\E[e^{\lambda Y}]}.
\]

Then
\[
\psi_Y''(\lambda)
=
\Var_\lambda(Y),
\]
where the variance is computed under the tilted distribution.

Since $Y$ is supported on an interval of length $b-a$,
\[
\Var_\lambda(Y)
\leq
\frac{(b-a)^2}{4}.
\]

By Taylor's theorem,
\[
\psi_Y(\lambda)
=
\psi_Y(0)
+
\psi_Y'(0)\lambda
+
\frac{\psi_Y''(\xi)}{2}\lambda^2
\]
for some $\xi$ between $0$ and $\lambda$.

Since
\[
\psi_Y(0)=0,
\qquad
\psi_Y'(0)=\E[Y]=0,
\]
we obtain
\[
\psi_Y(\lambda)
\leq
\frac{(b-a)^2}{8}\lambda^2.
\]

Thus $X$ is
\[
\frac{b-a}{2}
\]
-sub-Gaussian.
\end{proof}

\section{Sums of Sub-Gaussian Random Variables}

\begin{lemma}[Sum Rule]
Suppose
\[
X_1,\dots,X_n
\]
are independent and $X_i$ is $\sigma_i$-sub-Gaussian.

Then
\[
\sum_{i=1}^n X_i
\]
is
\[
\sqrt{\sum_{i=1}^n \sigma_i^2}
\]
-sub-Gaussian.
\end{lemma}

\begin{proof}
Using independence,
\[
\E\left[
e^{\lambda\sum_i(X_i-\E[X_i])}
\right]
=
\prod_{i=1}^n
\E[e^{\lambda(X_i-\E[X_i])}].
\]

Taking logarithms,
\[
\log\E\left[
e^{\lambda\sum_i(X_i-\E[X_i])}
\right]
=
\sum_{i=1}^n
\log\E[e^{\lambda(X_i-\E[X_i])}].
\]

Applying sub-Gaussianity,
\[
\leq
\sum_{i=1}^n
\frac{\sigma_i^2\lambda^2}{2}
=
\frac{\lambda^2}{2}
\sum_{i=1}^n \sigma_i^2.
\]
\end{proof}

\section{Hoeffding's Inequality}

\begin{theorem}[Hoeffding's Inequality]
Suppose
\[
X_1,\dots,X_n
\]
are independent with
\[
\E[X_i]=\mu_i,
\]
and each $X_i$ is $\sigma_i$-sub-Gaussian.

Then
\[
\Prob\left(
\sum_{i=1}^n (X_i-\mu_i)
\geq
t\norm{\sigma}_2
\right)
\leq
e^{-t^2/2},
\]
where
\[
\norm{\sigma}_2^2
=
\sum_{i=1}^n \sigma_i^2.
\]
\end{theorem}

$\mathbb{P}\left(\sum_{i=1}^n (X_i-\mu_i)\ge t\|\sigma\|_2\right)\le e^{-t^2/2}$

\begin{proof}
By the sum rule,
\[
\sum_{i=1}^n (X_i-\mu_i)
\]
is $\norm{\sigma}_2$-sub-Gaussian.

Applying the sub-Gaussian tail bound gives
\[
\Prob\left(
\sum_{i=1}^n (X_i-\mu_i)
\geq
t\norm{\sigma}_2
\right)
\leq
e^{-t^2/2}.
\]
\end{proof}

\section{Rademacher Random Variables}

\begin{definition}
A Rademacher random variable is a random variable $C$ such that
\[
\Prob(C=1)=\Prob(C=-1)=\frac12.
\]
\end{definition}

Since $C\in[-1,1]$, the bounded random variable theorem implies that $C$ is $1$-sub-Gaussian.

Thus, if
\[
C_1,\dots,C_n
\]
are iid Rademacher random variables, Hoeffding's inequality yields
\[
\Prob\left(
\frac{1}{n}\sum_{i=1}^n C_i
\geq
\frac12
\right)
\leq
e^{-n/8}.
\]

\begin{remark}
This shows exponential concentration of empirical averages around their mean.
\end{remark}

\section{Robust Mean Estimation}

Hoeffding's inequality gives confidence intervals for the sample mean when the observations are sub-Gaussian. If
\[
X_1,\dots,X_n \iid X
\]
are $\sigma$-sub-Gaussian with mean $\mu=\E[X]$, then the sample mean
\[
\widehat{\mu}
=
\frac{1}{n}\sum_{i=1}^n X_i
\]
satisfies
\[
\Prob\left(
|\widehat{\mu}-\mu|
\leq
\sigma\sqrt{\frac{2\log(2/p)}{n}}
\right)
\geq
1-p.
\]

This raises a natural question: can one obtain a similar high-probability guarantee without assuming sub-Gaussian tails?

The answer is yes, provided that the random variables have finite variance. The key idea is to replace the ordinary sample mean with a more robust estimator that ignores the effect of extreme outliers.

\subsection{The Median-of-Means Estimator}

Suppose
\[
X_1,\dots,X_n
\]
are iid with
\[
\E[X_i]=\mu,
\qquad
\Var(X_i)=\sigma^2.
\]

Fix $p\in(0,1]$. Split the samples into
\[
K=\left\lceil 18\log(1/p)\right\rceil
\]
equally sized bins. For simplicity, assume that $K$ divides $n$, and let
\[
m=\frac{n}{K}
\]
be the number of observations in each bin.

For each bin $j=1,\dots,K$, define the bin mean
\[
\widehat{\mu}_j
=
\frac{1}{m}
\sum_{i\in B_j} X_i.
\]

The median-of-means estimator is
\[
\widehat{\mu}_{\mathrm{MOM}}
=
\operatorname{median}
\left(
\widehat{\mu}_1,\dots,\widehat{\mu}_K
\right).
\]

\begin{theorem}[Median-of-Means Bound]
Let $X_1,\dots,X_n$ be iid with mean $\mu$ and variance $\sigma^2$. Let
\[
K=\left\lceil 18\log(1/p)\right\rceil.
\]
Then the median-of-means estimator satisfies
\[
\Prob\left(
|\widehat{\mu}_{\mathrm{MOM}}-\mu|
\leq
\sigma\sqrt{\frac{3K}{n}}
\right)
\geq
1-p.
\]
Equivalently, since $K\asymp \log(1/p)$,
\[
|\widehat{\mu}_{\mathrm{MOM}}-\mu|
=
O\left(
\sigma\sqrt{\frac{\log(1/p)}{n}}
\right)
\]
with probability at least $1-p$.
\end{theorem}

\begin{proof}
Let
\[
t
=
\sigma\sqrt{\frac{3K}{n}}.
\]

For each bin $j$, define the event
\[
E_j
=
\left\{
|\widehat{\mu}_j-\mu|\leq t
\right\}.
\]

Each bin mean is an average of $m=n/K$ iid observations, so
\[
\E[\widehat{\mu}_j]=\mu
\]
and
\[
\Var(\widehat{\mu}_j)
=
\frac{\sigma^2}{m}
=
\frac{\sigma^2K}{n}.
\]

By Chebyshev's inequality,
\[
\Prob(E_j^c)
=
\Prob\left(
|\widehat{\mu}_j-\mu|>t
\right)
\leq
\frac{\Var(\widehat{\mu}_j)}{t^2}.
\]

Substituting the variance and the definition of $t$ gives
\[
\Prob(E_j^c)
\leq
\frac{\sigma^2K/n}{3\sigma^2K/n}
=
\frac{1}{3}.
\]

Therefore,
\[
\Prob(E_j)\geq \frac{2}{3}.
\]

Let
\[
Z_j=\ind_{E_j}.
\]

Then $Z_1,\dots,Z_K$ are independent Bernoulli random variables with
\[
\E[Z_j]\geq \frac{2}{3}.
\]

If more than half of the events $E_j$ occur, then at least half of the bin means lie in the interval
\[
[\mu-t,\mu+t].
\]

Therefore the median of the bin means must also lie in this interval. Hence,
\[
|\widehat{\mu}_{\mathrm{MOM}}-\mu|\leq t
\]
whenever
\[
\frac{1}{K}\sum_{j=1}^K Z_j>\frac{1}{2}.
\]

Thus,
\[
\Prob\left(
|\widehat{\mu}_{\mathrm{MOM}}-\mu|>t
\right)
\leq
\Prob\left(
\frac{1}{K}\sum_{j=1}^K Z_j\leq \frac{1}{2}
\right).
\]

Since $\E[Z_j]\geq 2/3$, the event on the right implies
\[
\left|
\frac{1}{K}\sum_{j=1}^K Z_j
-
\E[Z_j]
\right|
\geq
\frac{1}{6}.
\]

The variables $Z_j$ are bounded in $[0,1]$, so Hoeffding's inequality gives
\[
\Prob\left(
\left|
\frac{1}{K}\sum_{j=1}^K Z_j
-
\E[Z_j]
\right|
\geq
\frac{1}{6}
\right)
\leq
2\exp\left(-\frac{K}{18}\right).
\]

With
\[
K=\left\lceil 18\log(1/p)\right\rceil,
\]
this probability is at most a constant multiple of $p$. Up to constants in the choice of $K$, we obtain
\[
\Prob\left(
|\widehat{\mu}_{\mathrm{MOM}}-\mu|>t
\right)
\leq
p.
\]

Therefore,
\[
\Prob\left(
|\widehat{\mu}_{\mathrm{MOM}}-\mu|
\leq
\sigma\sqrt{\frac{3K}{n}}
\right)
\geq
1-p.
\]
\end{proof}

\begin{remark}
The ordinary sample mean is sensitive to outliers when the distribution has heavy tails. The median-of-means estimator is robust because a few corrupted or extreme observations can only damage the bins in which they appear. As long as a majority of bin means remain accurate, the median remains accurate.
\end{remark}

\begin{remark}
The median-of-means estimator only requires a finite second moment. In particular, the moment generating function need not exist.
\end{remark}

\section{Chernoff's Inequality for Bernoulli Sums}

Hoeffding's inequality treats bounded random variables using only their worst-case range. For Bernoulli random variables, this can be loose.

For example, if
\[
X_i\sim \Bern(p),
\]
then $X_i\in[0,1]$, so Hoeffding gives a bound depending only on the interval length $1$. However, as $p\to 0$, the variable is usually zero, so the sum should concentrate more strongly near its mean. Chernoff's inequality captures this dependence on the mean.

\subsection{Upper Tail Bound}

\begin{theorem}[Chernoff's Inequality]
Let
\[
X_1,\dots,X_n
\]
be independent Bernoulli random variables with
\[
X_i\sim \Bern(p_i).
\]

Define
\[
S_n=\sum_{i=1}^n X_i
\]
and
\[
\mu=\E[S_n]=\sum_{i=1}^n p_i.
\]

Then, for all $t\geq \mu$,
\[
\Prob(S_n\geq t)
\leq
e^{-\mu}
\left(
\frac{e\mu}{t}
\right)^t.
\]
\end{theorem}

\begin{proof}
Let $\lambda>0$. Then
\[
\Prob(S_n\geq t)
=
\Prob(e^{\lambda S_n}\geq e^{\lambda t}).
\]

By Markov's inequality,
\[
\Prob(S_n\geq t)
\leq
e^{-\lambda t}\E[e^{\lambda S_n}].
\]

Using independence,
\[
\E[e^{\lambda S_n}]
=
\E\left[
e^{\lambda\sum_{i=1}^n X_i}
\right]
=
\prod_{i=1}^n
\E[e^{\lambda X_i}].
\]

Since $X_i\sim\Bern(p_i)$,
\[
\E[e^{\lambda X_i}]
=
p_i e^\lambda+(1-p_i)
=
1+p_i(e^\lambda-1).
\]

Using
\[
1+x\leq e^x,
\]
we obtain
\[
\E[e^{\lambda X_i}]
\leq
\exp\left(p_i(e^\lambda-1)\right).
\]

Therefore,
\[
\E[e^{\lambda S_n}]
\leq
\prod_{i=1}^n
\exp\left(p_i(e^\lambda-1)\right)
=
\exp\left((e^\lambda-1)\sum_{i=1}^n p_i\right).
\]

Since $\sum_i p_i=\mu$,
\[
\Prob(S_n\geq t)
\leq
\exp\left(
-\lambda t+(e^\lambda-1)\mu
\right).
\]

We minimize the exponent over $\lambda>0$. Define
\[
\phi(\lambda)
=
-\lambda t+(e^\lambda-1)\mu.
\]

Then
\[
\phi'(\lambda)
=
-t+\mu e^\lambda.
\]

Setting $\phi'(\lambda)=0$ gives
\[
e^\lambda=\frac{t}{\mu},
\qquad
\lambda=\log\left(\frac{t}{\mu}\right).
\]

This choice is valid when $t\geq\mu$. Substituting,
\[
\phi(\lambda)
=
-t\log\left(\frac{t}{\mu}\right)
+
t-\mu.
\]

Hence,
\[
\Prob(S_n\geq t)
\leq
\exp\left(
-t\log\left(\frac{t}{\mu}\right)+t-\mu
\right).
\]

Equivalently,
\[
\Prob(S_n\geq t)
\leq
e^{-\mu}
\left(
\frac{e\mu}{t}
\right)^t.
\]
\end{proof}

\subsection{Lower Tail Bound}

\begin{theorem}[Lower-Tail Chernoff Bound]
Let $S_n$ and $\mu$ be as above. Then, for all $0\leq t\leq\mu$,
\[
\Prob(S_n\leq t)
\leq
e^{-\mu}
\left(
\frac{e\mu}{t}
\right)^t.
\]
\end{theorem}

\begin{proof}
The argument is analogous to the upper-tail proof, but one applies Markov's inequality to $e^{-\lambda S_n}$ with $\lambda>0$:
\[
\Prob(S_n\leq t)
=
\Prob(e^{-\lambda S_n}\geq e^{-\lambda t}).
\]

Therefore,
\[
\Prob(S_n\leq t)
\leq
e^{\lambda t}\E[e^{-\lambda S_n}].
\]

Using independence,
\[
\E[e^{-\lambda S_n}]
=
\prod_{i=1}^n
\left(
1+p_i(e^{-\lambda}-1)
\right).
\]

Using $1+x\leq e^x$,
\[
\E[e^{-\lambda S_n}]
\leq
\exp\left(
(e^{-\lambda}-1)\mu
\right).
\]

Thus,
\[
\Prob(S_n\leq t)
\leq
\exp\left(
\lambda t+(e^{-\lambda}-1)\mu
\right).
\]

Optimizing gives
\[
e^{-\lambda}=\frac{t}{\mu}.
\]

Substitution yields
\[
\Prob(S_n\leq t)
\leq
\exp\left(
-t\log\left(\frac{t}{\mu}\right)+t-\mu
\right)
=
e^{-\mu}
\left(
\frac{e\mu}{t}
\right)^t.
\]
\end{proof}

\section{Large and Small Deviations}

Chernoff's inequality can be written as
\[
\Prob(S_n\geq t)
\leq
\exp\left(
-\mu h\left(\frac{t}{\mu}\right)
\right),
\]
where
\[
h(u)
=
u\log u-u+1.
\]

Indeed, setting $u=t/\mu$ gives
\[
e^{-\mu}\left(\frac{e\mu}{t}\right)^t
=
\exp\left(
-\mu+t-t\log(t/\mu)
\right)
=
\exp\left(
-\mu h(t/\mu)
\right).
\]

\subsection{Large Deviations}

When $t$ is much larger than $\mu$,
\[
e^{-\mu}
\left(
\frac{e\mu}{t}
\right)^t
\]
behaves like
\[
\exp(-t\log t)
\]
up to lower-order terms.

Thus, for large deviations, the Bernoulli Chernoff bound can be much sharper than a purely Gaussian-type bound.

\subsection{Small Deviations}

Now suppose
\[
t=(1+\delta)\mu
\]
with $\delta>0$ small. Then
\[
e^{-\mu}
\left(
\frac{e\mu}{t}
\right)^t
=
e^{-\mu}
\left(
\frac{e}{1+\delta}
\right)^{(1+\delta)\mu}.
\]

Taking logarithms,
\[
-\mu
+
(1+\delta)\mu
-
(1+\delta)\mu\log(1+\delta)
=
\mu\left[
\delta-(1+\delta)\log(1+\delta)
\right].
\]

Using the Taylor expansion
\[
\log(1+\delta)
=
\delta-\frac{\delta^2}{2}
+
O(\delta^3),
\]
we get
\[
(1+\delta)\log(1+\delta)
=
\delta+\frac{\delta^2}{2}
+
O(\delta^3).
\]

Therefore,
\[
\delta-(1+\delta)\log(1+\delta)
=
-\frac{\delta^2}{2}
+
O(\delta^3).
\]

Hence, for small $\delta$,
\[
\Prob(S_n\geq (1+\delta)\mu)
\lesssim
\exp\left(
-\frac{\delta^2\mu}{2}
\right).
\]

A common simplified version is
\[
\Prob(S_n\geq (1+\delta)\mu)
\leq
\exp\left(
-\frac{\delta^2\mu}{3}
\right),
\qquad
0<\delta<1.
\]

\begin{remark}
For Bernoulli sums, the relevant scale is the mean $\mu$, not the total number of samples $n$. This is why Chernoff bounds improve on Hoeffding's inequality when the Bernoulli probabilities are small.
\end{remark}

\section{Erdos--Renyi Graphs}

\begin{definition}[Erdos--Renyi Graph]
The Erdos--Renyi random graph $G(n,p)$ is the random graph formed by fixing a set of $n$ nodes and connecting each unordered pair of distinct nodes independently with probability $p$.
\end{definition}

For a node $i$, the degree of $i$ is the number of edges connected to $i$:
\[
\deg(i)
=
\#\{j\neq i : (i,j)\text{ is an edge}\}.
\]

Since each possible edge adjacent to node $i$ appears independently with probability $p$,
\[
\deg(i)
\sim
\Bin(n-1,p).
\]

Therefore,
\[
\E[\deg(i)]
=
(n-1)p.
\]

We denote the expected degree by
\[
d=(n-1)p.
\]

\begin{remark}
The quantity $d=(n-1)p$ is the average degree scale of the graph. Many graph properties change sharply depending on how $d$ grows with $n$.
\end{remark}

\subsection{Graph Regularity}

\begin{definition}
A graph is $d$-regular if
\[
\deg(i)=d
\]
for every node $i$.
\end{definition}

A random graph $G(n,p)$ is generally not exactly regular. However, when the expected degree is large enough, every degree is close to its expectation with high probability.

\begin{proposition}[Almost Regularity of $G(n,p)$]
Let $G\sim G(n,p)$ and define
\[
d=(n-1)p.
\]
Fix $\delta>0$. If
\[
d\geq \frac{8}{\delta^2}\log(2n),
\]
then
\[
\Prob\left(
\forall i\in\{1,\dots,n\},
\,
(1-\delta)d
\leq
\deg(i)
\leq
(1+\delta)d
\right)
\geq
1-\frac{1}{2n}.
\]
\end{proposition}

\begin{proof}
For a fixed node $i$,
\[
\deg(i)=\sum_{j\neq i} A_{ij},
\]
where $A_{ij}$ is the indicator that the edge $(i,j)$ is present.

Thus,
\[
A_{ij}\sim \Bern(p)
\]
and
\[
\E[\deg(i)]=d.
\]

By the Chernoff bound for Bernoulli sums,
\[
\Prob\left(
|\deg(i)-d|\geq \delta d
\right)
\leq
2\exp\left(
-\frac{\delta^2 d}{4}
\right),
\]
where the constant is chosen for a standard two-sided small-deviation form.

Let
\[
E_i
=
\left\{
|\deg(i)-d|<\delta d
\right\}.
\]

Then
\[
\Prob(E_i^c)
\leq
2\exp\left(
-\frac{\delta^2 d}{4}
\right).
\]

Using the union bound,
\[
\Prob\left(
\bigcap_{i=1}^n E_i
\right)
=
1-
\Prob\left(
\bigcup_{i=1}^n E_i^c
\right)
\geq
1-\sum_{i=1}^n \Prob(E_i^c).
\]

Therefore,
\[
\Prob\left(
\bigcap_{i=1}^n E_i
\right)
\geq
1-
2n\exp\left(
-\frac{\delta^2 d}{4}
\right).
\]

If
\[
d\geq \frac{8}{\delta^2}\log(2n),
\]
then
\[
\frac{\delta^2 d}{4}
\geq
2\log(2n).
\]

Hence,
\[
2n\exp\left(
-\frac{\delta^2 d}{4}
\right)
\leq
2n\exp(-2\log(2n))
=
\frac{1}{2n}.
\]

Thus,
\[
\Prob\left(
\forall i,
\,
(1-\delta)d
\leq
\deg(i)
\leq
(1+\delta)d
\right)
\geq
1-\frac{1}{2n}.
\]
\end{proof}

\begin{remark}
This result says that when the expected degree is at least logarithmic in $n$, all degrees concentrate uniformly around their common mean.
\end{remark}

\section{Sub-Exponential Random Variables}

Not all important random variables are sub-Gaussian. A key example is the square of a Gaussian random variable.

Let
\[
Z\sim \Normal(0,1).
\]

Then $Z^2$ has heavier tails than a Gaussian random variable. To see this, compute the centered moment generating function:
\[
\E\left[e^{\lambda(Z^2-1)}\right].
\]

We have
\[
\E\left[e^{\lambda(Z^2-1)}\right]
=
\frac{1}{\sqrt{2\pi}}
\int_{-\infty}^{\infty}
e^{\lambda(t^2-1)}
e^{-t^2/2}
\dd t.
\]

Combining exponents,
\[
\E\left[e^{\lambda(Z^2-1)}\right]
=
\frac{e^{-\lambda}}{\sqrt{2\pi}}
\int_{-\infty}^{\infty}
\exp\left(
-\frac{1-2\lambda}{2}t^2
\right)
\dd t.
\]

The integral is finite only if
\[
1-2\lambda>0,
\]
that is,
\[
\lambda<\frac{1}{2}.
\]

For $\lambda<1/2$,
\[
\int_{-\infty}^{\infty}
\exp\left(
-\frac{1-2\lambda}{2}t^2
\right)
\dd t
=
\sqrt{\frac{2\pi}{1-2\lambda}}.
\]

Therefore,
\[
\E\left[e^{\lambda(Z^2-1)}\right]
=
\frac{e^{-\lambda}}{\sqrt{1-2\lambda}},
\qquad
\lambda<\frac{1}{2}.
\]

This moment generating function is not finite for all $\lambda$, so $Z^2$ is not sub-Gaussian.

\subsection{A Useful Bound for Gaussian Squares}

For sufficiently small $\lambda$,
\[
\E\left[e^{\lambda(Z^2-1)}\right]
\]
can still be bounded by a Gaussian-like expression.

Indeed, for $|\lambda|\leq 1/4$,
\[
\E\left[e^{\lambda(Z^2-1)}\right]
=
\frac{e^{-\lambda}}{\sqrt{1-2\lambda}}
\leq
e^{2\lambda^2}.
\]

Equivalently,
\[
\log\E\left[e^{\lambda(Z^2-1)}\right]
\leq
2\lambda^2,
\qquad
|\lambda|\leq \frac{1}{4}.
\]

This motivates the following definition.

\begin{definition}[Sub-Exponential Random Variable]
A random variable $X$ with mean
\[
\mu=\E[X]
\]
is called $(\sigma,\alpha)$-sub-exponential if
\[
\E\left[e^{\lambda(X-\mu)}\right]
\leq
e^{\sigma^2\lambda^2/2}
\]
for all
\[
|\lambda|\leq \frac{1}{\alpha}.
\]
\end{definition}

\begin{example}
If
\[
Z\sim \Normal(0,1),
\]
then $Z^2$ is $(2,4)$-sub-exponential because
\[
\log\E\left[e^{\lambda(Z^2-1)}\right]
\leq
2\lambda^2
=
\frac{2^2\lambda^2}{2}
\]
for all $|\lambda|\leq 1/4$.
\end{example}

\section{Tail Bounds for Sub-Exponential Random Variables}

Sub-exponential random variables have two tail regimes:
\begin{itemize}
\item Gaussian-like decay for moderate deviations,
\item exponential decay for large deviations.
\end{itemize}

\begin{theorem}[Sub-Exponential Tail Bound]
Let $X$ be $(\sigma,\alpha)$-sub-exponential with mean $\mu$. Then, for every $t\geq0$,
\[
\Prob(X-\mu\geq t)
\leq
\begin{cases}
\exp\left(-\dfrac{t^2}{2\sigma^2}\right),
&
0\leq t\leq \dfrac{\sigma^2}{\alpha},
\\[1.25em]
\exp\left(-\dfrac{t}{\alpha}+\dfrac{\sigma^2}{2\alpha^2}\right),
&
t> \dfrac{\sigma^2}{\alpha}.
\end{cases}
\]
\end{theorem}

\begin{proof}
By the Chernoff method, for any $\lambda\in[0,1/\alpha]$,
\[
\Prob(X-\mu\geq t)
\leq
\exp\left(
-\lambda t+\frac{\sigma^2\lambda^2}{2}
\right).
\]

Thus,
\[
\log\Prob(X-\mu\geq t)
\leq
-\sup_{0\leq \lambda\leq 1/\alpha}
\left\{
\lambda t-\frac{\sigma^2\lambda^2}{2}
\right\}.
\]

Consider
\[
q(\lambda)
=
\lambda t-\frac{\sigma^2\lambda^2}{2}.
\]

The unconstrained maximizer satisfies
\[
q'(\lambda)=t-\sigma^2\lambda=0,
\]
so
\[
\lambda^*=\frac{t}{\sigma^2}.
\]

There are two cases.

If
\[
\frac{t}{\sigma^2}\leq \frac{1}{\alpha},
\]
equivalently
\[
t\leq \frac{\sigma^2}{\alpha},
\]
then the unconstrained maximizer is feasible. Hence,
\[
\sup_{0\leq \lambda\leq 1/\alpha}
q(\lambda)
=
q\left(\frac{t}{\sigma^2}\right)
=
\frac{t^2}{2\sigma^2}.
\]

Therefore,
\[
\Prob(X-\mu\geq t)
\leq
\exp\left(
-\frac{t^2}{2\sigma^2}
\right).
\]

If
\[
t> \frac{\sigma^2}{\alpha},
\]
then the unconstrained maximizer is not feasible, and the supremum over $[0,1/\alpha]$ occurs at the endpoint
\[
\lambda=\frac{1}{\alpha}.
\]

Thus,
\[
\sup_{0\leq \lambda\leq 1/\alpha}
q(\lambda)
=
\frac{t}{\alpha}
-
\frac{\sigma^2}{2\alpha^2}.
\]

Therefore,
\[
\Prob(X-\mu\geq t)
\leq
\exp\left(
-\frac{t}{\alpha}
+
\frac{\sigma^2}{2\alpha^2}
\right).
\]
\end{proof}

\begin{corollary}[Two-Sided Sub-Exponential Tail Bound]
Let $X$ be $(\sigma,\alpha)$-sub-exponential with mean $\mu$. Then
\[
\Prob(|X-\mu|\geq t)
\leq
2\exp\left(
-\min\left\{
\frac{t^2}{2\sigma^2},
\frac{t}{2\alpha}
\right\}
\right).
\]
\end{corollary}

\begin{proof}
The upper-tail theorem gives Gaussian-type decay for
\[
t\leq \frac{\sigma^2}{\alpha}
\]
and exponential-type decay for
\[
t> \frac{\sigma^2}{\alpha}.
\]

The same argument applies to $-X$, since the defining moment generating function bound is symmetric in $|\lambda|$.

Thus, applying the union bound,
\[
\Prob(|X-\mu|\geq t)
\leq
\Prob(X-\mu\geq t)
+
\Prob(\mu-X\geq t).
\]

Both tails are bounded by the same expression. Combining the two regimes into a single minimum gives
\[
\Prob(|X-\mu|\geq t)
\leq
2\exp\left(
-\min\left\{
\frac{t^2}{2\sigma^2},
\frac{t}{2\alpha}
\right\}
\right).
\]
\end{proof}

\section{Sums of Sub-Exponential Random Variables}

\begin{lemma}[Sum Rule]
Let
\[
X_1,\dots,X_n
\]
be independent random variables, where $X_i$ is $(\sigma_i,\alpha_i)$-sub-exponential.

Then
\[
\sum_{i=1}^n X_i
\]
is
\[
\left(
\norm{\sigma}_2,
\norm{\alpha}_\infty
\right)
\]
-sub-exponential, where
\[
\norm{\sigma}_2
=
\left(
\sum_{i=1}^n \sigma_i^2
\right)^{1/2}
\]
and
\[
\norm{\alpha}_\infty
=
\max_{1\leq i\leq n}\alpha_i.
\]
\end{lemma}

\begin{proof}
Let
\[
\mu_i=\E[X_i].
\]

For
\[
|\lambda|\leq \frac{1}{\norm{\alpha}_\infty},
\]
we have
\[
|\lambda|\leq \frac{1}{\alpha_i}
\]
for every $i$.

By independence,
\[
\E\left[
e^{\lambda\sum_{i=1}^n (X_i-\mu_i)}
\right]
=
\prod_{i=1}^n
\E\left[
e^{\lambda(X_i-\mu_i)}
\right].
\]

Using sub-exponentiality,
\[
\E\left[
e^{\lambda(X_i-\mu_i)}
\right]
\leq
e^{\sigma_i^2\lambda^2/2}.
\]

Therefore,
\[
\E\left[
e^{\lambda\sum_{i=1}^n (X_i-\mu_i)}
\right]
\leq
\prod_{i=1}^n
e^{\sigma_i^2\lambda^2/2}
=
e^{\lambda^2\sum_{i=1}^n\sigma_i^2/2}.
\]

Since
\[
\sum_{i=1}^n\sigma_i^2
=
\norm{\sigma}_2^2,
\]
the sum is
\[
\left(
\norm{\sigma}_2,
\norm{\alpha}_\infty
\right)
\]
-sub-exponential.
\end{proof}

\section{Bernstein's Inequality}

Combining the sub-exponential sum rule with the two-sided tail bound gives Bernstein's inequality.

\begin{theorem}[Bernstein's Inequality]
Let
\[
X_1,\dots,X_n
\]
be independent random variables such that $X_i$ is $(\sigma_i,\alpha_i)$-sub-exponential. Let
\[
\mu_i=\E[X_i].
\]

Then, for every $t\geq0$,
\[
\Prob\left(
\left|
\sum_{i=1}^n (X_i-\mu_i)
\right|
\geq t
\right)
\leq
2\exp\left(
-\min\left\{
\frac{t^2}{2\norm{\sigma}_2^2},
\frac{t}{2\norm{\alpha}_\infty}
\right\}
\right).
\]
\end{theorem}

\begin{proof}
By the sum rule,
\[
\sum_{i=1}^n X_i
\]
is
\[
\left(
\norm{\sigma}_2,
\norm{\alpha}_\infty
\right)
\]
-sub-exponential.

Applying the two-sided sub-exponential tail bound to
\[
\sum_{i=1}^n (X_i-\mu_i)
\]
gives
\[
\Prob\left(
\left|
\sum_{i=1}^n (X_i-\mu_i)
\right|
\geq t
\right)
\leq
2\exp\left(
-\min\left\{
\frac{t^2}{2\norm{\sigma}_2^2},
\frac{t}{2\norm{\alpha}_\infty}
\right\}
\right).
\]
\end{proof}

\begin{remark}
Bernstein's inequality interpolates between two regimes. For moderate deviations, the bound behaves like
\[
\exp\left(
-\frac{t^2}{2\norm{\sigma}_2^2}
\right).
\]

For large deviations, the bound behaves like
\[
\exp\left(
-\frac{t}{2\norm{\alpha}_\infty}
\right).
\]
\end{remark}

\chapter{Dimension Reduction and Orlicz Norms}

\section{Dimension Reduction}

Suppose we have a finite set of points
\[
\mathcal{X}=\{x_1,\dots,x_m\}\subseteq \R^d.
\]

The goal of dimension reduction is to construct a map
\[
f:\R^d\to\R^n,
\qquad n<d,
\]
such that the geometry of $\mathcal{X}$ is approximately preserved.

In particular, we want
\[
(1-\varepsilon)\norm{x-y}_2^2
\leq
\norm{f(x)-f(y)}_2^2
\leq
(1+\varepsilon)\norm{x-y}_2^2
\]
for all
\[
x,y\in\mathcal{X}.
\]

The central question is: how small can the embedding dimension $n$ be?

\section{The Johnson--Lindenstrauss Lemma}

\begin{theorem}[Johnson--Lindenstrauss Lemma]
Fix
\[
\varepsilon,\delta\in(0,1),
\]
and let
\[
\mathcal{X}\subseteq \R^d
\]
be a set with $m$ points.

Suppose
\[
n
\geq
\frac{8}{\varepsilon^2}
\log\left(\frac{m^2}{\delta}\right).
\]

Let
\[
M\in\R^{n\times d}
\]
be a random matrix with iid $\Normal(0,1)$ entries. Define
\[
f(x)=\frac{1}{\sqrt{n}}Mx.
\]

Then, with probability at least $1-\delta$,
\[
(1-\varepsilon)\norm{x-y}_2^2
\leq
\norm{f(x)-f(y)}_2^2
\leq
(1+\varepsilon)\norm{x-y}_2^2
\]
for all $x,y\in\mathcal{X}$.
\end{theorem}

\begin{remark}
The embedding dimension depends logarithmically on the number of points $m$, but does not depend on the original dimension $d$.
\end{remark}

\begin{remark}
The map is oblivious to the data: the random matrix $M$ is sampled without using the structure of $\mathcal{X}$.
\end{remark}

\begin{remark}
The theorem is an example of the probabilistic method. It proves that a good low-dimensional embedding exists by showing that a random embedding succeeds with positive probability.
\end{remark}

\begin{proof}
Fix two distinct points
\[
x,y\in\mathcal{X},
\]
and define
\[
z=x-y.
\]

Since $f$ is linear,
\[
f(x)-f(y)
=
\frac{1}{\sqrt{n}}M(x-y)
=
\frac{1}{\sqrt{n}}Mz.
\]

Therefore,
\[
\norm{f(x)-f(y)}_2^2
=
\frac{1}{n}\norm{Mz}_2^2.
\]

Let $M_i$ denote the $i$-th row of $M$. Then
\[
\norm{Mz}_2^2
=
\sum_{i=1}^n \inner{M_i}{z}^2.
\]

Since $M_i\sim \Normal(0,I_d)$, rotational invariance of the Gaussian distribution gives
\[
\inner{M_i}{z}
\sim
\Normal(0,\norm{z}_2^2).
\]

Equivalently,
\[
\frac{\inner{M_i}{z}}{\norm{z}_2}
\sim
\Normal(0,1).
\]

Thus,
\[
\frac{\norm{Mz}_2^2}{\norm{z}_2^2}
=
\sum_{i=1}^n
\left(
\frac{\inner{M_i}{z}}{\norm{z}_2}
\right)^2
\]
is a sum of $n$ independent squared standard Gaussians.

Let
\[
Z_i
=
\left(
\frac{\inner{M_i}{z}}{\norm{z}_2}
\right)^2.
\]

Then
\[
Z_i\sim \chi_1^2,
\qquad
\E[Z_i]=1.
\]

From the sub-exponential calculation for squared Gaussians,
\[
Z_i-1
\]
is sub-exponential. Hence
\[
\frac{1}{n}
\sum_{i=1}^n (Z_i-1)
=
\frac{\norm{Mz}_2^2}{n\norm{z}_2^2}
-
1
\]
satisfies a Bernstein-type bound. In particular, for $0<\varepsilon<1$,
\[
\Prob\left(
\left|
\frac{\norm{Mz}_2^2}{n\norm{z}_2^2}
-
1
\right|
\geq
\varepsilon
\right)
\leq
2\exp\left(
-\frac{n\varepsilon^2}{8}
\right).
\]

Equivalently,
\[
\Prob\left(
\left|
\frac{\norm{f(x)-f(y)}_2^2}{\norm{x-y}_2^2}
-
1
\right|
\geq
\varepsilon
\right)
\leq
2\exp\left(
-\frac{n\varepsilon^2}{8}
\right).
\]

Thus, for this fixed pair $x,y$,
\[
(1-\varepsilon)\norm{x-y}_2^2
\leq
\norm{f(x)-f(y)}_2^2
\leq
(1+\varepsilon)\norm{x-y}_2^2
\]
fails with probability at most
\[
2\exp\left(
-\frac{n\varepsilon^2}{8}
\right).
\]

There are at most
\[
\binom{m}{2}
\leq
\frac{m^2}{2}
\]
distinct pairs of points. By the union bound, the probability that the distortion bound fails for at least one pair is at most
\[
2\binom{m}{2}
\exp\left(
-\frac{n\varepsilon^2}{8}
\right)
\leq
m^2
\exp\left(
-\frac{n\varepsilon^2}{8}
\right).
\]

If
\[
n
\geq
\frac{8}{\varepsilon^2}
\log\left(\frac{m^2}{\delta}\right),
\]
then
\[
m^2
\exp\left(
-\frac{n\varepsilon^2}{8}
\right)
\leq
\delta.
\]

Therefore, with probability at least $1-\delta$, the desired inequality holds simultaneously for all pairs $x,y\in\mathcal{X}$.
\end{proof}

\section{Equivalent Characterizations of Sub-Gaussian Random Variables}

There are several equivalent ways to define sub-Gaussian random variables. These equivalences are useful because different definitions are convenient in different arguments.

\begin{proposition}[Equivalent Sub-Gaussian Conditions]
Let $X$ be a random variable. The following conditions are equivalent up to universal constant factors.

\begin{enumerate}
\item There exists $K_1>0$ such that
\[
\Prob(|X|>t)
\leq
2\exp\left(
-\frac{t^2}{K_1^2}
\right)
\]
for all $t\geq0$.

\item There exists $K_2>0$ such that
\[
\norm{X}_{L^p}
=
\left(\E|X|^p\right)^{1/p}
\leq
K_2\sqrt{p}
\]
for all $p\geq1$.

\item There exists $K_3>0$ such that
\[
\E\left[
\exp\left(
\frac{X^2}{K_3^2}
\right)
\right]
\leq
2.
\]
\end{enumerate}

Moreover, if $\E[X]=0$, then these are also equivalent to the moment generating function condition: there exists $K_4>0$ such that
\[
\E[e^{\lambda X}]
\leq
\exp(K_4^2\lambda^2)
\]
for all $\lambda\in\R$.
\end{proposition}

\begin{proof}
We prove the chain of implications.

First suppose the tail bound holds. By rescaling, assume $K_1=1$, so
\[
\Prob(|X|>t)
\leq
2e^{-t^2}.
\]

Using the tail integral formula,
\[
\E|X|^p
=
\int_0^\infty
p t^{p-1}\Prob(|X|>t)\dd t.
\]

Hence,
\[
\E|X|^p
\leq
2p
\int_0^\infty
t^{p-1}e^{-t^2}\dd t.
\]

Make the change of variables
\[
s=t^2,
\qquad
\dd s=2t\dd t.
\]

Then
\[
\int_0^\infty
t^{p-1}e^{-t^2}\dd t
=
\frac{1}{2}
\int_0^\infty
s^{p/2-1}e^{-s}\dd s
=
\frac{1}{2}\Gamma(p/2).
\]

Therefore,
\[
\E|X|^p
\leq
p\Gamma(p/2).
\]

Using the standard Gamma-function bound
\[
\Gamma(p/2)
\leq
C^p p^{p/2}
\]
for a universal constant $C>0$, we get
\[
\E|X|^p
\leq
(C\sqrt{p})^p.
\]

Thus,
\[
\norm{X}_{L^p}
\leq
C\sqrt{p}.
\]

This proves that the tail condition implies the moment-growth condition.

Next assume
\[
\norm{X}_{L^p}
\leq
K_2\sqrt{p}
\]
for all $p\geq1$. By rescaling, assume $K_2=1$. Then
\[
\E|X|^{2p}
\leq
(2p)^p.
\]

Using the power series expansion of the exponential,
\[
\E e^{\lambda^2X^2}
=
1+
\sum_{p=1}^\infty
\frac{\lambda^{2p}\E X^{2p}}{p!}.
\]

Thus,
\[
\E e^{\lambda^2X^2}
\leq
1+
\sum_{p=1}^\infty
\frac{\lambda^{2p}(2p)^p}{p!}.
\]

Using the lower bound
\[
p!\geq \left(\frac{p}{e}\right)^p,
\]
we obtain
\[
\frac{(2p)^p}{p!}
\leq
(2e)^p.
\]

Therefore,
\[
\E e^{\lambda^2X^2}
\leq
1+
\sum_{p=1}^\infty
(2e\lambda^2)^p.
\]

Choose $\lambda>0$ sufficiently small so that
\[
2e\lambda^2\leq \frac{1}{2}.
\]

Then the geometric series is bounded, and
\[
\E e^{\lambda^2X^2}
\leq
2.
\]

Equivalently, for some universal constant $K_3>0$,
\[
\E\left[
\exp\left(
\frac{X^2}{K_3^2}
\right)
\right]
\leq
2.
\]

This proves the exponential-square integrability condition.

Now assume
\[
\E\left[
\exp\left(
\frac{X^2}{K_3^2}
\right)
\right]
\leq
2.
\]

By rescaling, assume $K_3=1$. Then by Markov's inequality,
\[
\Prob(|X|>t)
=
\Prob(e^{X^2}>e^{t^2})
\leq
e^{-t^2}\E[e^{X^2}]
\leq
2e^{-t^2}.
\]

This recovers the sub-Gaussian tail bound.

Finally, if $\E[X]=0$, the moment generating function condition implies the tail condition by the Chernoff method. Conversely, the preceding moment-growth bounds imply convergence and Gaussian-type control of the moment generating function, yielding
\[
\E[e^{\lambda X}]
\leq
\exp(CK^2\lambda^2)
\]
for a universal constant $C>0$.
\end{proof}

\section{The Sub-Gaussian Norm}

The preceding characterization motivates a norm-like quantity.

\begin{definition}[Sub-Gaussian Norm]
The sub-Gaussian norm of a random variable $X$ is
\[
\norm{X}_{\psi_2}
=
\inf\left\{
K>0:
\E\left[
\exp\left(
\frac{X^2}{K^2}
\right)
\right]
\leq
2
\right\}.
\]
\end{definition}

\begin{remark}
The function
\[
\norm{\cdot}_{\psi_2}
\]
is a norm on the space of random variables with finite sub-Gaussian norm.
\end{remark}

\begin{lemma}[Centering a Sub-Gaussian Random Variable]
There exists a universal constant $C>0$ such that
\[
\norm{X-\E[X]}_{\psi_2}
\leq
C\norm{X}_{\psi_2}.
\]
\end{lemma}

\begin{proof}
By the triangle inequality for $\norm{\cdot}_{\psi_2}$,
\[
\norm{X-\E[X]}_{\psi_2}
\leq
\norm{X}_{\psi_2}
+
\norm{\E[X]}_{\psi_2}.
\]

Since $\E[X]$ is a constant random variable,
\[
\norm{\E[X]}_{\psi_2}
=
|\E[X]|\norm{1}_{\psi_2}.
\]

By Jensen's inequality,
\[
|\E[X]|
\leq
\E|X|.
\]

The $L^1$ norm is controlled by the sub-Gaussian norm, so
\[
\E|X|
\leq
C'\norm{X}_{\psi_2}.
\]

Thus,
\[
\norm{\E[X]}_{\psi_2}
\leq
C''\norm{X}_{\psi_2}.
\]

Combining the bounds,
\[
\norm{X-\E[X]}_{\psi_2}
\leq
C\norm{X}_{\psi_2}.
\]
\end{proof}

\begin{theorem}[Hoeffding Inequality via the Sub-Gaussian Norm]
Let
\[
X_1,\dots,X_n
\]
be independent mean-zero sub-Gaussian random variables. Then there exists a universal constant $C>0$ such that
\[
\norm{\sum_{i=1}^n X_i}_{\psi_2}^2
\leq
C\sum_{i=1}^n \norm{X_i}_{\psi_2}^2.
\]
\end{theorem}

\begin{remark}
This is a norm-based version of the sum rule for independent sub-Gaussian random variables.
\end{remark}

\section{Sub-Exponential Norms}

Sub-exponential random variables also admit equivalent characterizations.

\begin{proposition}[Equivalent Sub-Exponential Conditions]
Let $X$ be a random variable. The following are equivalent up to universal constant factors.

\begin{enumerate}
\item There exists $K_1>0$ such that
\[
\Prob(|X|>t)
\leq
2\exp\left(
-\frac{t}{K_1}
\right)
\]
for all $t\geq0$.

\item There exists $K_2>0$ such that
\[
\norm{X}_{L^p}
\leq
K_2p
\]
for all $p\geq1$.

\item There exists $K_3>0$ such that
\[
\E\left[
\exp\left(
\frac{|X|}{K_3}
\right)
\right]
\leq
2.
\]
\end{enumerate}

Moreover, if $\E[X]=0$, then these are also equivalent to the local moment generating function condition: there exists $K_4>0$ such that
\[
\E[e^{\lambda X}]
\leq
\exp(K_4^2\lambda^2)
\]
for all
\[
|\lambda|\leq \frac{1}{K_4}.
\]
\end{proposition}

\begin{definition}[Sub-Exponential Norm]
The sub-exponential norm of a random variable $X$ is
\[
\norm{X}_{\psi_1}
=
\inf\left\{
K>0:
\E\left[
\exp\left(
\frac{|X|}{K}
\right)
\right]
\leq
2
\right\}.
\]
\end{definition}

\begin{remark}
The function
\[
\norm{\cdot}_{\psi_1}
\]
is a norm on the space of sub-exponential random variables.
\end{remark}

\begin{lemma}[Centering a Sub-Exponential Random Variable]
There exists a universal constant $C>0$ such that
\[
\norm{X-\E[X]}_{\psi_1}
\leq
C\norm{X}_{\psi_1}.
\]
\end{lemma}

\section{Products of Sub-Gaussian Random Variables}

The square of a Gaussian motivates sub-exponential random variables. More generally, products of sub-Gaussian random variables are sub-exponential.

\begin{lemma}
Suppose $X$ and $Y$ are sub-Gaussian random variables. Then there exists a universal constant $C>0$ such that
\[
\norm{XY}_{\psi_1}
\leq
C\norm{X}_{\psi_2}\norm{Y}_{\psi_2}.
\]
\end{lemma}

\begin{proof}
By homogeneity, it suffices to prove the result when
\[
\norm{X}_{\psi_2}=1,
\qquad
\norm{Y}_{\psi_2}=1.
\]

Then
\[
\E[e^{X^2}]\leq 2,
\qquad
\E[e^{Y^2}]\leq 2.
\]

Using Young's inequality,
\[
|XY|
\leq
\frac{X^2}{2}
+
\frac{Y^2}{2}.
\]

Therefore,
\[
e^{|XY|}
\leq
\exp\left(
\frac{X^2}{2}
+
\frac{Y^2}{2}
\right).
\]

By the arithmetic-geometric mean inequality,
\[
\exp\left(
\frac{X^2}{2}
+
\frac{Y^2}{2}
\right)
=
e^{X^2/2}e^{Y^2/2}
\leq
\frac{1}{2}e^{X^2}
+
\frac{1}{2}e^{Y^2}.
\]

Taking expectations,
\[
\E[e^{|XY|}]
\leq
\frac{1}{2}\E[e^{X^2}]
+
\frac{1}{2}\E[e^{Y^2}]
\leq
2.
\]

Thus,
\[
\norm{XY}_{\psi_1}
\leq
1
\]
under the normalized assumption. Rescaling gives the general result up to a universal constant.
\end{proof}

\section{Orlicz Norms}

The sub-Gaussian and sub-exponential norms are special cases of a more general construction.

\begin{definition}[Orlicz Norm]
Let
\[
\psi:\R_+\to\R_+
\]
be convex and nondecreasing, with
\[
\psi(0)=0
\]
and
\[
\lim_{t\to\infty}\psi(t)=\infty.
\]

The Orlicz norm of a random variable $X$ associated with $\psi$ is
\[
\norm{X}_{\psi}
=
\inf\left\{
K>0:
\E\left[
\psi\left(
\frac{|X|}{K}
\right)
\right]
\leq
1
\right\}.
\]
\end{definition}

\begin{remark}
Under standard assumptions on $\psi$, this defines a norm on the space
\[
\{X:\norm{X}_{\psi}<\infty\}.
\]
\end{remark}

\begin{example}
If
\[
\psi(t)=t^p
\]
for $p\geq1$, then the corresponding Orlicz norm is the usual $L^p$ norm:
\[
\norm{X}_{\psi}
=
\norm{X}_{L^p}.
\]
\end{example}

\begin{example}
If
\[
\psi_2(t)=e^{t^2}-1,
\]
then the corresponding Orlicz norm is equivalent to the sub-Gaussian norm.
\end{example}

\begin{example}
If
\[
\psi_1(t)=e^t-1,
\]
then the corresponding Orlicz norm is equivalent to the sub-exponential norm.
\end{example}

\section{Concentration of Functions of Independent Random Variables}

So far, we have mainly studied concentration of sums. However, concentration is a more general phenomenon.

The guiding principle is:

\begin{quote}
If $X_1,\dots,X_n$ are independent random variables, then $f(X_1,\dots,X_n)$ concentrates near its mean provided that $f$ is not too sensitive to any single coordinate.
\end{quote}

This section studies one formal version of this principle: bounded differences.

\section{Bounded Differences}

\begin{definition}[Bounded Differences Property]
Let
\[
f:\R^n\to\R.
\]
We say that $f$ satisfies the bounded differences property with constants
\[
c_1,\dots,c_n>0
\]
if, for every $j\in\{1,\dots,n\}$,
\[
\left|
f(z_1,\dots,z_j,\dots,z_n)
-
f(z_1,\dots,\widetilde{z}_j,\dots,z_n)
\right|
\leq
c_j
\]
for all
\[
z_1,\dots,z_n,\widetilde{z}_j\in\R.
\]
\end{definition}

This means that changing only the $j$-th coordinate can change the value of $f$ by at most $c_j$.

\section{McDiarmid's Inequality}

\begin{theorem}[McDiarmid's Inequality]
Let
\[
X_1,\dots,X_n
\]
be independent random variables, and let
\[
f:\R^n\to\R
\]
satisfy the bounded differences property with constants
\[
c_1,\dots,c_n.
\]

Then, for every $t\geq0$,
\[
\Prob\left(
\left|
f(X_1,\dots,X_n)
-
\E[f(X_1,\dots,X_n)]
\right|
\geq t
\right)
\leq
2\exp\left(
-\frac{2t^2}{\sum_{j=1}^n c_j^2}
\right).
\]
\end{theorem}

\begin{remark}
McDiarmid's inequality extends Hoeffding-type concentration from sums to general functions of independent random variables.
\end{remark}

\section{Martingales}

To prove McDiarmid's inequality, we use the martingale method.

\begin{definition}[Martingale]
Let
\[
X_0,X_1,X_2,\dots
\]
be a sequence of random variables. A sequence
\[
Y_0,Y_1,Y_2,\dots
\]
is called a martingale with respect to
\[
X_0,X_1,X_2,\dots
\]
if, for every $n$,
\[
\E|Y_n|<\infty,
\]
\[
Y_n
\]
is measurable with respect to
\[
X_0,\dots,X_n,
\]
and
\[
\E[Y_{n+1}\mid X_0,\dots,X_n]
=
Y_n.
\]
\end{definition}

\begin{remark}
Martingales model fair games. They are useful when full independence is unavailable, and they appear in central limit theorems, concentration inequalities, and stochastic process theory.
\end{remark}

\subsection{Conditional Expectation Facts}

We will use two basic facts about conditional expectation.

\begin{lemma}[Tower Property]
If
\[
j<k,
\]
then
\[
\E\left[
\E[Y\mid X_1,\dots,X_k]
\mid X_1,\dots,X_j
\right]
=
\E[Y\mid X_1,\dots,X_j].
\]
\end{lemma}

\begin{lemma}[Measurability Property]
If $Y$ is measurable with respect to
\[
X_1,\dots,X_k,
\]
then
\[
\E[Y\mid X_1,\dots,X_k]
=
Y.
\]
\end{lemma}

\section{The Doob Martingale}

Let
\[
X=(X_1,\dots,X_n).
\]

Define
\[
Y_0
=
\E[f(X)]
\]
and, for $j=1,\dots,n$,
\[
Y_j
=
\E[f(X)\mid X_1,\dots,X_j].
\]

Then
\[
Y_n
=
f(X_1,\dots,X_n)
\]
because $f(X)$ is measurable with respect to $X_1,\dots,X_n$.

Hence,
\[
f(X)-\E[f(X)]
=
Y_n-Y_0.
\]

This difference can be decomposed as
\[
Y_n-Y_0
=
\sum_{j=0}^{n-1}(Y_{j+1}-Y_j).
\]

\begin{proposition}
The sequence
\[
Y_0,Y_1,\dots,Y_n
\]
defined by
\[
Y_j=\E[f(X)\mid X_1,\dots,X_j]
\]
is a martingale with respect to
\[
X_1,\dots,X_n.
\]
\end{proposition}

\begin{proof}
First, $Y_j$ is measurable with respect to $X_1,\dots,X_j$ by definition.

Next, by the tower property,
\[
\E[Y_{j+1}\mid X_1,\dots,X_j]
=
\E\left[
\E[f(X)\mid X_1,\dots,X_{j+1}]
\mid X_1,\dots,X_j
\right]
=
\E[f(X)\mid X_1,\dots,X_j].
\]

Therefore,
\[
\E[Y_{j+1}\mid X_1,\dots,X_j]
=
Y_j.
\]

Thus, $\{Y_j\}_{j=0}^n$ is a martingale.
\end{proof}

\section{Martingale Differences}

Define the martingale differences
\[
\Delta_j
=
Y_j-Y_{j-1},
\qquad
j=1,\dots,n.
\]

Then
\[
f(X)-\E[f(X)]
=
\sum_{j=1}^n \Delta_j.
\]

Therefore, controlling the concentration of $f(X)-\E[f(X)]$ reduces to controlling a sum of martingale differences.

\section{Azuma's Inequality}

\begin{lemma}[Azuma's Inequality: Sub-Gaussian Difference Form]
Suppose
\[
Y_0,Y_1,\dots,Y_n
\]
is a martingale with respect to
\[
X_1,\dots,X_n.
\]

Let
\[
\Delta_j=Y_j-Y_{j-1}.
\]

Assume that, for each $j$,
\[
\E\left[
e^{\lambda \Delta_j}
\mid
X_1,\dots,X_{j-1}
\right]
\leq
e^{\lambda^2\sigma_j^2/2}
\]
almost surely for all $\lambda\in\R$.

Then
\[
\sum_{j=1}^n \Delta_j
\]
is
\[
\norm{\sigma}_2
\]
-sub-Gaussian, where
\[
\norm{\sigma}_2^2
=
\sum_{j=1}^n \sigma_j^2.
\]

Consequently,
\[
\Prob\left(
\left|
\sum_{j=1}^n \Delta_j
\right|
\geq t
\right)
\leq
2\exp\left(
-\frac{t^2}{2\sum_{j=1}^n \sigma_j^2}
\right).
\]
\end{lemma}

\begin{proof}
We prove the moment generating function bound. For $\lambda\in\R$,
\[
\E\left[
e^{\lambda\sum_{j=1}^n \Delta_j}
\right]
=
\E\left[
e^{\lambda\sum_{j=1}^{n-1} \Delta_j}
e^{\lambda\Delta_n}
\right].
\]

Using conditional expectation,
\[
\E\left[
e^{\lambda\sum_{j=1}^n \Delta_j}
\right]
=
\E\left[
e^{\lambda\sum_{j=1}^{n-1} \Delta_j}
\E\left[
e^{\lambda\Delta_n}
\mid
X_1,\dots,X_{n-1}
\right]
\right].
\]

By assumption,
\[
\E\left[
e^{\lambda\Delta_n}
\mid
X_1,\dots,X_{n-1}
\right]
\leq
e^{\lambda^2\sigma_n^2/2}.
\]

Therefore,
\[
\E\left[
e^{\lambda\sum_{j=1}^n \Delta_j}
\right]
\leq
e^{\lambda^2\sigma_n^2/2}
\E\left[
e^{\lambda\sum_{j=1}^{n-1}\Delta_j}
\right].
\]

Iterating this argument gives
\[
\E\left[
e^{\lambda\sum_{j=1}^n \Delta_j}
\right]
\leq
\exp\left(
\frac{\lambda^2}{2}\sum_{j=1}^n\sigma_j^2
\right).
\]

Thus,
\[
\sum_{j=1}^n \Delta_j
\]
is $\norm{\sigma}_2$-sub-Gaussian.

The two-sided tail bound follows from the usual sub-Gaussian tail inequality.
\end{proof}

\begin{remark}
Azuma's inequality is the martingale analogue of the sum rule for independent sub-Gaussian random variables.
\end{remark}

\section{Proof of Azuma's Inequality}

Recall the setup.

Let
\[
Y_0,Y_1,\dots,Y_n
\]
be a martingale with respect to
\[
X_1,\dots,X_n,
\]
and define the martingale differences
\[
\Delta_k
=
Y_k-Y_{k-1}.
\]

Assume that for every $k$,
\[
\E\left[
e^{\lambda \Delta_k}
\mid
X_1,\dots,X_{k-1}
\right]
\leq
e^{\lambda^2\sigma_k^2/2}
\]
almost surely.

We prove that
\[
\sum_{k=1}^n \Delta_k
\]
is
\[
\norm{\sigma}_2
\]
-sub-Gaussian.

\begin{proof}
For any $\lambda\in\R$,
\[
\E\left[
e^{\lambda\sum_{k=1}^n \Delta_k}
\right]
=
\E\left[
\E\left[
e^{\lambda\sum_{k=1}^n \Delta_k}
\mid
X_1,\dots,X_{n-1}
\right]
\right].
\]

Using conditional expectation,
\[
\E\left[
e^{\lambda\sum_{k=1}^n \Delta_k}
\mid
X_1,\dots,X_{n-1}
\right]
=
e^{\lambda\sum_{k=1}^{n-1}\Delta_k}
\E\left[
e^{\lambda\Delta_n}
\mid
X_1,\dots,X_{n-1}
\right].
\]

Applying the assumption,
\[
\E\left[
e^{\lambda\sum_{k=1}^n \Delta_k}
\right]
\leq
e^{\lambda^2\sigma_n^2/2}
\E\left[
e^{\lambda\sum_{k=1}^{n-1}\Delta_k}
\right].
\]

Repeating this argument recursively gives
\[
\E\left[
e^{\lambda\sum_{k=1}^n \Delta_k}
\right]
\leq
\exp\left(
\frac{\lambda^2}{2}
\sum_{k=1}^n \sigma_k^2
\right).
\]

Therefore,
\[
\sum_{k=1}^n \Delta_k
\]
is
\[
\norm{\sigma}_2
\]
-sub-Gaussian.
\end{proof}

\begin{remark}
Although the martingale property
\[
\E[\Delta_k\mid X_1,\dots,X_{k-1}]=0
\]
did not explicitly appear in the iteration argument, it is essential for the conditional moment generating function bound.
\end{remark}

\section{Proof of McDiarmid's Inequality}

Recall the Doob martingale
\[
Y_k
=
\E[f(X_1,\dots,X_n)\mid X_1,\dots,X_k].
\]

We showed that
\[
f(X)-\E[f(X)]
=
\sum_{k=1}^n \Delta_k,
\qquad
\Delta_k=Y_k-Y_{k-1}.
\]

To apply Azuma's inequality, we need to control the conditional moment generating functions of the increments.

\begin{proof}[Proof of McDiarmid's Inequality]
Fix
\[
k\in\{1,\dots,n\}.
\]

Recall
\[
\Delta_k
=
\E[f(X)\mid X_1,\dots,X_k]
-
\E[f(X)\mid X_1,\dots,X_{k-1}].
\]

Condition on
\[
X_1,\dots,X_{k-1}.
\]

Define
\[
A_k
=
\E\left[
\inf_t
f(X_1,\dots,X_{k-1},t,X_{k+1},\dots,X_n)
\mid
X_1,\dots,X_{k-1}
\right]
\]
and
\[
B_k
=
\E\left[
\sup_t
f(X_1,\dots,X_{k-1},t,X_{k+1},\dots,X_n)
\mid
X_1,\dots,X_{k-1}
\right].
\]

Then
\[
A_k
\leq
\Delta_k
\leq
B_k
\]
almost surely.

Indeed, replacing the $k$-th coordinate by an arbitrary value $t$ gives lower and upper bounds on the conditional expectation.

By the bounded differences assumption,
\[
B_k-A_k
\leq
c_k.
\]

Thus, conditioned on
\[
X_1,\dots,X_{k-1},
\]
the random variable $\Delta_k$ lies in an interval of length at most $c_k$.

Applying Hoeffding's lemma conditionally,
\[
\E\left[
e^{\lambda\Delta_k}
\mid
X_1,\dots,X_{k-1}
\right]
\leq
e^{\lambda^2 c_k^2/8}.
\]

Therefore, Azuma's inequality applies with
\[
\sigma_k^2=\frac{c_k^2}{4}.
\]

Hence,
\[
f(X)-\E[f(X)]
=
\sum_{k=1}^n \Delta_k
\]
is
\[
\frac{\norm{c}_2}{2}
\]
-sub-Gaussian.

Consequently,
\[
\Prob\left(
|f(X)-\E[f(X)]|
\geq t
\right)
\leq
2\exp\left(
-\frac{t^2}{2(\norm{c}_2^2/4)}
\right).
\]

Equivalently,
\[
\Prob\left(
|f(X)-\E[f(X)]|
\geq t
\right)
\leq
2\exp\left(
-\frac{2t^2}{\norm{c}_2^2}
\right).
\]
\end{proof}

\section{U-Statistics}

McDiarmid's inequality is especially useful when independence is not directly available but dependence is weak.

\begin{definition}[U-Statistic]
Let
\[
X_1,\dots,X_n
\]
be iid random variables, and let
\[
g:\R^2\to\R
\]
be a symmetric measurable function.

The associated U-statistic is
\[
U(X_1,\dots,X_n)
=
\frac{1}{\binom{n}{2}}
\sum_{i<j} g(X_i,X_j).
\]
\end{definition}

\begin{remark}
The terms
\[
g(X_i,X_j)
\]
are not independent because different pairs share coordinates. Nevertheless, the dependence is weak enough for concentration inequalities to apply.
\end{remark}

\begin{proposition}[Concentration of U-Statistics]
Suppose
\[
|g(x,y)|\leq C
\]
for all
\[
x,y\in\R.
\]

Then, for every $t\geq0$,
\[
\Prob\left(
|U-\E[U]|
\geq t
\right)
\leq
2\exp\left(
-\frac{nt^2}{8C^2}
\right).
\]
\end{proposition}

\begin{proof}
Fix an index $k$ and replace $X_k$ by an alternative value $\widetilde{X}_k$.

Only terms involving the index $k$ can change. There are exactly
\[
n-1
\]
such terms.

Thus,
\[
|U(X_1,\dots,X_n)-U(X_1,\dots,\widetilde{X}_k,\dots,X_n)|
\]
is bounded by
\[
\frac{1}{\binom{n}{2}}
\sum_{j\neq k}
|g(X_k,X_j)-g(\widetilde{X}_k,X_j)|.
\]

Since
\[
|g|\leq C,
\]
each difference satisfies
\[
|g(X_k,X_j)-g(\widetilde{X}_k,X_j)|
\leq
2C.
\]

Therefore,
\[
|U(X)-U(\widetilde{X})|
\leq
\frac{(n-1)2C}{\binom{n}{2}}
=
\frac{4C}{n}.
\]

Hence the bounded differences constants are
\[
c_k=\frac{4C}{n},
\qquad
k=1,\dots,n.
\]

Therefore,
\[
\sum_{k=1}^n c_k^2
=
n\left(\frac{4C}{n}\right)^2
=
\frac{16C^2}{n}.
\]

Applying McDiarmid's inequality,
\[
\Prob(|U-\E[U]|\geq t)
\leq
2\exp\left(
-\frac{2t^2}{16C^2/n}
\right).
\]

Simplifying,
\[
\Prob(|U-\E[U]|\geq t)
\leq
2\exp\left(
-\frac{nt^2}{8C^2}
\right).
\]
\end{proof}

\section{Lipschitz Functions of Gaussian Random Variables}

Next we study concentration of Lipschitz functions of Gaussian vectors.

\begin{definition}[Lipschitz Function]
A function
\[
f:\R^n\to\R
\]
is called $L$-Lipschitz if
\[
|f(x)-f(y)|
\leq
L\norm{x-y}_2
\]
for all
\[
x,y\in\R^n.
\]
\end{definition}

\begin{theorem}[Gaussian Concentration]
Let
\[
X=(X_1,\dots,X_n),
\qquad
X_i\iid \Normal(0,1).
\]

Suppose
\[
f:\R^n\to\R
\]
is $L$-Lipschitz. Then
\[
\Prob\left(
|f(X)-\E[f(X)]|
\geq t
\right)
\leq
2\exp\left(
-\frac{t^2}{2L^2}
\right).
\]
\end{theorem}

\begin{remark}
The lecture derived a slightly weaker constant:
\[
\Prob\left(
|f(X)-\E[f(X)]|
\geq t
\right)
\leq
2\exp\left(
-\frac{2t^2}{\pi^2L^2}
\right).
\]

The sharp constant requires deeper results from Gaussian isoperimetry.
\end{remark}

\subsection{A Key Gaussian Inequality}

Assume temporarily that $f$ is continuously differentiable.

\begin{lemma}
Let
\[
X,Y
\]
be independent standard Gaussian vectors in $\R^n$.

For every convex function
\[
\psi:\R\to\R,
\]
we have
\[
\E\left[
\psi(f(X)-\E[f(X)])
\right]
\leq
\E\left[
\psi\left(
\frac{\pi}{2}\inner{\grad f(X)}{Y}
\right)
\right].
\]
\end{lemma}

\begin{remark}
This result compares fluctuations of $f(X)$ to Gaussian linear fluctuations determined by the gradient.
\end{remark}

\subsection{Derivation of Gaussian Concentration}

We now show how the previous lemma implies sub-Gaussian concentration.

\begin{proof}
Apply the lemma with
\[
\psi(t)=e^{\lambda t}.
\]

Then
\[
\E\left[
e^{\lambda(f(X)-\E[f(X)])}
\right]
\leq
\E\left[
e^{\lambda \frac{\pi}{2}\inner{\grad f(X)}{Y}}
\right].
\]

Condition on $X$.

For fixed $X$,
\[
\inner{\grad f(X)}{Y}
\]
is Gaussian with mean zero and variance
\[
\norm{\grad f(X)}_2^2.
\]

Therefore,
\[
\E_Y\left[
e^{\lambda \frac{\pi}{2}\inner{\grad f(X)}{Y}}
\right]
=
\exp\left(
\frac{\lambda^2\pi^2}{8}
\norm{\grad f(X)}_2^2
\right).
\]

Since $f$ is $L$-Lipschitz,
\[
\norm{\grad f(X)}_2
\leq
L
\]
almost everywhere.

Hence,
\[
\E\left[
e^{\lambda(f(X)-\E[f(X)])}
\right]
\leq
\exp\left(
\frac{\lambda^2\pi^2L^2}{8}
\right).
\]

Thus,
\[
f(X)-\E[f(X)]
\]
is sub-Gaussian with parameter
\[
\frac{\pi L}{2}.
\]

Applying the standard sub-Gaussian tail bound gives
\[
\Prob\left(
|f(X)-\E[f(X)]|
\geq t
\right)
\leq
2\exp\left(
-\frac{2t^2}{\pi^2L^2}
\right).
\]
\end{proof}

\section{Proof of Gaussian Concentration}

In the previous lecture, we used the following claim to prove concentration of Lipschitz functions of Gaussian random vectors.

\begin{lemma}
Let
\[
X,Y\iid \Normal(0,I_n),
\]
and let
\[
\psi:\R\to\R
\]
be convex.

Then
\[
\E\left[
\psi(f(X)-\E[f(X)])
\right]
\leq
\E\left[
\psi\left(
\frac{\pi}{2}
\inner{\grad f(X)}{Y}
\right)
\right].
\]
\end{lemma}

\begin{proof}
Since
\[
Y\stackrel{d}{=}X,
\]
we have
\[
\E[f(Y)]
=
\E[f(X)].
\]

By Jensen's inequality,
\[
\psi(f(X)-\E[f(X)])
=
\psi\left(
f(X)-\E_Y[f(Y)]
\right)
\leq
\E_Y[\psi(f(X)-f(Y))].
\]

Taking expectation over $X$ gives
\[
\E[\psi(f(X)-\E[f(X)])]
\leq
\E_{X,Y}[\psi(f(X)-f(Y))].
\]

For
\[
\theta\in[0,\pi/2],
\]
define
\[
\gamma(\theta)
=
X\cos\theta+Y\sin\theta
\]
and
\[
\dot{\gamma}(\theta)
=
-X\sin\theta+Y\cos\theta.
\]

By the fundamental theorem of calculus,
\[
f(Y)-f(X)
=
\int_0^{\pi/2}
\frac{\dd}{\dd\theta}
f(\gamma(\theta))
\dd\theta.
\]

Using the chain rule,
\[
\frac{\dd}{\dd\theta}
f(\gamma(\theta))
=
\inner{
\grad f(\gamma(\theta))
}{
\dot{\gamma}(\theta)
}.
\]

Therefore,
\[
f(Y)-f(X)
=
\int_0^{\pi/2}
\inner{
\grad f(\gamma(\theta))
}{
\dot{\gamma}(\theta)
}
\dd\theta.
\]

Rewrite this as
\[
f(Y)-f(X)
=
\frac{2}{\pi}
\int_0^{\pi/2}
\frac{\pi}{2}
\inner{
\grad f(\gamma(\theta))
}{
\dot{\gamma}(\theta)
}
\dd\theta.
\]

Applying Jensen's inequality to the convex function $\psi$,
\[
\psi(f(Y)-f(X))
\leq
\frac{2}{\pi}
\int_0^{\pi/2}
\psi\left(
\frac{\pi}{2}
\inner{
\grad f(\gamma(\theta))
}{
\dot{\gamma}(\theta)
}
\right)
\dd\theta.
\]

Taking expectations,
\[
\E[\psi(f(Y)-f(X))]
\leq
\frac{2}{\pi}
\int_0^{\pi/2}
\E\left[
\psi\left(
\frac{\pi}{2}
\inner{
\grad f(\gamma(\theta))
}{
\dot{\gamma}(\theta)
}
\right)
\right]
\dd\theta.
\]

Now observe that
\[
(\gamma(\theta),\dot{\gamma}(\theta))
\]
has the same distribution as
\[
(X,Y).
\]

Indeed,
\[
\begin{pmatrix}
\gamma(\theta)
\\
\dot{\gamma}(\theta)
\end{pmatrix}
=
G_\theta
\begin{pmatrix}
X
\\
Y
\end{pmatrix},
\]
where
\[
G_\theta
=
\begin{pmatrix}
\cos\theta & \sin\theta
\\
-\sin\theta & \cos\theta
\end{pmatrix}
\]
is an orthogonal matrix.

Since Gaussian vectors are rotationally invariant,
\[
(\gamma(\theta),\dot{\gamma}(\theta))
\stackrel{d}{=}
(X,Y).
\]

Thus the integrand does not depend on $\theta$, so
\[
\E[\psi(f(Y)-f(X))]
\leq
\E\left[
\psi\left(
\frac{\pi}{2}
\inner{
\grad f(X)
}{Y}
\right)
\right].
\]

Combining with the earlier inequality yields
\[
\E[\psi(f(X)-\E[f(X)])]
\leq
\E\left[
\psi\left(
\frac{\pi}{2}
\inner{
\grad f(X)
}{Y}
\right)
\right].
\]
\end{proof}

\section{Order Statistics}

\begin{definition}[Order Statistics]
Given
\[
x_1,\dots,x_n\in\R,
\]
their order statistics are the reordered values
\[
x_{(1)}
\leq
x_{(2)}
\leq
\dots
\leq
x_{(n)}.
\]
\end{definition}

In particular,
\[
x_{(1)}=\min_i x_i,
\qquad
x_{(n)}=\max_i x_i.
\]

\begin{lemma}
For any
\[
x,y\in\R^n
\]
and every
\[
k\in\{1,\dots,n\},
\]
we have
\[
|x_{(k)}-y_{(k)}|
\leq
\norm{x-y}_2.
\]
\end{lemma}

\begin{proof}
The Euclidean norm dominates every coordinate difference:
\[
|x_i-y_i|
\leq
\norm{x-y}_2.
\]

Thus every coordinate of $x$ can move by at most
\[
\norm{x-y}_2.
\]

After reordering, the $k$-th order statistic can therefore move by at most the same amount:
\[
|x_{(k)}-y_{(k)}|
\leq
\norm{x-y}_2.
\]
\end{proof}

\begin{corollary}
Suppose
\[
X_1,\dots,X_n
\iid
\Normal(0,1).
\]

Then, for every $k$,
\[
\Prob\left(
|X_{(k)}-\E[X_{(k)}]|
\geq t
\right)
\leq
2e^{-t^2/2}.
\]
\end{corollary}

\begin{proof}
The map
\[
x\mapsto x_{(k)}
\]
is $1$-Lipschitz by the previous lemma.

Applying Gaussian concentration yields
\[
\Prob\left(
|X_{(k)}-\E[X_{(k)}]|
\geq t
\right)
\leq
2e^{-t^2/2}.
\]
\end{proof}

\section{Concentration of the Norm}

Random vectors in high dimensions behave very differently from low-dimensional intuition.

\subsection{Mass Concentration in the Hypercube}

Consider
\[
X\sim \mathrm{Unif}([0,1]^d).
\]

Fix
\[
\delta\in(0,1/2).
\]

The boundary shell of width $\delta$ is
\[
[0,1]^d
\setminus
[\delta,1-\delta]^d.
\]

Its probability mass is
\[
p_d
=
1-(1-2\delta)^d.
\]

For fixed $\delta>0$,
\[
(1-2\delta)^d\to0
\]
as
\[
d\to\infty.
\]

Therefore,
\[
p_d\to1.
\]

\begin{remark}
In high dimensions, most of the mass of the cube concentrates near the boundary.
\end{remark}

\section{Concentration of Euclidean Norms}

We now study concentration of
\[
\norm{X}_2
\]
for high-dimensional random vectors.

\begin{theorem}
Suppose
\[
X=(X_1,\dots,X_d)\in\R^d
\]
has independent coordinates satisfying:
\begin{itemize}
\item each $X_i$ is sub-Gaussian,
\item
\[
\E[X_i]=0,
\]
\item
\[
\E[X_i^2]=1.
\]
\end{itemize}

Let
\[
\sigma
=
\max_i \norm{X_i}_{\psi_2}.
\]

Then there exists a universal constant $c>0$ such that
\[
\Prob\left(
\left|
\norm{X}_2^2-d
\right|
\geq td
\right)
\leq
2\exp\left(
-cd\min\left(
\frac{t^2}{\sigma^4},
\frac{t}{\sigma^2}
\right)
\right).
\]
\end{theorem}

\begin{proof}
First observe that
\[
\E[\norm{X}_2^2]
=
\sum_{i=1}^d \E[X_i^2]
=
d.
\]

Thus,
\[
\frac{1}{d}\E[\norm{X}_2^2]=1.
\]

We previously proved that if
\[
Y,Z
\]
are sub-Gaussian, then
\[
YZ
\]
is sub-exponential.

Applying this with
\[
Y=Z=X_i,
\]
we conclude that
\[
X_i^2
\]
is sub-exponential.

Moreover,
\[
\norm{X_i^2}_{\psi_1}
\leq
C\norm{X_i}_{\psi_2}^2
\leq
C\sigma^2.
\]

Centering preserves the sub-exponential norm up to constants, so
\[
\norm{X_i^2-1}_{\psi_1}
\leq
C\sigma^2.
\]

Now write
\[
\frac{1}{d}\norm{X}_2^2-1
=
\frac{1}{d}
\sum_{i=1}^d (X_i^2-1).
\]

Since the variables
\[
X_i^2-1
\]
are independent centered sub-exponential random variables, Bernstein's inequality gives
\[
\Prob\left(
\left|
\frac{1}{d}\sum_{i=1}^d (X_i^2-1)
\right|
\geq t
\right)
\leq
2\exp\left(
-c
\min\left(
\frac{d^2t^2}{d\sigma^4},
\frac{dt}{\sigma^2}
\right)
\right).
\]

Simplifying,
\[
\Prob\left(
\left|
\frac{1}{d}\norm{X}_2^2-1
\right|
\geq t
\right)
\leq
2\exp\left(
-cd
\min\left(
\frac{t^2}{\sigma^4},
\frac{t}{\sigma^2}
\right)
\right).
\]

Equivalently,
\[
\Prob\left(
\left|
\norm{X}_2^2-d
\right|
\geq td
\right)
\leq
2\exp\left(
-cd
\min\left(
\frac{t^2}{\sigma^4},
\frac{t}{\sigma^2}
\right)
\right).
\]
\end{proof}

\begin{corollary}
Under the assumptions of the theorem,
there exists a universal constant $c>0$ such that
\[
\Prob\left(
\left|
\norm{X}_2-\sqrt{d}
\right|
\geq t\sqrt{d}
\right)
\leq
2\exp\left(
-\frac{cdt^2}{\sigma^4}
\right).
\]
\end{corollary}

\begin{proof}
Let
\[
z=\frac{\norm{X}_2}{\sqrt{d}}.
\]

Then
\[
|z-1|\geq t
\]
implies
\[
|z^2-1|
=
|z-1||z+1|.
\]

Since
\[
|z+1|\geq1
\]
and
\[
|z+1|\geq |z-1|,
\]
we obtain
\[
|z^2-1|
\geq
\max(t,t^2).
\]

Therefore,
\[
\Prob(|z-1|\geq t)
\leq
\Prob(|z^2-1|\geq \max(t,t^2)).
\]

Applying the previous theorem with
\[
\max(t,t^2)
\]
in place of $t$ yields
\[
\Prob(|z-1|\geq t)
\leq
2\exp\left(
-cd
\min\left(
\frac{t^2}{\sigma^4},
\frac{t}{\sigma^2}
\right)
\right).
\]

For small deviations,
\[
t\leq1,
\]
the quadratic term dominates, giving
\[
\Prob(|z-1|\geq t)
\leq
2\exp\left(
-\frac{cdt^2}{\sigma^4}
\right).
\]

Substituting
\[
z=\frac{\norm{X}_2}{\sqrt{d}}
\]
gives
\[
\Prob\left(
\left|
\norm{X}_2-\sqrt{d}
\right|
\geq t\sqrt{d}
\right)
\leq
2\exp\left(
-\frac{cdt^2}{\sigma^4}
\right).
\]
\end{proof}

\begin{remark}
In high dimensions,
\[
\norm{X}_2
\]
typically fluctuates only by constant-order deviations around
\[
\sqrt{d}.
\]

Thus high-dimensional random vectors are concentrated near a thin spherical shell.
\end{remark}

\chapter{Matrix Norms and Matrix Factorizations}

\section{Matrix Inner Products and Norms}

We now move from scalars and vectors to matrices. Many high-dimensional probability results are naturally expressed in terms of matrix norms.

\begin{definition}[Trace Inner Product]
For matrices
\[
A,B\in\R^{n\times m},
\]
the trace inner product is
\[
\inner{A}{B}
=
\tr(A^\top B).
\]
\end{definition}

The trace inner product induces the Frobenius norm.

\begin{definition}[Frobenius Norm]
For
\[
A\in\R^{n\times m},
\]
the Frobenius norm is
\[
\norm{A}_F
=
\sqrt{\inner{A}{A}}
=
\sqrt{\tr(A^\top A)}.
\]
\end{definition}

Equivalently,
\[
\norm{A}_F^2
=
\sum_{i=1}^n\sum_{j=1}^m A_{ij}^2.
\]

If $\operatorname{vec}(A)$ denotes the vector obtained by stacking the entries of $A$, then
\[
\inner{A}{B}
=
\operatorname{vec}(A)^\top \operatorname{vec}(B)
\]
and
\[
\norm{A}_F
=
\norm{\operatorname{vec}(A)}_2.
\]

\section{The Operator Norm}

A matrix
\[
A\in\R^{n\times m}
\]
can be viewed as a linear map
\[
A:\R^m\to\R^n.
\]

This induces another important norm.

\begin{definition}[Operator Norm]
The operator norm of
\[
A\in\R^{n\times m}
\]
is
\[
\norm{A}_{\mathrm{op}}
=
\sup_{\norm{v}_2=1}
\norm{Av}_2.
\]
\end{definition}

Equivalently,
\[
\norm{A}_{\mathrm{op}}
=
\sup_{v\neq 0}
\frac{\norm{Av}_2}{\norm{v}_2}.
\]

\begin{remark}
The Frobenius norm measures the overall size of all matrix entries, while the operator norm measures the maximum stretching effect of the linear map $x\mapsto Ax$.
\end{remark}

\section{Rotational Invariance of Matrix Norms}

\begin{proposition}
Let
\[
A\in\R^{n\times m},
\qquad
U\in O(n),
\qquad
V\in O(m),
\]
where $O(k)$ denotes the set of $k\times k$ orthogonal matrices. Then
\[
\norm{UAV}_F=\norm{A}_F
\]
and
\[
\norm{UAV}_{\mathrm{op}}=\norm{A}_{\mathrm{op}}.
\]
\end{proposition}

\begin{proof}
For the Frobenius norm,
\[
\norm{UAV}_F^2
=
\tr((UAV)^\top(UAV)).
\]

Expanding,
\[
\norm{UAV}_F^2
=
\tr(V^\top A^\top U^\top U A V).
\]

Since
\[
U^\top U=I,
\]
we get
\[
\norm{UAV}_F^2
=
\tr(V^\top A^\top A V).
\]

Using cyclic invariance of trace,
\[
\tr(V^\top A^\top A V)
=
\tr(A^\top A VV^\top).
\]

Since
\[
VV^\top=I,
\]
we obtain
\[
\norm{UAV}_F^2
=
\tr(A^\top A)
=
\norm{A}_F^2.
\]

For the operator norm,
\[
\norm{UAV}_{\mathrm{op}}
=
\sup_{\norm{x}_2=1}
\norm{UAVx}_2.
\]

Since $U$ is orthogonal,
\[
\norm{UAVx}_2
=
\norm{AVx}_2.
\]

As $x$ ranges over the unit sphere, so does $Vx$. Therefore,
\[
\norm{UAV}_{\mathrm{op}}
=
\sup_{\norm{z}_2=1}
\norm{Az}_2
=
\norm{A}_{\mathrm{op}}.
\]
\end{proof}

\section{Singular Value Decomposition}

Matrix factorizations are one of the main tools for understanding matrices geometrically and algebraically.

\begin{theorem}[Singular Value Decomposition]
Let
\[
A\in\R^{n\times m}.
\]

Then there exist orthogonal matrices
\[
U\in O(n),
\qquad
V\in O(m),
\]
and a diagonal matrix
\[
\Sigma\in\R^{n\times m}
\]
such that
\[
A=U\Sigma V^\top.
\]

The only possibly nonzero entries of $\Sigma$ are
\[
\sigma_1\geq \sigma_2\geq \cdots \geq \sigma_p\geq 0,
\qquad
p=\min\{n,m\}.
\]
\end{theorem}

\begin{definition}[Singular Values and Singular Vectors]
The numbers
\[
\sigma_1,\dots,\sigma_p
\]
are called the singular values of $A$.

The columns
\[
u_1,\dots,u_n
\]
of $U$ are called the left singular vectors, and the columns
\[
v_1,\dots,v_m
\]
of $V$ are called the right singular vectors.
\end{definition}

\begin{remark}
The singular value decomposition says that, after appropriate orthogonal changes of coordinates in the domain and codomain, the map $x\mapsto Ax$ acts like a coordinate-wise scaling.
\end{remark}

\subsection{Proof of the Singular Value Decomposition}

\begin{proof}
Consider the optimization problem
\[
\sigma_1
=
\max_{\norm{u}_2=1,\ \norm{v}_2=1}
u^\top A v.
\]

The objective function
\[
(u,v)\mapsto u^\top A v
\]
is continuous, and the product of unit spheres
\[
S^{n-1}\times S^{m-1}
\]
is compact. Hence a maximizer exists.

Let
\[
(u_1,v_1)
\]
be a maximizer, and define
\[
\sigma_1=u_1^\top A v_1.
\]

We claim that
\[
Av_1=\sigma_1u_1.
\]

Indeed, for fixed $v_1$, the quantity
\[
u^\top A v_1
\]
is maximized over $\norm{u}_2=1$ when $u$ is in the direction of $Av_1$. Therefore,
\[
u_1=\frac{Av_1}{\norm{Av_1}_2}
\]
whenever $Av_1\neq0$, and consequently
\[
Av_1=\norm{Av_1}_2u_1.
\]

Moreover,
\[
\sigma_1
=
u_1^\top Av_1
=
\norm{Av_1}_2.
\]

Hence
\[
Av_1=\sigma_1u_1.
\]

Complete $u_1$ to an orthonormal basis of $\R^n$:
\[
u_1,\widetilde{u}_2,\dots,\widetilde{u}_n,
\]
and complete $v_1$ to an orthonormal basis of $\R^m$:
\[
v_1,\widetilde{v}_2,\dots,\widetilde{v}_m.
\]

Let
\[
\widetilde{U}
=
\begin{pmatrix}
u_1 & \widetilde{u}_2 & \cdots & \widetilde{u}_n
\end{pmatrix}
\]
and
\[
\widetilde{V}
=
\begin{pmatrix}
v_1 & \widetilde{v}_2 & \cdots & \widetilde{v}_m
\end{pmatrix}.
\]

Then
\[
\widetilde{U}^\top A\widetilde{V}
=
\begin{pmatrix}
\sigma_1 & w^\top \\
0 & B
\end{pmatrix}
\]
for some vector $w$ and matrix $B$.

We now show that
\[
w=0.
\]

Suppose not. Then there exists a unit vector $q$ such that
\[
w^\top q=\norm{w}_2.
\]

Consider the unit vector
\[
v_\theta
=
\cos\theta\, e_1+\sin\theta
\begin{pmatrix}
0\\
q
\end{pmatrix}.
\]

Then
\[
\norm{v_\theta}_2=1.
\]

Using the block form,
\[
\widetilde{U}^\top A\widetilde{V}v_\theta
=
\begin{pmatrix}
\sigma_1\cos\theta+\norm{w}_2\sin\theta\\
Bq\sin\theta
\end{pmatrix}.
\]

For small positive $\theta$, the first coordinate is larger than $\sigma_1$, contradicting the maximality of $\sigma_1$ as the maximum value of
\[
\norm{Av}_2
\]
over unit vectors $v$.

Therefore,
\[
w=0.
\]

Hence
\[
\widetilde{U}^\top A\widetilde{V}
=
\begin{pmatrix}
\sigma_1 & 0 \\
0 & B
\end{pmatrix}.
\]

Applying the same argument inductively to $B$ yields orthogonal matrices $U$ and $V$ such that
\[
U^\top A V=\Sigma,
\]
where $\Sigma$ is diagonal with nonnegative diagonal entries ordered as
\[
\sigma_1\geq \sigma_2\geq\cdots\geq\sigma_p\geq0.
\]

Thus,
\[
A=U\Sigma V^\top.
\]
\end{proof}

\section{Properties of the Singular Value Decomposition}

Let
\[
A=U\Sigma V^\top
\]
be the singular value decomposition of
\[
A\in\R^{n\times m}.
\]

Let
\[
u_i
\]
denote the columns of $U$, and let
\[
v_i
\]
denote the columns of $V$.

\begin{proposition}[Basic SVD Properties]
The following properties hold.

\begin{enumerate}
\item For every $i$,
\[
Av_i=\sigma_i u_i
\]
and
\[
A^\top u_i=\sigma_i v_i.
\]

\item For every $x\in\R^m$,
\[
\sigma_{\min}(A)\norm{x}_2
\leq
\norm{Ax}_2
\leq
\sigma_{\max}(A)\norm{x}_2,
\]
where
\[
\sigma_{\max}(A)=\sigma_1.
\]

If $A$ is rank-deficient, then $\sigma_{\min}(A)$ is interpreted as the smallest singular value among the relevant domain directions, and the lower bound is zero on the nullspace.

\item If
\[
r=\rank(A),
\]
then
\[
\sigma_r(A)>0
\]
and
\[
\sigma_{r+1}(A)=0.
\]

\item The nullspace and range are
\[
\operatorname{null}(A)
=
\operatorname{span}\{v_{r+1},\dots,v_m\},
\]
and
\[
\operatorname{range}(A)
=
\operatorname{span}\{u_1,\dots,u_r\}.
\]

\item The matrix $A$ can be written as
\[
A
=
\sum_{i=1}^r
\sigma_i u_i v_i^\top.
\]

\item The Frobenius and operator norms satisfy
\[
\norm{A}_F
=
\left(
\sum_{i=1}^r \sigma_i^2
\right)^{1/2}
\]
and
\[
\norm{A}_{\mathrm{op}}
=
\sigma_1.
\]
\end{enumerate}
\end{proposition}

\begin{proof}
Since
\[
A=U\Sigma V^\top,
\]
we have
\[
AV=U\Sigma.
\]

Comparing columns gives
\[
Av_i=\sigma_i u_i.
\]

Similarly,
\[
A^\top
=
V\Sigma^\top U^\top,
\]
so
\[
A^\top U=V\Sigma^\top.
\]

Comparing columns gives
\[
A^\top u_i=\sigma_i v_i.
\]

For the norm bounds, write
\[
x=Vz
\]
for some $z\in\R^m$ with
\[
\norm{z}_2=\norm{x}_2.
\]

Then
\[
Ax
=
U\Sigma z.
\]

Since $U$ is orthogonal,
\[
\norm{Ax}_2
=
\norm{\Sigma z}_2.
\]

Because $\Sigma$ scales coordinates by the singular values,
\[
\norm{\Sigma z}_2^2
=
\sum_i \sigma_i^2 z_i^2.
\]

Therefore,
\[
\norm{\Sigma z}_2
\leq
\sigma_{\max}(A)\norm{z}_2
=
\sigma_{\max}(A)\norm{x}_2.
\]

The lower bound follows similarly on the subspace orthogonal to the nullspace. If $A$ has a nontrivial nullspace, the global lower bound is zero.

The rank statement follows because the rank of $\Sigma$ equals the number of nonzero diagonal entries, and orthogonal multiplication does not change rank.

The nullspace statement follows from
\[
Ax=0
\iff
\Sigma V^\top x=0.
\]

Thus $V^\top x$ must be supported only on coordinates corresponding to zero singular values. Hence
\[
\operatorname{null}(A)
=
\operatorname{span}\{v_{r+1},\dots,v_m\}.
\]

The range statement follows from
\[
Ax=U\Sigma V^\top x,
\]
so every vector in the range is a linear combination of
\[
u_1,\dots,u_r.
\]

The rank-one expansion follows from writing
\[
U\Sigma V^\top
=
\sum_{i=1}^r
\sigma_i u_i v_i^\top.
\]

Finally, rotational invariance of the Frobenius and operator norms gives
\[
\norm{A}_F=\norm{\Sigma}_F
=
\left(
\sum_{i=1}^r\sigma_i^2
\right)^{1/2},
\]
and
\[
\norm{A}_{\mathrm{op}}
=
\norm{\Sigma}_{\mathrm{op}}
=
\sigma_1.
\]
\end{proof}

\section{Spectral Decomposition of Symmetric Matrices}

For symmetric matrices, there is a more specialized factorization.

Let
\[
\mathbb{S}^n
=
\{A\in\R^{n\times n}:A=A^\top\}
\]
denote the set of real symmetric matrices.

\begin{theorem}[Spectral Decomposition]
Let
\[
A\in\mathbb{S}^n.
\]

Then there exists an orthogonal matrix
\[
U\in O(n)
\]
and a diagonal matrix
\[
\Lambda
=
\diag(\lambda_1,\dots,\lambda_n)
\]
such that
\[
A=U\Lambda U^\top.
\]
\end{theorem}

\begin{proof}
We prove the result using two facts:

\begin{enumerate}
\item Every real symmetric matrix has real eigenvalues.
\item Eigenvectors corresponding to distinct eigenvalues are orthogonal.
\end{enumerate}

Let
\[
\lambda_1
\]
be an eigenvalue of $A$ with unit eigenvector $u_1$:
\[
Au_1=\lambda_1u_1.
\]

Let
\[
W_1=\{u_1\}^{\perp}.
\]

We claim that $W_1$ is invariant under $A$. If $w\in W_1$, then
\[
\inner{Aw}{u_1}
=
\inner{w}{A^\top u_1}
=
\inner{w}{Au_1}
=
\inner{w}{\lambda_1u_1}
=
0.
\]

Thus
\[
Aw\in W_1.
\]

Therefore, the restriction of $A$ to $W_1$ is again symmetric. Repeating the argument inductively, we obtain an orthonormal basis
\[
u_1,\dots,u_n
\]
of eigenvectors of $A$ with corresponding real eigenvalues
\[
\lambda_1,\dots,\lambda_n.
\]

Let
\[
U=
\begin{pmatrix}
u_1 & \cdots & u_n
\end{pmatrix}.
\]

Then $U$ is orthogonal and
\[
AU
=
U\Lambda.
\]

Multiplying by $U^\top$ on the right gives
\[
A
=
U\Lambda U^\top.
\]
\end{proof}

\begin{remark}
For symmetric matrices, singular values are the absolute values of eigenvalues:
\[
\sigma_i(A)=|\lambda_i(A)|
\]
after appropriate ordering.
\end{remark}

\chapter{Low-Rank Approximation and Perturbation Theory}

\section{Low-Rank Approximation}

Singular value and spectral decompositions are useful because they allow us to construct low-rank approximations of matrices.

\begin{definition}[Best Rank-$k$ Approximation]
Let
\[
A\in\R^{n\times m}.
\]
The best rank-$k$ approximation of $A$ in operator norm is
\[
A_{[k]}
\in
\argmin_{\rank(B)\leq k}
\norm{A-B}_{\mathrm{op}}.
\]
\end{definition}

Low-rank approximations are useful in data science for several reasons:
\begin{itemize}
\item Low-rank matrices require fewer parameters to store.
\item Matrix-vector products can be computed more quickly.
\item Many real datasets are approximately low-rank.
\end{itemize}

For example, if
\[
A\approx UV^\top
\]
where
\[
U\in\R^{d\times r},
\qquad
V\in\R^{d\times r},
\]
then storage requires approximately
\[
2dr
\]
parameters rather than
\[
d^2.
\]

Similarly, multiplying
\[
A\approx UV^\top
\]
by a vector can be done as
\[
Ax\approx U(V^\top x),
\]
which costs
\[
O(dr)
\]
rather than
\[
O(d^2).
\]

\section{The Eckart--Young Theorem}

Let
\[
A=U\Sigma V^\top
\]
be the singular value decomposition of $A$, with singular values
\[
\sigma_1\geq \sigma_2\geq \cdots \geq \sigma_r>0,
\]
where
\[
r=\rank(A).
\]

Define
\[
A_{[k]}
=
\sum_{i=1}^k \sigma_i u_i v_i^\top.
\]

\begin{theorem}[Eckart--Young Theorem: Operator Norm]
Let
\[
A\in\R^{n\times m}
\]
and let
\[
k\leq \rank(A).
\]

Then
\[
A_{[k]}
=
\sum_{i=1}^k \sigma_i u_i v_i^\top
\]
has rank at most $k$, and
\[
\min_{\rank(B)\leq k}
\norm{A-B}_{\mathrm{op}}
=
\norm{A-A_{[k]}}_{\mathrm{op}}
=
\sigma_{k+1}.
\]
\end{theorem}

\begin{proof}
Using the singular value decomposition,
\[
A=U\Sigma V^\top.
\]

By rotational invariance of the operator norm,
\[
\norm{A-A_{[k]}}_{\mathrm{op}}
=
\norm{\Sigma-\Sigma_{[k]}}_{\mathrm{op}},
\]
where
\[
\Sigma_{[k]}
=
\diag(\sigma_1,\dots,\sigma_k,0,\dots,0).
\]

Therefore,
\[
\norm{A-A_{[k]}}_{\mathrm{op}}
=
\sigma_{k+1}.
\]

It remains to show that no rank-$k$ matrix can do better.

Let $B$ be any matrix with
\[
\rank(B)\leq k.
\]

Then
\[
\dim(\operatorname{null}(B))\geq m-k.
\]

Also,
\[
\dim(\operatorname{span}\{v_1,\dots,v_{k+1}\})=k+1.
\]

By dimension counting,
\[
\operatorname{null}(B)
\cap
\operatorname{span}\{v_1,\dots,v_{k+1}\}
\neq
\{0\}.
\]

Choose a unit vector
\[
z\in
\operatorname{null}(B)
\cap
\operatorname{span}\{v_1,\dots,v_{k+1}\}.
\]

Then
\[
Bz=0.
\]

Hence,
\[
\norm{A-B}_{\mathrm{op}}
\geq
\norm{(A-B)z}_2
=
\norm{Az}_2.
\]

Write
\[
z=\sum_{i=1}^{k+1} a_i v_i,
\qquad
\sum_{i=1}^{k+1} a_i^2=1.
\]

Then
\[
Az
=
\sum_{i=1}^{k+1} a_i\sigma_i u_i.
\]

Therefore,
\[
\norm{Az}_2^2
=
\sum_{i=1}^{k+1} a_i^2\sigma_i^2.
\]

Since
\[
\sigma_i\geq \sigma_{k+1}
\]
for
\[
i\leq k+1,
\]
we obtain
\[
\norm{Az}_2^2
\geq
\sigma_{k+1}^2
\sum_{i=1}^{k+1}a_i^2
=
\sigma_{k+1}^2.
\]

Thus,
\[
\norm{A-B}_{\mathrm{op}}
\geq
\sigma_{k+1}.
\]

Since $B$ was arbitrary,
\[
\min_{\rank(B)\leq k}
\norm{A-B}_{\mathrm{op}}
\geq
\sigma_{k+1}.
\]

Together with the construction of $A_{[k]}$, this proves the result.
\end{proof}

\begin{theorem}[Eckart--Young Theorem: Frobenius Norm]
Let
\[
A\in\R^{n\times m}
\]
and let
\[
k\leq \rank(A).
\]

Then
\[
A_{[k]}
=
\sum_{i=1}^k \sigma_i u_i v_i^\top
\]
solves
\[
\min_{\rank(B)\leq k}
\norm{A-B}_F,
\]
and
\[
\min_{\rank(B)\leq k}
\norm{A-B}_F
=
\norm{A-A_{[k]}}_F
=
\left(
\sum_{i=k+1}^r \sigma_i^2
\right)^{1/2}.
\]
\end{theorem}

\begin{proof}
By rotational invariance of the Frobenius norm,
\[
\norm{A-A_{[k]}}_F
=
\norm{\Sigma-\Sigma_{[k]}}_F.
\]

Since
\[
\Sigma-\Sigma_{[k]}
\]
has diagonal entries
\[
0,\dots,0,\sigma_{k+1},\dots,\sigma_r,
\]
we get
\[
\norm{A-A_{[k]}}_F^2
=
\sum_{i=k+1}^r \sigma_i^2.
\]

The optimality proof follows from the same singular-vector dimension-counting argument as in the operator-norm case, with the Frobenius norm measuring the squared sum of discarded singular values.
\end{proof}

\section{Symmetric Low-Rank Approximation}

If
\[
A\in\mathbb{S}^n
\]
has spectral decomposition
\[
A=U\Lambda U^\top,
\]
where
\[
\Lambda=\diag(\lambda_1,\dots,\lambda_n),
\]
then a rank-$k$ approximation is obtained by keeping the $k$ eigenvalues of largest magnitude.

Let
\[
|\lambda_{i_1}|\geq |\lambda_{i_2}|\geq \cdots \geq |\lambda_{i_n}|.
\]

Then
\[
A_{[k]}
=
\sum_{\ell=1}^k
\lambda_{i_\ell}u_{i_\ell}u_{i_\ell}^\top.
\]

\begin{remark}
For symmetric matrices, the singular values are the absolute values of the eigenvalues. Therefore, best low-rank approximation keeps the spectral components corresponding to the largest absolute eigenvalues.
\end{remark}

\section{Perturbation Theory}

In applications, matrices are often only approximately low-rank. A common model is
\[
M=A+E,
\]
where
\[
A
\]
is the structured or low-rank signal matrix and
\[
E
\]
is a perturbation or noise matrix.

This leads to the following questions:
\begin{itemize}
\item How far are the eigenvalues of $M$ from those of $A$?
\item How far are the singular values of $M$ from those of $A$?
\item How far are the eigenspaces or singular subspaces of $M$ from those of $A$?
\end{itemize}

\section{Weyl's Inequality}

\begin{lemma}[Weyl's Inequality for Eigenvalues]
Let
\[
A,E\in\mathbb{S}^n.
\]

Then, for every
\[
i\in\{1,\dots,n\},
\]
\[
|\lambda_i(A)-\lambda_i(A+E)|
\leq
\norm{E}_{\mathrm{op}}.
\]
\end{lemma}

\begin{lemma}[Weyl's Inequality for Singular Values]
Let
\[
A,E\in\R^{n\times m}.
\]

Then, for every
\[
i\in\{1,\dots,\min(n,m)\},
\]
\[
|\sigma_i(A)-\sigma_i(A+E)|
\leq
\norm{E}_{\mathrm{op}}.
\]
\end{lemma}

\begin{remark}
Eigenvalues and singular values are Lipschitz with respect to perturbations in operator norm. Thus small operator-norm perturbations cannot significantly change the spectrum.
\end{remark}

\section{Distances Between Subspaces}

When studying perturbations of eigenvectors, it is often better to compare subspaces rather than individual vectors.

This is necessary because:
\begin{itemize}
\item eigenvectors are only defined up to sign;
\item repeated eigenvalues generate eigenspaces rather than individual distinguished eigenvectors;
\item different orthonormal bases can span the same subspace.
\end{itemize}

Let
\[
\mathcal{U}
=
\operatorname{span}(U)
\]
and
\[
\mathcal{U}_*
=
\operatorname{span}(U_*),
\]
where
\[
U,U_*\in\R^{n\times r}
\]
have orthonormal columns.

Let
\[
U_\perp
\]
and
\[
U_{*,\perp}
\]
be orthonormal bases for the orthogonal complements, so that
\[
[U\ \ U_\perp]\in O(n)
\]
and
\[
[U_*\ \ U_{*,\perp}]\in O(n).
\]

\subsection{Why Direct Basis Comparison Fails}

A naive attempt to measure distance between subspaces is
\[
\norm{U-U_*}.
\]

This is not a good subspace distance because the same subspace can be represented by many different orthonormal bases. If
\[
R\in O(r),
\]
then
\[
UR
\]
spans the same subspace as $U$, but
\[
\norm{UR-U_*}
\]
can differ from
\[
\norm{U-U_*}.
\]

Thus a useful subspace distance should be invariant under changes of basis.

\section{Subspace Distance Metrics}

\subsection{Distance with Optimal Rotation}

One option is to define
\[
\operatorname{dist}(\mathcal{U},\mathcal{U}_*)
=
\min_{R\in O(r)}
\norm{UR-U_*},
\]
where the norm may be either the Frobenius norm or the operator norm.

This removes the arbitrary choice of basis by optimally rotating one basis toward the other.

\subsection{Distance Between Projections}

Another basis-invariant choice is
\[
\norm{UU^\top-U_*U_*^\top}.
\]

The matrix
\[
UU^\top
\]
is the orthogonal projector onto
\[
\mathcal{U}.
\]

Thus this distance compares the subspaces themselves rather than the chosen bases.

\subsection{Principal Angles}

A third approach uses principal angles.

Let
\[
\sigma_1\geq \cdots \geq \sigma_r
\]
be the singular values of
\[
U^\top U_*.
\]

Since
\[
\norm{U^\top U_*}_{\mathrm{op}}
\leq
\norm{U}_{\mathrm{op}}\norm{U_*}_{\mathrm{op}}
=
1,
\]
we have
\[
\sigma_i\in[0,1].
\]

Define the principal angles
\[
\theta_i=\arccos(\sigma_i),
\qquad
i=1,\dots,r.
\]

Let
\[
\Theta
=
\diag(\theta_1,\dots,\theta_r),
\]
and define
\[
\sin\Theta
=
\diag(\sin\theta_1,\dots,\sin\theta_r).
\]

A common subspace distance is then
\[
\norm{\sin\Theta}.
\]

\begin{remark}
Principal angles measure how aligned two subspaces are. If the subspaces are identical, then all $\theta_i=0$. If they contain orthogonal directions, some angles approach $\pi/2$.
\end{remark}

\section{A Useful Rotation-Invariant Norm Fact}

\begin{lemma}
Suppose
\[
\norm{\cdot}
\]
is invariant under orthogonal transformations. Then, for compatible matrices $A$ and $B$,
\[
\norm{A}\sigma_{\min}(B)
\leq
\norm{AB}
\leq
\norm{A}\norm{B}_{\mathrm{op}},
\]
and
\[
\sigma_{\min}(A)\norm{B}
\leq
\norm{AB}
\leq
\norm{A}_{\mathrm{op}}\norm{B}.
\]
\end{lemma}

\begin{remark}
This fact is useful for proving equivalences between different subspace distances and for bounding matrix products when one factor is measured by singular values.
\end{remark}

\section{Basis Invariance of Subspace Distances}

\begin{proposition}
The following quantities do not depend on the choice of orthonormal bases for
\[
\mathcal{U}
\]
and
\[
\mathcal{U}_*:
\]
\begin{enumerate}
\item
\[
\min_{R\in O(r)}
\norm{UR-U_*};
\]

\item
\[
\norm{UU^\top-U_*U_*^\top};
\]

\item
\[
\norm{\sin\Theta}.
\]
\end{enumerate}
\end{proposition}

\begin{proof}
Suppose we replace
\[
U
\]
by
\[
UQ
\]
and
\[
U_*
\]
by
\[
U_*Q_*,
\]
where
\[
Q,Q_*\in O(r).
\]

For the optimal rotation distance,
\[
\min_{R\in O(r)}
\norm{UQR-U_*Q_*}.
\]

Let
\[
\widetilde{R}=QRQ_*^\top.
\]

As $R$ ranges over $O(r)$, so does $\widetilde{R}$. By rotational invariance of the norm, the minimized value is unchanged.

For the projection distance,
\[
(UQ)(UQ)^\top
=
UQQ^\top U^\top
=
UU^\top.
\]

Similarly,
\[
(U_*Q_*)(U_*Q_*)^\top
=
U_*U_*^\top.
\]

Thus the projection distance is unchanged.

For principal angles,
\[
(UQ)^\top(U_*Q_*)
=
Q^\top U^\top U_* Q_*.
\]

Multiplication on the left and right by orthogonal matrices does not change singular values. Therefore the singular values, and hence the principal angles, are unchanged.
\end{proof}

\section{Principal Angles Between Subspaces}

Recall that
\[
\mathcal{U}=\operatorname{span}(U),
\qquad
\mathcal{U}_*=\operatorname{span}(U_*),
\]
where
\[
U,U_*\in\R^{n\times r}
\]
have orthonormal columns.

The singular values of
\[
U^\top U_*
\]
play a central role in measuring distances between subspaces.

Let
\[
\sigma_1\geq \sigma_2\geq \cdots \geq \sigma_r\geq 0
\]
be the singular values of
\[
U^\top U_*.
\]

Since
\[
\norm{U^\top U_*}_{\mathrm{op}}
\leq
\norm{U}_{\mathrm{op}}
\norm{U_*}_{\mathrm{op}}
=
1,
\]
we have
\[
\sigma_i\in[0,1].
\]

\begin{definition}[Principal Angles]
The principal angles between
\[
\mathcal{U}
\]
and
\[
\mathcal{U}_*
\]
are defined by
\[
\theta_i
=
\arccos(\sigma_i),
\qquad
i=1,\dots,r.
\]
\end{definition}

Define
\[
\Theta
=
\diag(\theta_1,\dots,\theta_r)
\]
and
\[
\sin\Theta
=
\diag(\sin\theta_1,\dots,\sin\theta_r).
\]

A common subspace distance is
\[
\norm{\sin\Theta}.
\]

\subsection{Geometric Interpretation}

For one-dimensional subspaces, the angle is uniquely defined.

If
\[
u,u_*\in\R^n
\]
are unit vectors spanning one-dimensional subspaces, then
\[
\cos\theta
=
|\inner{u}{u_*}|.
\]

For higher-dimensional subspaces, there is no single canonical angle. Instead, one defines a sequence of principal angles.

The first principal angle is
\[
\theta_1
=
\min_{\substack{
u\in \mathcal{U}\cap S^{n-1}\\
v\in \mathcal{U}_*\cap S^{n-1}
}}
\arccos(\inner{u}{v}).
\]

Equivalently,
\[
\theta_1
=
\arccos
\left(
\max_{\substack{
u\in \mathcal{U}\cap S^{n-1}\\
v\in \mathcal{U}_*\cap S^{n-1}
}}
\inner{u}{v}
\right).
\]

If
\[
u_1,v_1
\]
are maximizers, then the second principal angle is defined recursively by restricting to orthogonal complements:
\[
\theta_2
=
\min_{\substack{
u\in \mathcal{U}\cap \operatorname{span}(u_1)^\perp\cap S^{n-1}\\
v\in \mathcal{U}_*\cap \operatorname{span}(v_1)^\perp\cap S^{n-1}
}}
\arccos(\inner{u}{v}).
\]

Continuing inductively yields
\[
\theta_1\leq \theta_2\leq \cdots \leq \theta_r.
\]

\begin{proposition}
The vectors
\[
u_i,v_i
\]
obtained from the recursive principal-angle construction are precisely the left and right singular vectors of
\[
U^\top U_*.
\]
\end{proposition}

\begin{proof}
The first principal angle satisfies
\[
\cos\theta_1
=
\max_{\substack{
u\in\mathcal{U}\cap S^{n-1}\\
v\in\mathcal{U}_*\cap S^{n-1}
}}
\inner{u}{v}.
\]

Write
\[
u=Ua,
\qquad
v=U_*b,
\]
where
\[
\norm{a}_2=\norm{b}_2=1.
\]

Then
\[
\inner{u}{v}
=
a^\top U^\top U_* b.
\]

Thus,
\[
\cos\theta_1
=
\max_{\norm{a}_2=\norm{b}_2=1}
a^\top(U^\top U_*)b.
\]

By the variational characterization of singular values, this maximum equals the top singular value of
\[
U^\top U_*.
\]

The maximizing vectors are the corresponding singular vectors. Repeating the argument inductively on orthogonal complements yields the remaining singular vectors and principal angles.
\end{proof}

\section{Equivalence of Subspace Distances}

We previously introduced several notions of distance between subspaces:
\begin{enumerate}
\item optimal rotation distance;
\item projection distance;
\item principal-angle distance.
\end{enumerate}

These distances are all equivalent up to constants.

\begin{lemma}
Let
\[
U,U_*\in\R^{n\times r}
\]
have orthonormal columns. Then
\[
\norm{UU^\top-U_*U_*^\top}_{\mathrm{op}}
=
\norm{\sin\Theta}_{\mathrm{op}}
=
\norm{U_\perp^\top U_*}_{\mathrm{op}}
=
\norm{U^\top U_{*,\perp}}_{\mathrm{op}}.
\]
\end{lemma}

\begin{proof}
We prove the identity
\[
\norm{UU^\top-U_*U_*^\top}_{\mathrm{op}}
=
\norm{\sin\Theta}_{\mathrm{op}}.
\]

Choose orthonormal bases so that
\[
U^\top U_*
=
\diag(\cos\theta_1,\dots,\cos\theta_r).
\]

In this basis, the projection matrices decompose into independent $2\times2$ blocks corresponding to each principal angle.

For a single angle $\theta$, the associated block difference has eigenvalues
\[
\pm \sin\theta.
\]

Therefore, the operator norm equals the largest singular value:
\[
\norm{UU^\top-U_*U_*^\top}_{\mathrm{op}}
=
\max_i \sin\theta_i
=
\norm{\sin\Theta}_{\mathrm{op}}.
\]

The identities involving
\[
U_\perp^\top U_*
\]
and
\[
U^\top U_{*,\perp}
\]
follow from the block orthogonal decomposition
\[
[U\ \ U_\perp]^\top
[U_*\ \ U_{*,\perp}]
\in O(n),
\]
whose off-diagonal blocks encode the sines of the principal angles.
\end{proof}

\begin{lemma}
Let
\[
U,U_*\in\R^{n\times r}
\]
have orthonormal columns. Then
\[
\frac{1}{\sqrt{2}}
\norm{UU^\top-U_*U_*^\top}_F
\leq
\min_{R\in O(r)}
\norm{UR-U_*}_F
\leq
\norm{UU^\top-U_*U_*^\top}_F.
\]
\end{lemma}

\begin{proof}
Let
\[
U^\top U_*
=
P\Sigma Q^\top
\]
be the singular value decomposition.

Choose
\[
R=PQ^\top.
\]

Then
\[
U^\top U_*R
=
P\Sigma P^\top.
\]

Using orthogonal invariance,
\[
\norm{UR-U_*}_F^2
=
2r-2\tr(R^\top U^\top U_*).
\]

Since
\[
R=PQ^\top,
\]
we obtain
\[
\tr(R^\top U^\top U_*)
=
\tr(\Sigma)
=
\sum_{i=1}^r \cos\theta_i.
\]

Therefore,
\[
\norm{UR-U_*}_F^2
=
2\sum_{i=1}^r (1-\cos\theta_i).
\]

On the other hand,
\[
\norm{UU^\top-U_*U_*^\top}_F^2
=
2\sum_{i=1}^r \sin^2\theta_i.
\]

Using
\[
1-\cos\theta
\leq
\sin^2\theta
\leq
2(1-\cos\theta),
\]
we obtain
\[
\frac{1}{2}
\norm{UU^\top-U_*U_*^\top}_F^2
\leq
\norm{UR-U_*}_F^2
\leq
\norm{UU^\top-U_*U_*^\top}_F^2.
\]

Taking square roots yields the result.
\end{proof}

\begin{remark}
These equivalences imply that controlling any one of the subspace distances automatically controls the others.
\end{remark}

\section{Perturbation of Eigenvectors}

Unlike eigenvalues, eigenvectors are not always stable under perturbation.

\subsection{A Warning Example}

Consider
\[
M_*=
\begin{pmatrix}
1+\varepsilon & 0\\
0 & 1-\varepsilon
\end{pmatrix},
\qquad
E=
\begin{pmatrix}
-\varepsilon & \varepsilon\\
\varepsilon & \varepsilon
\end{pmatrix},
\]
and define
\[
M=M_*+E
=
\begin{pmatrix}
1 & \varepsilon\\
\varepsilon & 1
\end{pmatrix}.
\]

The leading eigenvector of
\[
M_*
\]
is
\[
u_1^*=
\begin{pmatrix}
1\\
0
\end{pmatrix},
\]
while the leading eigenvector of
\[
M
\]
is
\[
u_1
=
\frac{1}{\sqrt{2}}
\begin{pmatrix}
1\\
1
\end{pmatrix}.
\]

Thus the leading eigenvectors remain far apart even as
\[
\varepsilon\to0.
\]

The issue is that the top two eigenvalues of
\[
M_*
\]
are nearly equal:
\[
\lambda_1(M_*)-\lambda_2(M_*)
=
2\varepsilon.
\]

Therefore, eigenvectors are stable only when the relevant eigenvalues are sufficiently separated.

\section{The Davis--Kahan $\sin\Theta$ Theorem}

Let
\[
M_*=M-E\in\mathbb{S}^n.
\]

Suppose
\[
M_*
=
[U_*\ \ U_{*,\perp}]
\begin{pmatrix}
\Lambda_* & 0\\
0 & \Lambda_{*,\perp}
\end{pmatrix}
\begin{pmatrix}
U_*^\top\\
U_{*,\perp}^\top
\end{pmatrix},
\]
and similarly
\[
M
=
[U\ \ U_\perp]
\begin{pmatrix}
\Lambda & 0\\
0 & \Lambda_\perp
\end{pmatrix}
\begin{pmatrix}
U^\top\\
U_\perp^\top
\end{pmatrix}.
\]

\begin{theorem}[Davis--Kahan $\sin\Theta$ Theorem]
Suppose there exist
\[
a\leq b
\]
and
\[
\Delta>0
\]
such that one of the following holds:

\begin{enumerate}
\item
\[
\{\lambda_1^*,\dots,\lambda_r^*\}
\subseteq
[a,b]
\]
and
\[
\{\lambda_{r+1},\dots,\lambda_n\}
\subseteq
(-\infty,a-\Delta]
\cup
[b+\Delta,\infty).
\]

\item
\[
\{\lambda_1,\dots,\lambda_r\}
\subseteq
[a,b]
\]
and
\[
\{\lambda_{r+1}^*,\dots,\lambda_n^*\}
\subseteq
(-\infty,a-\Delta]
\cup
[b+\Delta,\infty).
\]
\end{enumerate}

Then
\[
\norm{\sin\Theta}_{\mathrm{op}}
\leq
\frac{\norm{E}_{\mathrm{op}}}{\Delta}
\]
and
\[
\norm{\sin\Theta}_F
\leq
\frac{\sqrt{r}\norm{E}_{\mathrm{op}}}{\Delta}.
\]
\end{theorem}

\begin{remark}
The theorem states that eigenspaces are stable under perturbation provided there is a spectral gap
\[
\Delta.
\]
The perturbation size is controlled by
\[
\norm{E}_{\mathrm{op}}.
\]
\end{remark}

\begin{corollary}
Under the assumptions of the Davis--Kahan theorem,
\[
\operatorname{dist}_{\mathrm{op}}(U,U_*)
\leq
\frac{\sqrt{2}\norm{E}_{\mathrm{op}}}{\Delta},
\]
and
\[
\operatorname{dist}_F(U,U_*)
\leq
\frac{\sqrt{2r}\norm{E}_{\mathrm{op}}}{\Delta}.
\]
\end{corollary}

\begin{proof}
From the equivalence of subspace distances,
\[
\operatorname{dist}_{\mathrm{op}}(U,U_*)
\leq
\sqrt{2}\norm{\sin\Theta}_{\mathrm{op}},
\]
and
\[
\operatorname{dist}_F(U,U_*)
\leq
\sqrt{2}\norm{\sin\Theta}_F.
\]

Applying the Davis--Kahan theorem yields
\[
\operatorname{dist}_{\mathrm{op}}(U,U_*)
\leq
\frac{\sqrt{2}\norm{E}_{\mathrm{op}}}{\Delta},
\]
and
\[
\operatorname{dist}_F(U,U_*)
\leq
\frac{\sqrt{2r}\norm{E}_{\mathrm{op}}}{\Delta}.
\]
\end{proof}

\section{A Gap-Based Corollary}

Suppose
\[
|\lambda_1^*|
\geq
\cdots
\geq
|\lambda_n^*|,
\]
and
\[
|\lambda_1|
\geq
\cdots
\geq
|\lambda_n|.
\]

Define the eigengap
\[
\Delta
=
|\lambda_r^*|-|\lambda_{r+1}^*|.
\]

\begin{corollary}
Suppose
\[
\norm{E}_{\mathrm{op}}
<
\left(
1-\frac{1}{\sqrt{2}}
\right)\Delta.
\]

Then
\[
\norm{\sin\Theta}_{\mathrm{op}}
\leq
\frac{2\norm{E}_{\mathrm{op}}}{\Delta},
\]
and
\[
\norm{\sin\Theta}_F
\leq
\frac{2\sqrt{r}\norm{E}_{\mathrm{op}}}{\Delta}.
\]
\end{corollary}

\begin{proof}
By Weyl's inequality,
\[
|\lambda_i-\lambda_i^*|
\leq
\norm{E}_{\mathrm{op}}.
\]

Therefore the perturbed eigenvalues remain separated provided
\[
\norm{E}_{\mathrm{op}}
<
\left(
1-\frac{1}{\sqrt{2}}
\right)\Delta.
\]

This guarantees that the spectral gap condition in Davis--Kahan holds with effective gap at least
\[
\Delta/2.
\]

Applying the theorem yields
\[
\norm{\sin\Theta}_{\mathrm{op}}
\leq
\frac{\norm{E}_{\mathrm{op}}}{\Delta/2}
=
\frac{2\norm{E}_{\mathrm{op}}}{\Delta},
\]
and similarly,
\[
\norm{\sin\Theta}_F
\leq
\frac{2\sqrt{r}\norm{E}_{\mathrm{op}}}{\Delta}.
\]
\end{proof}

\section{Proof of the Davis--Kahan $\sin\Theta$ Theorem}

We now prove the Davis--Kahan theorem. The argument is based on controlling
\[
U_\perp^\top U_*,
\]
which, by the identities from the previous lecture, determines the subspace distance.

Recall the setup:
\[
M=M_*+E\in\mathbb{S}^n.
\]

Suppose
\[
M_*
=
[U_*\ \ U_{*,\perp}]
\begin{pmatrix}
\Lambda_* & 0\\
0 & \Lambda_{*,\perp}
\end{pmatrix}
\begin{pmatrix}
U_*^\top\\
U_{*,\perp}^\top
\end{pmatrix},
\]
and
\[
M
=
[U\ \ U_\perp]
\begin{pmatrix}
\Lambda & 0\\
0 & \Lambda_\perp
\end{pmatrix}
\begin{pmatrix}
U^\top\\
U_\perp^\top
\end{pmatrix}.
\]

Assume condition (1) of the Davis--Kahan theorem:
\[
\{\lambda_1^*,\dots,\lambda_r^*\}\subseteq[a,b]
\]
and
\[
\{\lambda_{r+1},\dots,\lambda_n\}
\subseteq
(-\infty,a-\Delta]
\cup
[b+\Delta,\infty).
\]

\subsection{Reduction to a Symmetric Interval}

Without loss of generality, assume
\[
a=-b\leq 0.
\]

Indeed, otherwise define shifted matrices
\[
\widetilde{M}_*
=
M_*-\frac{a+b}{2}I_n,
\qquad
\widetilde{M}
=
M-\frac{a+b}{2}I_n.
\]

Subtracting a scalar multiple of the identity shifts eigenvalues but does not change eigenvectors.

Thus we may assume the target eigenvalues lie in
\[
[-b,b].
\]

Consequently,
\[
\norm{\Lambda_*}_{\mathrm{op}}
\leq b,
\]
while the complementary eigenvalues satisfy
\[
\sigma_{\min}(\Lambda_\perp)
\geq
b+\Delta.
\]

\section{Main Decomposition}

We begin with the identity
\[
U_\perp^\top(M-M_*)U_*
=
\Lambda_\perp U_\perp^\top U_*
-
U_\perp^\top U_*\Lambda_*.
\]

\begin{proof}
Since
\[
MU_\perp
=
U_\perp\Lambda_\perp,
\]
we have
\[
U_\perp^\top M
=
\Lambda_\perp U_\perp^\top.
\]

Similarly,
\[
M_*U_*
=
U_*\Lambda_*.
\]

Therefore,
\[
U_\perp^\top(M-M_*)U_*
=
U_\perp^\top MU_*
-
U_\perp^\top M_*U_*.
\]

Substituting the eigenvalue relations gives
\[
U_\perp^\top(M-M_*)U_*
=
\Lambda_\perp U_\perp^\top U_*
-
U_\perp^\top U_*\Lambda_*.
\]
\end{proof}

Since
\[
E=M-M_*,
\]
the identity becomes
\[
U_\perp^\top E U_*
=
\Lambda_\perp U_\perp^\top U_*
-
U_\perp^\top U_*\Lambda_*.
\]

\section{Operator Norm Bound}

We now derive the operator norm estimate.

Using the reverse triangle inequality,
\[
\norm{U_\perp^\top E U_*}
\geq
\norm{\Lambda_\perp U_\perp^\top U_*}
-
\norm{U_\perp^\top U_*\Lambda_*}.
\]

Using the singular-value inequalities for rotationally invariant norms,
\[
\norm{\Lambda_\perp U_\perp^\top U_*}
\geq
\sigma_{\min}(\Lambda_\perp)
\norm{U_\perp^\top U_*},
\]
and
\[
\norm{U_\perp^\top U_*\Lambda_*}
\leq
\norm{U_\perp^\top U_*}
\norm{\Lambda_*}_{\mathrm{op}}.
\]

Hence,
\[
\norm{U_\perp^\top E U_*}
\geq
\left(
\sigma_{\min}(\Lambda_\perp)
-
\norm{\Lambda_*}_{\mathrm{op}}
\right)
\norm{U_\perp^\top U_*}.
\]

Using
\[
\sigma_{\min}(\Lambda_\perp)\geq b+\Delta
\]
and
\[
\norm{\Lambda_*}_{\mathrm{op}}\leq b,
\]
we obtain
\[
\norm{U_\perp^\top E U_*}
\geq
\Delta
\norm{U_\perp^\top U_*}.
\]

Therefore,
\[
\norm{U_\perp^\top U_*}
\leq
\frac{
\norm{U_\perp^\top E U_*}
}{\Delta}.
\]

Finally,
\[
\norm{U_\perp^\top E U_*}
\leq
\norm{U_\perp}_{\mathrm{op}}
\norm{E U_*}.
\]

Since
\[
\norm{U_\perp}_{\mathrm{op}}=1,
\]
we conclude
\[
\norm{U_\perp^\top U_*}
\leq
\frac{\norm{EU_*}}{\Delta}.
\]

Using
\[
\norm{EU_*}_{\mathrm{op}}
\leq
\norm{E}_{\mathrm{op}},
\]
we obtain
\[
\norm{U_\perp^\top U_*}_{\mathrm{op}}
\leq
\frac{\norm{E}_{\mathrm{op}}}{\Delta}.
\]

Recalling that
\[
\norm{\sin\Theta}_{\mathrm{op}}
=
\norm{U_\perp^\top U_*}_{\mathrm{op}},
\]
we arrive at
\[
\norm{\sin\Theta}_{\mathrm{op}}
\leq
\frac{\norm{E}_{\mathrm{op}}}{\Delta}.
\]

\section{Frobenius Norm Bound}

Using the Frobenius norm instead, the same argument gives
\[
\norm{U_\perp^\top U_*}_F
\leq
\frac{\norm{EU_*}_F}{\Delta}.
\]

Now,
\[
\norm{EU_*}_F
\leq
\norm{E}_{\mathrm{op}}
\norm{U_*}_F.
\]

Since
\[
U_*
\]
has orthonormal columns,
\[
\norm{U_*}_F=\sqrt{r}.
\]

Therefore,
\[
\norm{\sin\Theta}_F
=
\norm{U_\perp^\top U_*}_F
\leq
\frac{\sqrt{r}\norm{E}_{\mathrm{op}}}{\Delta}.
\]

This completes the proof of the Davis--Kahan theorem.

\section{Proof of the Gap Corollary}

We now prove the earlier corollary based on the eigengap.

Assume
\[
|\lambda_1^*|
\geq
\cdots
\geq
|\lambda_n^*|,
\]
and define
\[
\Delta
=
|\lambda_r^*|-|\lambda_{r+1}^*|.
\]

Suppose
\[
\norm{E}_{\mathrm{op}}
<
\left(
1-\frac{1}{\sqrt{2}}
\right)\Delta.
\]

We show that condition (2) of the Davis--Kahan theorem holds.

\subsection{Weyl's Inequality}

Let
\[
\lambda_i(M_*),
\qquad
\lambda_i(M)
\]
denote the eigenvalues sorted by value.

By Weyl's inequality,
\[
|\lambda_i(M)-\lambda_i(M_*)|
\leq
\norm{E}_{\mathrm{op}}
\]
for all
\[
i.
\]

Let
\[
r_1
=
\#\left\{
i:\lambda_i(M)\geq |\lambda_{r+1}^*|
\right\},
\]
and
\[
r_2
=
\#\left\{
i:\lambda_i(M)\leq -|\lambda_{r+1}^*|
\right\}.
\]

Then
\[
r=r_1+r_2.
\]

For indices
\[
i\leq r_1
\]
or
\[
i>n-r_2,
\]
Weyl's inequality gives
\[
|\lambda_i(M)|
\geq
|\lambda_i(M_*)|
-
\norm{E}_{\mathrm{op}}.
\]

Using the perturbation assumption,
\[
\norm{E}_{\mathrm{op}}
<
\left(
1-\frac{1}{\sqrt{2}}
\right)
(|\lambda_r^*|-|\lambda_{r+1}^*|),
\]
we obtain
\[
|\lambda_i(M)|
>
|\lambda_{r+1}^*|
+
\norm{E}_{\mathrm{op}}.
\]

On the other hand, for
\[
r_1<i\leq n-r_2,
\]
Weyl's inequality gives
\[
|\lambda_i(M)|
\leq
|\lambda_{r+1}^*|
+
\norm{E}_{\mathrm{op}}.
\]

Thus setting
\[
b
=
|\lambda_{r+1}^*|
+
\norm{E}_{\mathrm{op}}
\]
implies
\[
\{\lambda_{r+1},\dots,\lambda_n\}
\subseteq
[-b,b].
\]

Meanwhile,
\[
\{\lambda_1^*,\dots,\lambda_r^*\}
\subseteq
(-\infty,-|\lambda_r^*|]
\cup
[|\lambda_r^*|,\infty).
\]

Hence we may choose
\[
a=-b
\]
and
\[
\widetilde{\Delta}
=
|\lambda_r^*|
-
|\lambda_{r+1}^*|
-
\norm{E}_{\mathrm{op}}.
\]

Since
\[
\norm{E}_{\mathrm{op}}
<
\left(
1-\frac{1}{\sqrt{2}}
\right)\Delta,
\]
we have
\[
\widetilde{\Delta}
>
\frac{\Delta}{\sqrt{2}}.
\]

Applying the Davis--Kahan theorem gives
\[
\norm{\sin\Theta}_{\mathrm{op}}
\leq
\frac{\norm{E}_{\mathrm{op}}}{\widetilde{\Delta}}
\leq
\frac{\sqrt{2}\norm{E}_{\mathrm{op}}}{\Delta}.
\]

Using the earlier equivalence relations between subspace distances yields the stated bounds.

\section{Wedin's $\sin\Theta$ Theorem}

Davis--Kahan controls perturbations of eigenspaces for symmetric matrices. Wedin's theorem gives the analogous result for singular vectors.

Suppose
\[
M=M_*+E\in\R^{n\times m},
\]
and assume
\[
n\leq m.
\]

Write the singular value decompositions
\[
M_*
=
[U_*\ \ U_{*,\perp}]
\begin{pmatrix}
\Sigma_* & 0 & 0\\
0 & \Sigma_{*,\perp} & 0
\end{pmatrix}
\begin{pmatrix}
V_*^\top\\
V_{*,\perp}^\top
\end{pmatrix},
\]
and
\[
M
=
[U\ \ U_\perp]
\begin{pmatrix}
\Sigma & 0 & 0\\
0 & \Sigma_\perp & 0
\end{pmatrix}
\begin{pmatrix}
V^\top\\
V_\perp^\top
\end{pmatrix}.
\]

Here,
\[
\Sigma_*=\diag(\sigma_1^*,\dots,\sigma_r^*),
\]
and similarly for the remaining blocks.

\begin{theorem}[Wedin's $\sin\Theta$ Theorem]
Suppose
\[
\norm{E}_{\mathrm{op}}
\leq
\frac{\sigma_r^*-\sigma_{r+1}^*}{\sqrt{2}}.
\]

Then
\[
\operatorname{dist}_{\mathrm{op}}(U,U_*)
\vee
\operatorname{dist}_{\mathrm{op}}(V,V_*)
\leq
\frac{
2
\left(
\norm{E^\top U_*}_{\mathrm{op}}
\vee
\norm{EV_*}_{\mathrm{op}}
\right)
}{
\sigma_r^*-\sigma_{r+1}^*
},
\]
and
\[
\operatorname{dist}_F(U,U_*)
\vee
\operatorname{dist}_F(V,V_*)
\leq
\frac{
2
\left(
\norm{E^\top U_*}_F
\vee
\norm{EV_*}_F
\right)
}{
\sigma_r^*-\sigma_{r+1}^*
}.
\]
\end{theorem}

\begin{remark}
Wedin's theorem is the singular-vector analogue of Davis--Kahan. The proof follows the same general strategy:
\begin{enumerate}
\item derive a decomposition involving perturbation terms;
\item exploit the singular value gap;
\item bound subspace distances using operator or Frobenius norms.
\end{enumerate}
\end{remark}

\section{Community Detection in Networks}

We now apply perturbation theory to a statistical learning problem.

\subsection{The Stochastic Block Model}

Consider a graph
\[
G(n,p,q)
\]
with
\[
n
\]
vertices split into two communities of equal size
\[
n/2.
\]

Edges are generated independently according to:
\[
\Prob((v_i,v_j)\text{ is an edge})
=
\begin{cases}
p, & v_i,v_j \text{ in same community},\\
q, & \text{otherwise}.
\end{cases}
\]

The goal is to recover the hidden communities from one realization of the graph.

\section{Adjacency Matrix Representation}

Let
\[
A\in\R^{n\times n}
\]
be the adjacency matrix:
\[
A_{ij}
=
\begin{cases}
1, & (i,j)\text{ is an edge},\\
0, & \text{otherwise}.
\end{cases}
\]

Assume the first
\[
n/2
\]
vertices belong to the first community.

Then
\[
\E[A]
=
\begin{pmatrix}
p\mathbf{1}\mathbf{1}^\top & q\mathbf{1}\mathbf{1}^\top\\
q\mathbf{1}\mathbf{1}^\top & p\mathbf{1}\mathbf{1}^\top
\end{pmatrix},
\]
where each block has dimension
\[
n/2\times n/2.
\]

\section{Spectral Structure of the Population Matrix}

The matrix
\[
\E[A]
\]
has rank $2$.

Define
\[
u_1
=
\frac{1}{\sqrt{n}}
\begin{pmatrix}
1\\
\vdots\\
1
\end{pmatrix},
\]
and
\[
u_2
=
\frac{1}{\sqrt{n}}
\begin{pmatrix}
1\\
\vdots\\
1\\
-1\\
\vdots\\
-1
\end{pmatrix},
\]
where the first
\[
n/2
\]
entries are
\[
+1
\]
and the last
\[
n/2
\]
entries are
\[
-1.
\]

\begin{proposition}
The vectors
\[
u_1,u_2
\]
are eigenvectors of
\[
\E[A].
\]

Their eigenvalues are
\[
\lambda_1
=
\frac{n}{2}(p+q),
\]
and
\[
\lambda_2
=
\frac{n}{2}(p-q).
\]
\end{proposition}

\begin{proof}
For
\[
u_1,
\]
every row sum of
\[
\E[A]
\]
equals
\[
\frac{n}{2}p+\frac{n}{2}q
=
\frac{n}{2}(p+q).
\]

Hence
\[
\E[A]u_1
=
\frac{n}{2}(p+q)u_1.
\]

For
\[
u_2,
\]
rows inside the first community contribute
\[
\frac{n}{2}p-\frac{n}{2}q
=
\frac{n}{2}(p-q),
\]
while rows inside the second community contribute the negative of this quantity. Therefore
\[
\E[A]u_2
=
\frac{n}{2}(p-q)u_2.
\]
\end{proof}

\begin{remark}
The second eigenvector
\[
u_2
\]
encodes the community structure through its signs.

If the observed adjacency matrix
\[
A
\]
is close to
\[
\E[A]
\]
in operator norm, then Davis--Kahan implies that the empirical second eigenvector of
\[
A
\]
is close to
\[
u_2.
\]

Thus the communities can be approximately recovered by thresholding the signs of the empirical second eigenvector.
\end{remark}

\section{Spectral Clustering for Community Detection}

Recall the two-community stochastic block model
\[
G\sim G(n,p,q),
\]
where the vertex set is split into two equal communities of size $n/2$. Edges are sampled independently with probability
\[
\Prob((v_i,v_j)\text{ is an edge})
=
\begin{cases}
p, & v_i,v_j \text{ are in the same community},\\
q, & v_i,v_j \text{ are in different communities}.
\end{cases}
\]

Let
\[
A\in\R^{n\times n}
\]
be the adjacency matrix of $G$.

Assume, without loss of generality, that the first community is
\[
\{1,\dots,n/2\}.
\]

Then the population adjacency matrix has block form
\[
\E[A]
=
\begin{pmatrix}
p\mathbf{1}\mathbf{1}^\top & q\mathbf{1}\mathbf{1}^\top\\
q\mathbf{1}\mathbf{1}^\top & p\mathbf{1}\mathbf{1}^\top
\end{pmatrix},
\]
where each block is of size
\[
\frac{n}{2}\times \frac{n}{2}.
\]

Define
\[
u_1
=
\frac{1}{\sqrt{n}}
\begin{pmatrix}
1\\
\vdots\\
1
\end{pmatrix},
\qquad
u_2
=
\frac{1}{\sqrt{n}}
\begin{pmatrix}
1\\
\vdots\\
1\\
-1\\
\vdots\\
-1
\end{pmatrix},
\]
where $u_2$ has $n/2$ positive entries followed by $n/2$ negative entries.

Then
\[
\E[A]u_1
=
\frac{n}{2}(p+q)u_1
\]
and
\[
\E[A]u_2
=
\frac{n}{2}(p-q)u_2.
\]

Thus,
\[
\lambda_1(\E[A])
=
\frac{n}{2}(p+q),
\qquad
\lambda_2(\E[A])
=
\frac{n}{2}(p-q).
\]

The second eigenvector encodes the communities through its signs.

\subsection{Spectral Clustering Algorithm}

\begin{algorithm}[H]
\caption{Spectral Clustering for the Two-Community SBM}
\begin{algorithmic}[1]
\State \textbf{Input:} Graph $G$
\State Compute the adjacency matrix $A$
\State Compute $u_2(A)$, the eigenvector associated with the second-largest eigenvalue of $A$
\State Return the community labels
\[
\operatorname{sign}(u_2(A)).
\]
\end{algorithmic}
\end{algorithm}

\section{Misclassification Guarantee}

Define the signal strength
\[
\mu
=
q\wedge (p-q),
\]
and assume
\[
\mu>0.
\]

\begin{theorem}[Spectral Clustering Misclassification Bound]
Let
\[
G\sim G(n,p,q)
\]
with two equal-sized communities. Suppose
\[
q\wedge(p-q)=\mu>0.
\]

Then, with probability at least
\[
1-4e^{-n},
\]
the spectral clustering algorithm misclassifies at most
\[
\frac{C}{\mu^2}
\]
vertices, where $C>0$ is a universal constant.
\end{theorem}

\begin{proof}
Let
\[
u_2(\E[A])
\]
denote the unit-norm eigenvector of $\E[A]$ associated with
\[
\lambda_2(\E[A]).
\]

By Davis--Kahan,
\[
\sin\angle(u_2(\E[A]),u_2(A))
\leq
\frac{2\norm{A-\E[A]}_{\mathrm{op}}}
{
\min\{\lambda_2(\E[A]),\lambda_1(\E[A])-\lambda_2(\E[A])\}
}.
\]

Now,
\[
\lambda_2(\E[A])
=
\frac{n}{2}(p-q),
\]
and
\[
\lambda_1(\E[A])-\lambda_2(\E[A])
=
\frac{n}{2}(p+q)-\frac{n}{2}(p-q)
=
nq.
\]

Therefore,
\[
\min\{\lambda_2(\E[A]),\lambda_1(\E[A])-\lambda_2(\E[A])\}
\geq
c\mu n
\]
for a universal constant $c>0$.

Hence,
\[
\sin\angle(u_2(\E[A]),u_2(A))
\leq
C\frac{\norm{A-\E[A]}_{\mathrm{op}}}{\mu n}.
\]

We use the matrix concentration fact
\[
\Prob\left(
\norm{A-\E[A]}_{\mathrm{op}}
\geq
C\sqrt{n}
\right)
\leq
4e^{-n}.
\]

On this event,
\[
\sin\angle(u_2(\E[A]),u_2(A))
\leq
\frac{C}{\mu\sqrt{n}}.
\]

Since both vectors have unit norm,
\[
\min_{s\in\{-1,1\}}
\norm{u_2(\E[A])-s u_2(A)}_2
\leq
C\sin\angle(u_2(\E[A]),u_2(A)).
\]

Thus,
\[
\min_{s\in\{-1,1\}}
\norm{u_2(\E[A])-s u_2(A)}_2
\leq
\frac{C}{\mu\sqrt{n}}.
\]

Now define
\[
v_2
=
\sqrt{n}u_2(\E[A]).
\]

Then $v_2$ has entries
\[
+1
\]
on one community and
\[
-1
\]
on the other.

Choose
\[
s\in\{-1,1\}
\]
so that
\[
\norm{u_2(\E[A])-s u_2(A)}_2
\]
is minimized. Multiplying by $\sqrt{n}$ gives
\[
\norm{v_2-s\sqrt{n}u_2(A)}_2
\leq
\frac{C}{\mu}.
\]

A vertex is misclassified only if the sign of the empirical vector disagrees with the sign of $v_2$. For such a coordinate $i$,
\[
|v_2^{(i)}-s\sqrt{n}u_2^{(i)}(A)|
\geq
1.
\]

Therefore,
\[
\#\left\{
i:
\operatorname{sign}(s u_2^{(i)}(A))\neq v_2^{(i)}
\right\}
\leq
\sum_{i=1}^n
\left(
v_2^{(i)}-s\sqrt{n}u_2^{(i)}(A)
\right)^2.
\]

Hence,
\[
\#\{\text{misclassified vertices}\}
\leq
\norm{v_2-s\sqrt{n}u_2(A)}_2^2
\leq
\frac{C}{\mu^2}.
\]
\end{proof}

\begin{remark}
The proof has two main ingredients:
\begin{enumerate}
\item Davis--Kahan converts matrix perturbation into eigenvector perturbation.
\item A sign-thresholding argument converts eigenvector perturbation into a misclassification bound.
\end{enumerate}
\end{remark}

\section{Nets, Coverings, and Packings}

We now develop tools for proving operator norm bounds such as
\[
\norm{A-\E[A]}_{\mathrm{op}}
=
O(\sqrt{n})
\]
with high probability.

For a matrix
\[
A\in\R^{n\times m},
\]
the operator norm is
\[
\norm{A}_{\mathrm{op}}
=
\max_{v\in S^{m-1}}
\norm{Av}_2.
\]

This maximum ranges over infinitely many vectors. To use union bounds, we replace the sphere by a finite approximation.

\section{Metric Spaces}

\begin{definition}[Metric Space]
A metric space is a pair
\[
(T,d),
\]
where
\[
d:T\times T\to\R_+
\]
satisfies, for all $x,y,z\in T$:
\begin{enumerate}
\item
\[
d(x,x)=0;
\]
\item if $x\neq y$, then
\[
d(x,y)>0;
\]
\item
\[
d(x,y)=d(y,x);
\]
\item
\[
d(x,y)\leq d(x,z)+d(z,y).
\]
\end{enumerate}
\end{definition}

\begin{example}
The Euclidean space
\[
\R^d
\]
with
\[
d(x,y)=\norm{x-y}_2
\]
is a metric space.
\end{example}

\section{\texorpdfstring{$\varepsilon$}{epsilon}-Nets}

\begin{definition}[$\varepsilon$-Net]
Let
\[
(T,d)
\]
be a metric space, let
\[
K\subseteq T,
\]
and let
\[
\varepsilon>0.
\]

A subset
\[
\mathcal{N}\subseteq K
\]
is called an $\varepsilon$-net of $K$ if, for every
\[
x\in K,
\]
there exists
\[
x_0\in\mathcal{N}
\]
such that
\[
d(x,x_0)\leq \varepsilon.
\]
\end{definition}

\begin{definition}[Covering Number]
The covering number of $K$ at scale $\varepsilon$ is
\[
N(K,d,\varepsilon),
\]
defined as the cardinality of the smallest $\varepsilon$-net of $K$.
\end{definition}

\begin{remark}
In Euclidean space, an $\varepsilon$-net covers a set using Euclidean balls of radius $\varepsilon$ centered at the points of the net.
\end{remark}

\section{Packing Numbers}

\begin{definition}[$\varepsilon$-Separated Set]
Let
\[
(T,d)
\]
be a metric space. A subset
\[
\mathcal{P}\subseteq T
\]
is called $\varepsilon$-separated if
\[
d(x,y)>\varepsilon
\]
for all distinct
\[
x,y\in\mathcal{P}.
\]
\end{definition}

\begin{definition}[Packing Number]
The packing number of $K$ at scale $\varepsilon$ is
\[
P(K,d,\varepsilon),
\]
defined as the cardinality of the largest $\varepsilon$-separated subset of $K$.
\end{definition}

\section{Relationship Between Covering and Packing Numbers}

Covering and packing numbers are closely related.

\begin{lemma}
For any subset
\[
K\subseteq T
\]
and any
\[
\varepsilon>0,
\]
we have
\[
P(K,d,2\varepsilon)
\leq
N(K,d,\varepsilon)
\leq
P(K,d,\varepsilon).
\]
\end{lemma}

\begin{proof}
We first prove the lower bound.

Let
\[
\mathcal{P}
\]
be a $2\varepsilon$-separated subset of $K$, and let
\[
\mathcal{N}
\]
be an $\varepsilon$-net of $K$.

For each
\[
x\in\mathcal{P},
\]
there exists
\[
y\in\mathcal{N}
\]
such that
\[
d(x,y)\leq \varepsilon.
\]

We claim that distinct points in $\mathcal{P}$ must map to distinct points in $\mathcal{N}$. Suppose two distinct points
\[
x,z\in\mathcal{P}
\]
were both within distance $\varepsilon$ of the same $y\in\mathcal{N}$. Then by the triangle inequality,
\[
d(x,z)
\leq
d(x,y)+d(y,z)
\leq
2\varepsilon,
\]
contradicting the fact that $\mathcal{P}$ is $2\varepsilon$-separated.

Thus,
\[
|\mathcal{P}|\leq |\mathcal{N}|.
\]

Since $\mathcal{P}$ and $\mathcal{N}$ were arbitrary,
\[
P(K,d,2\varepsilon)
\leq
N(K,d,\varepsilon).
\]

Now we prove the upper bound.

Let
\[
\mathcal{N}
\]
be a maximal $\varepsilon$-separated subset of $K$. Then
\[
|\mathcal{N}|=P(K,d,\varepsilon).
\]

We claim that $\mathcal{N}$ is an $\varepsilon$-net of $K$.

Suppose not. Then there exists
\[
x\in K
\]
such that
\[
d(x,y)>\varepsilon
\]
for every
\[
y\in\mathcal{N}.
\]

Then
\[
\mathcal{N}\cup\{x\}
\]
is a larger $\varepsilon$-separated subset of $K$, contradicting maximality of $\mathcal{N}$.

Therefore $\mathcal{N}$ is an $\varepsilon$-net, and so
\[
N(K,d,\varepsilon)
\leq
|\mathcal{N}|
=
P(K,d,\varepsilon).
\]
\end{proof}

\begin{remark}
This lemma allows us to bound covering numbers using packing arguments, which are often easier because separated balls are disjoint.
\end{remark}

\section{Covering Numbers via Volume}

In many applications, we need quantitative bounds on covering numbers. We now derive such bounds in Euclidean space using volume.

Throughout this section, let
\[
B_2^n
=
\{x\in\R^n:\norm{x}_2\leq 1\}.
\]

\begin{definition}[Minkowski Sum]
Let
\[
A,B\subseteq \R^n.
\]
The Minkowski sum of $A$ and $B$ is
\[
A+B
=
\{a+b:a\in A,\ b\in B\}.
\]
\end{definition}

\begin{proposition}[Volume Bounds for Covering Numbers]
Let
\[
K\subseteq \R^n
\]
and let
\[
\varepsilon>0.
\]
Then
\[
\frac{\operatorname{Vol}(K)}
{\operatorname{Vol}(\varepsilon B_2^n)}
\leq
N(K,\norm{\cdot}_2,\varepsilon)
\leq
P(K,\norm{\cdot}_2,\varepsilon)
\leq
\frac{
\operatorname{Vol}\left(K+\frac{\varepsilon}{2}B_2^n\right)
}{
\operatorname{Vol}\left(\frac{\varepsilon}{2}B_2^n\right)
}.
\]
\end{proposition}

\begin{proof}
For the lower bound, let
\[
N=N(K,\norm{\cdot}_2,\varepsilon).
\]
Then $K$ can be covered by $N$ Euclidean balls of radius $\varepsilon$. Therefore,
\[
\operatorname{Vol}(K)
\leq
N\operatorname{Vol}(\varepsilon B_2^n).
\]
Rearranging gives
\[
\frac{\operatorname{Vol}(K)}
{\operatorname{Vol}(\varepsilon B_2^n)}
\leq
N(K,\norm{\cdot}_2,\varepsilon).
\]

For the upper bound, let
\[
\mathcal{P}
\]
be a maximal $\varepsilon$-separated subset of $K$, so that
\[
|\mathcal{P}|=P(K,\norm{\cdot}_2,\varepsilon).
\]

The balls
\[
x+\frac{\varepsilon}{2}B_2^n,
\qquad x\in\mathcal{P},
\]
are pairwise disjoint. Moreover, since $x\in K$,
\[
x+\frac{\varepsilon}{2}B_2^n
\subseteq
K+\frac{\varepsilon}{2}B_2^n.
\]

Hence,
\[
|\mathcal{P}|
\operatorname{Vol}\left(
\frac{\varepsilon}{2}B_2^n
\right)
\leq
\operatorname{Vol}\left(
K+\frac{\varepsilon}{2}B_2^n
\right).
\]

Thus,
\[
P(K,\norm{\cdot}_2,\varepsilon)
\leq
\frac{
\operatorname{Vol}\left(K+\frac{\varepsilon}{2}B_2^n\right)
}{
\operatorname{Vol}\left(\frac{\varepsilon}{2}B_2^n\right)
}.
\]

Finally, from the relationship between covering and packing numbers,
\[
N(K,\norm{\cdot}_2,\varepsilon)
\leq
P(K,\norm{\cdot}_2,\varepsilon).
\]
\end{proof}

\section{Covering the Euclidean Unit Ball and Sphere}

\begin{corollary}
For the Euclidean unit ball
\[
B_2^n,
\]
the covering number satisfies
\[
\left(\frac{1}{\varepsilon}\right)^n
\leq
N(B_2^n,\norm{\cdot}_2,\varepsilon)
\leq
\left(1+\frac{2}{\varepsilon}\right)^n.
\]

In particular, for
\[
\varepsilon\in(0,1),
\]
\[
N(B_2^n,\norm{\cdot}_2,\varepsilon)
\leq
\left(\frac{3}{\varepsilon}\right)^n.
\]

The same upper bound holds for the sphere
\[
S^{n-1}.
\]
\end{corollary}

\begin{proof}
Recall that volume scales as
\[
\operatorname{Vol}(\varepsilon B_2^n)
=
\varepsilon^n\operatorname{Vol}(B_2^n).
\]

Applying the lower bound from the previous proposition with
\[
K=B_2^n
\]
gives
\[
N(B_2^n,\norm{\cdot}_2,\varepsilon)
\geq
\frac{\operatorname{Vol}(B_2^n)}
{\operatorname{Vol}(\varepsilon B_2^n)}
=
\frac{1}{\varepsilon^n}
=
\left(\frac{1}{\varepsilon}\right)^n.
\]

For the upper bound, observe that
\[
B_2^n+\frac{\varepsilon}{2}B_2^n
=
\left(1+\frac{\varepsilon}{2}\right)B_2^n.
\]

Thus,
\[
N(B_2^n,\norm{\cdot}_2,\varepsilon)
\leq
P(B_2^n,\norm{\cdot}_2,\varepsilon)
\leq
\frac{
\operatorname{Vol}\left(
\left(1+\frac{\varepsilon}{2}\right)B_2^n
\right)
}{
\operatorname{Vol}\left(
\frac{\varepsilon}{2}B_2^n
\right)
}.
\]

Using volume scaling,
\[
\frac{
\operatorname{Vol}\left(
\left(1+\frac{\varepsilon}{2}\right)B_2^n
\right)
}{
\operatorname{Vol}\left(
\frac{\varepsilon}{2}B_2^n
\right)
}
=
\left(
\frac{1+\varepsilon/2}{\varepsilon/2}
\right)^n
=
\left(
1+\frac{2}{\varepsilon}
\right)^n.
\]

If
\[
\varepsilon\in(0,1),
\]
then
\[
1+\frac{2}{\varepsilon}
\leq
\frac{3}{\varepsilon}.
\]

Therefore,
\[
N(B_2^n,\norm{\cdot}_2,\varepsilon)
\leq
\left(\frac{3}{\varepsilon}\right)^n.
\]

Since
\[
S^{n-1}\subseteq B_2^n,
\]
any $\varepsilon$-net of $B_2^n$ also controls $S^{n-1}$ up to the same order, giving the same upper bound.
\end{proof}

\begin{definition}[Metric Entropy]
The metric entropy of a set $K$ at scale $\varepsilon$ is
\[
\log N(K,d,\varepsilon).
\]
\end{definition}

\begin{remark}
Covering numbers measure the geometric complexity of a set. In many probabilistic arguments, the quantity
\[
\log N(K,d,\varepsilon)
\]
is more important than the covering number itself.
\end{remark}

\section{Error-Correcting Codes}

Covering and packing ideas also arise in coding theory.

Suppose one wants to send a message
\[
x
\]
through a noisy channel. An adversary may corrupt the message by changing some letters. The goal is to encode the original message with redundancy so that the receiver can recover the original message even after some corruptions.

\subsection{Repeating Codes}

A simple strategy is to repeat the message many times. If the message is repeated
\[
2r+1
\]
times, then majority decoding can recover the original message even if up to
\[
r
\]
positions are corrupted.

However, this is inefficient because a $k$-symbol message becomes a message of length at least
\[
(2r+1)k.
\]

We will show that much better scaling is possible.

\section{Binary Error-Correcting Codes}

For simplicity, consider binary strings.

\begin{definition}[Error-Correcting Code]
An error-correcting code that encodes $k$-bit strings into $n$-bit strings and corrects up to $r$ errors consists of:
\begin{itemize}
\item an encoding map
\[
E:\{0,1\}^k\to\{0,1\}^n,
\]
\item a decoding map
\[
D:\{0,1\}^n\to\{0,1\}^k,
\]
\end{itemize}
such that
\[
D(y)=x
\]
for every
\[
x\in\{0,1\}^k
\]
and every
\[
y\in\{0,1\}^n
\]
that differs from
\[
E(x)
\]
in at most $r$ coordinates.
\end{definition}

\section{The Hamming Cube}

\begin{definition}[Hamming Distance]
The Hamming distance between
\[
x,y\in\{0,1\}^n
\]
is
\[
d_H(x,y)
=
\#\{i:x_i\neq y_i\}.
\]
\end{definition}

\begin{lemma}
The set
\[
\{0,1\}^n
\]
equipped with the Hamming distance
\[
d_H
\]
is a metric space.
\end{lemma}

\begin{proof}
Clearly,
\[
d_H(x,x)=0.
\]

If
\[
x\neq y,
\]
then $x$ and $y$ differ in at least one coordinate, so
\[
d_H(x,y)>0.
\]

Symmetry follows because
\[
x_i\neq y_i
\]
if and only if
\[
y_i\neq x_i.
\]

For the triangle inequality, fix
\[
x,y,z\in\{0,1\}^n.
\]
If
\[
x_i\neq y_i,
\]
then at least one of the two inequalities
\[
x_i\neq z_i
\qquad\text{or}\qquad
z_i\neq y_i
\]
must hold. Hence,
\[
\ind_{\{x_i\neq y_i\}}
\leq
\ind_{\{x_i\neq z_i\}}
+
\ind_{\{z_i\neq y_i\}}.
\]

Summing over $i$ gives
\[
d_H(x,y)
\leq
d_H(x,z)+d_H(z,y).
\]
\end{proof}

\section{Covering and Packing in the Hamming Cube}

Let
\[
K=\{0,1\}^n.
\]

For
\[
m\in\{0,1,\dots,n\},
\]
the size of a Hamming ball of radius $m$ is
\[
\sum_{i=0}^m \binom{n}{i}.
\]

\begin{proposition}
For
\[
K=\{0,1\}^n
\]
with Hamming distance $d_H$ and
\[
m\in\{0,\dots,n\},
\]
we have
\[
\frac{2^n}{\sum_{i=0}^m \binom{n}{i}}
\leq
N(K,d_H,m)
\leq
P(K,d_H,m)
\leq
\frac{2^n}{\sum_{i=0}^{\lfloor (m-1)/2\rfloor}\binom{n}{i}}.
\]
\end{proposition}

\begin{proof}
The lower bound follows from the same volume argument as in Euclidean space, replacing volume by cardinality.

An $m$-net covers
\[
\{0,1\}^n
\]
by Hamming balls of radius $m$. Each such ball contains
\[
\sum_{i=0}^m\binom{n}{i}
\]
points. Therefore,
\[
2^n
\leq
N(K,d_H,m)
\sum_{i=0}^m\binom{n}{i}.
\]

Rearranging gives
\[
N(K,d_H,m)
\geq
\frac{2^n}{\sum_{i=0}^m\binom{n}{i}}.
\]

The middle inequality
\[
N(K,d_H,m)\leq P(K,d_H,m)
\]
is the general covering-packing relationship.

For the packing upper bound, let
\[
\mathcal{P}
\]
be an $m$-separated set. Around each point of $\mathcal{P}$, place a Hamming ball of radius
\[
\left\lfloor\frac{m-1}{2}\right\rfloor.
\]

These balls are disjoint. Hence,
\[
|\mathcal{P}|
\sum_{i=0}^{\lfloor (m-1)/2\rfloor}
\binom{n}{i}
\leq
2^n.
\]

Therefore,
\[
P(K,d_H,m)
\leq
\frac{2^n}{\sum_{i=0}^{\lfloor (m-1)/2\rfloor}\binom{n}{i}}.
\]
\end{proof}

\section{Codes from Packings}

The key observation is that a large separated set in the Hamming cube gives an error-correcting code.

\begin{lemma}
Suppose
\[
k,n,r\in\N
\]
satisfy
\[
\log_2 P(\{0,1\}^n,d_H,2r)>k.
\]

Then there exists an error-correcting code that encodes $k$-bit strings into $n$-bit strings and corrects up to $r$ errors.
\end{lemma}

\begin{proof}
The assumption implies that there exists a $2r$-separated set
\[
\mathcal{N}\subseteq\{0,1\}^n
\]
with
\[
|\mathcal{N}|\geq 2^k.
\]

Choose an injective encoding map
\[
E:\{0,1\}^k\to\mathcal{N}.
\]

Define the decoder by nearest neighbor decoding:
\[
D(y)
=
\argmin_{x\in\{0,1\}^k}
d_H(E(x),y).
\]

Now suppose
\[
y
\]
is obtained from
\[
E(x)
\]
by corrupting at most $r$ coordinates. Then
\[
d_H(y,E(x))\leq r.
\]

For any other message
\[
x'\neq x,
\]
the separation property gives
\[
d_H(E(x),E(x'))>2r.
\]

By the triangle inequality,
\[
d_H(y,E(x'))
\geq
d_H(E(x),E(x'))-d_H(y,E(x))
>
2r-r
=
r.
\]

Thus $E(x)$ is strictly closer to $y$ than any other codeword. Therefore,
\[
D(y)=x.
\]
\end{proof}

\section{Existence of Efficient Error-Correcting Codes}

\begin{theorem}
Suppose
\[
k,n,r\in\N
\]
satisfy
\[
n
\geq
k+2r\log_2\left(\frac{en}{2r}\right).
\]

Then there exists an error-correcting code that encodes $k$-bit strings into $n$-bit strings and corrects up to $r$ errors.
\end{theorem}

\begin{proof}
By the previous lemma, it suffices to show that
\[
P(\{0,1\}^n,d_H,2r)\geq 2^k.
\]

Using the covering-packing inequality,
\[
P(\{0,1\}^n,d_H,2r)
\geq
N(\{0,1\}^n,d_H,2r).
\]

From the Hamming cube covering bound,
\[
N(\{0,1\}^n,d_H,2r)
\geq
\frac{2^n}
{\sum_{i=0}^{2r}\binom{n}{i}}.
\]

We use the standard binomial sum bound
\[
\sum_{i=0}^{2r}\binom{n}{i}
\leq
\left(\frac{en}{2r}\right)^{2r}.
\]

Therefore,
\[
P(\{0,1\}^n,d_H,2r)
\geq
2^n
\left(\frac{2r}{en}\right)^{2r}.
\]

Taking base-$2$ logarithms,
\[
\log_2 P(\{0,1\}^n,d_H,2r)
\geq
n-2r\log_2\left(\frac{en}{2r}\right).
\]

By assumption,
\[
n-2r\log_2\left(\frac{en}{2r}\right)
\geq
k.
\]

Hence,
\[
P(\{0,1\}^n,d_H,2r)
\geq
2^k.
\]

Thus there exists an error-correcting code encoding $k$ bits into $n$ bits and correcting up to $r$ errors.
\end{proof}

\begin{remark}
Ignoring logarithmic factors, the required code length grows roughly linearly in
\[
k+r,
\]
which is far more efficient than the naive repeating code length
\[
(2r+1)k.
\]
\end{remark}

\section{Singular Values of Sub-Gaussian Matrices}

We now use nets to obtain high-probability bounds on the singular values of random matrices.

The guiding idea is that if
\[
\max_{x\in S^{m-1}}\norm{Ax}_2
\approx
\max_{x\in\mathcal{N}}\norm{Ax}_2
\]
for a finite net
\[
\mathcal{N}\subseteq S^{m-1},
\]
then the operator norm can be controlled by a union bound over the finite set $\mathcal{N}$.

\section{Approximating the Operator Norm by a Net}

\begin{lemma}
Let
\[
A\in\R^{n\times m}
\]
and let
\[
\varepsilon\in[0,1).
\]

If
\[
\mathcal{N}
\]
is an $\varepsilon$-net of
\[
S^{m-1},
\]
then
\[
\sup_{x\in\mathcal{N}}\norm{Ax}_2
\leq
\norm{A}_{\mathrm{op}}
\leq
\frac{1}{1-\varepsilon}
\sup_{x\in\mathcal{N}}\norm{Ax}_2.
\]
\end{lemma}

\begin{proof}
The lower bound is immediate because
\[
\mathcal{N}\subseteq S^{m-1}.
\]

For the upper bound, choose
\[
x\in S^{m-1}
\]
such that
\[
\norm{A}_{\mathrm{op}}=\norm{Ax}_2.
\]

Since $\mathcal{N}$ is an $\varepsilon$-net, there exists
\[
y\in\mathcal{N}
\]
such that
\[
\norm{x-y}_2\leq \varepsilon.
\]

Then
\[
\norm{Ax}_2
\leq
\norm{A(x-y)}_2+\norm{Ay}_2.
\]

Using the definition of the operator norm,
\[
\norm{A(x-y)}_2
\leq
\norm{A}_{\mathrm{op}}\norm{x-y}_2
\leq
\varepsilon\norm{A}_{\mathrm{op}}.
\]

Thus,
\[
\norm{A}_{\mathrm{op}}
\leq
\varepsilon\norm{A}_{\mathrm{op}}
+
\norm{Ay}_2.
\]

Rearranging,
\[
(1-\varepsilon)\norm{A}_{\mathrm{op}}
\leq
\norm{Ay}_2
\leq
\sup_{z\in\mathcal{N}}\norm{Az}_2.
\]

Therefore,
\[
\norm{A}_{\mathrm{op}}
\leq
\frac{1}{1-\varepsilon}
\sup_{z\in\mathcal{N}}\norm{Az}_2.
\]
\end{proof}

\section{Net Approximation for Symmetric Quadratic Forms}

For symmetric matrices, the operator norm can be written in terms of quadratic forms.

\begin{lemma}
Let
\[
A\in\mathbb{S}^n
\]
and let
\[
\varepsilon\in[0,1/2).
\]

If
\[
\mathcal{N}
\]
is an $\varepsilon$-net of
\[
S^{n-1},
\]
then
\[
\sup_{x\in\mathcal{N}}|\inner{x}{Ax}|
\leq
\norm{A}_{\mathrm{op}}
\leq
\frac{1}{1-2\varepsilon}
\sup_{x\in\mathcal{N}}|\inner{x}{Ax}|.
\]
\end{lemma}

\begin{proof}
The lower bound follows since
\[
|\inner{x}{Ax}|
\leq
\norm{A}_{\mathrm{op}}
\]
for every unit vector $x$.

For the upper bound, choose
\[
x\in S^{n-1}
\]
such that
\[
|\inner{x}{Ax}|
=
\norm{A}_{\mathrm{op}}.
\]

Let
\[
y\in\mathcal{N}
\]
satisfy
\[
\norm{x-y}_2\leq \varepsilon.
\]

Then
\[
\inner{x}{Ax}
=
\inner{y}{Ay}
+
\inner{x-y}{Ax}
+
\inner{y}{A(x-y)}.
\]

Since $A$ is symmetric,
\[
|\inner{x-y}{Ax}|
\leq
\norm{x-y}_2\norm{A}_{\mathrm{op}}\norm{x}_2
\leq
\varepsilon\norm{A}_{\mathrm{op}},
\]
and similarly,
\[
|\inner{y}{A(x-y)}|
\leq
\varepsilon\norm{A}_{\mathrm{op}}.
\]

Therefore,
\[
\norm{A}_{\mathrm{op}}
\leq
|\inner{y}{Ay}|
+
2\varepsilon\norm{A}_{\mathrm{op}}.
\]

Rearranging gives
\[
(1-2\varepsilon)\norm{A}_{\mathrm{op}}
\leq
|\inner{y}{Ay}|
\leq
\sup_{z\in\mathcal{N}}|\inner{z}{Az}|.
\]

Thus,
\[
\norm{A}_{\mathrm{op}}
\leq
\frac{1}{1-2\varepsilon}
\sup_{z\in\mathcal{N}}|\inner{z}{Az}|.
\]
\end{proof}

\section{Sub-Gaussian Random Vectors}

\begin{definition}[Sub-Gaussian Random Vector]
A random vector
\[
Z\in\R^m
\]
is called sub-Gaussian if
\[
\norm{Z}_{\psi_2}
:=
\sup_{v\in S^{m-1}}
\norm{\inner{Z}{v}}_{\psi_2}
<\infty.
\]
\end{definition}

\begin{definition}[Isotropic Random Vector]
A random vector
\[
Z\in\R^m
\]
is called isotropic if
\[
\E[ZZ^\top]=I_m.
\]
Equivalently,
\[
\E[\inner{Z}{v}^2]=\norm{v}_2^2
\]
for every
\[
v\in\R^m.
\]
\end{definition}

\begin{lemma}
If
\[
Z\sim \Normal(0,I_m),
\]
then $Z$ is sub-Gaussian and isotropic.
\end{lemma}

\begin{proof}
For any
\[
v\in S^{m-1},
\]
we have
\[
\inner{Z}{v}\sim \Normal(0,1).
\]

Thus
\[
\norm{\inner{Z}{v}}_{\psi_2}
\]
is bounded by a universal constant uniformly over $v$, so $Z$ is sub-Gaussian.

Also,
\[
\E[ZZ^\top]
=
I_m
\]
by the covariance structure of the standard Gaussian. Hence $Z$ is isotropic.
\end{proof}

\section{Singular Value Bounds for Sub-Gaussian Matrices}

\begin{theorem}
Let
\[
A\in\R^{n\times m},
\qquad m\leq n,
\]
be a random matrix with independent, mean-zero, isotropic, sub-Gaussian rows
\[
A_1,\dots,A_n\in\R^m.
\]

Let
\[
K=\max_{1\leq i\leq n}\norm{A_i}_{\psi_2}.
\]

Then for every
\[
t\geq0,
\]
with probability at least
\[
1-2e^{-t^2},
\]
we have
\[
\sqrt{n}-CK^2(\sqrt{m}+t)
\leq
\sigma_m(A)
\leq
\sigma_1(A)
\leq
\sqrt{n}+CK^2(\sqrt{m}+t),
\]
where $C>0$ is a universal constant.
\end{theorem}

\section{An Equivalent Isometry Formulation}

The singular value bound is equivalent to saying that
\[
\frac{1}{n}A^\top A
\]
is close to the identity.

\begin{lemma}[Isometry Equivalence]
Let
\[
A\in\R^{n\times m},
\qquad m\leq n,
\]
and let
\[
\varepsilon>0.
\]

The following are equivalent:
\begin{enumerate}
\item
\[
\norm{A^\top A-I_m}_{\mathrm{op}}\leq \varepsilon.
\]

\item For every
\[
x\in\R^m,
\]
\[
(1-\varepsilon)\norm{x}_2^2
\leq
\norm{Ax}_2^2
\leq
(1+\varepsilon)\norm{x}_2^2.
\]

\item
\[
\sqrt{1-\varepsilon}
\leq
\sigma_m(A)
\leq
\sigma_1(A)
\leq
\sqrt{1+\varepsilon}.
\]
\end{enumerate}
\end{lemma}

\begin{proof}
It suffices to check the statement on unit vectors.

For
\[
\norm{x}_2=1,
\]
we have
\[
x^\top(A^\top A-I_m)x
=
\norm{Ax}_2^2-1.
\]

Thus,
\[
\norm{A^\top A-I_m}_{\mathrm{op}}
=
\sup_{x\in S^{m-1}}
\left|
\norm{Ax}_2^2-1
\right|.
\]

Therefore, the operator norm bound is equivalent to
\[
1-\varepsilon
\leq
\norm{Ax}_2^2
\leq
1+\varepsilon
\]
for all
\[
x\in S^{m-1}.
\]

By homogeneity, this is equivalent to the same statement for all $x\in\R^m$:
\[
(1-\varepsilon)\norm{x}_2^2
\leq
\norm{Ax}_2^2
\leq
(1+\varepsilon)\norm{x}_2^2.
\]

Finally, by the variational characterization of singular values,
\[
\sigma_1(A)
=
\max_{x\in S^{m-1}}\norm{Ax}_2,
\qquad
\sigma_m(A)
=
\min_{x\in S^{m-1}}\norm{Ax}_2.
\]

Thus the previous display is equivalent to
\[
\sqrt{1-\varepsilon}
\leq
\sigma_m(A)
\leq
\sigma_1(A)
\leq
\sqrt{1+\varepsilon}.
\]
\end{proof}

\section{Proof of the Singular Value Bound}

We prove the theorem by showing that
\[
\frac{1}{n}A^\top A
\]
is close to the identity in operator norm.

It is enough to prove that
\[
\norm{\frac{1}{n}A^\top A-I_m}_{\mathrm{op}}
\leq
K^2(\delta\vee \delta^2),
\]
where
\[
\delta
=
C\left(
\sqrt{\frac{m}{n}}+\frac{t}{\sqrt{n}}
\right).
\]

\subsection{Step 1: Approximation by a Net}

Let
\[
\mathcal{N}
\]
be a $1/4$-net of
\[
S^{m-1}.
\]

From the covering number bound,
\[
|\mathcal{N}|
\leq
9^m.
\]

Apply the symmetric quadratic form net lemma to
\[
\frac{1}{n}A^\top A-I_m.
\]

Since
\[
\frac{1}{1-2(1/4)}=2,
\]
we obtain
\[
\norm{\frac{1}{n}A^\top A-I_m}_{\mathrm{op}}
\leq
2
\max_{x\in\mathcal{N}}
\left|
\inner{x}{
\left(
\frac{1}{n}A^\top A-I_m
\right)x}
\right|.
\]

But
\[
\inner{x}{
\left(
\frac{1}{n}A^\top A-I_m
\right)x}
=
\frac{1}{n}\norm{Ax}_2^2-1.
\]

Thus,
\[
\norm{\frac{1}{n}A^\top A-I_m}_{\mathrm{op}}
\leq
2
\max_{x\in\mathcal{N}}
\left|
\frac{1}{n}\norm{Ax}_2^2-1
\right|.
\]

Therefore, it suffices to show
\[
\max_{x\in\mathcal{N}}
\left|
\frac{1}{n}\norm{Ax}_2^2-1
\right|
\leq
\frac{\varepsilon}{2},
\]
where
\[
\varepsilon=K^2(\delta\vee\delta^2).
\]

\subsection{Step 2: Concentration for a Fixed Net Point}

Fix
\[
x\in\mathcal{N}.
\]

Write
\[
\norm{Ax}_2^2
=
\sum_{i=1}^n \inner{A_i}{x}^2.
\]

Define
\[
X_i=\inner{A_i}{x}.
\]

Since the rows $A_i$ are independent, the variables $X_i$ are independent. Since each row is isotropic,
\[
\E[X_i^2]=1.
\]

Since each row is sub-Gaussian,
\[
\norm{X_i}_{\psi_2}
=
\norm{\inner{A_i}{x}}_{\psi_2}
\leq
K.
\]

Therefore,
\[
X_i^2-1
\]
are independent, mean-zero, sub-exponential random variables satisfying
\[
\norm{X_i^2-1}_{\psi_1}
\leq
CK^2.
\]

By Bernstein's inequality,
\[
\Prob\left(
\left|
\frac{1}{n}\sum_{i=1}^n (X_i^2-1)
\right|
\geq
\frac{\varepsilon}{2}
\right)
\leq
2\exp\left(
-c_1n
\min\left\{
\frac{\varepsilon^2}{K^4},
\frac{\varepsilon}{K^2}
\right\}
\right).
\]

Because
\[
\varepsilon
=
K^2(\delta\vee\delta^2),
\]
we have
\[
\min\left\{
\frac{\varepsilon^2}{K^4},
\frac{\varepsilon}{K^2}
\right\}
\geq
c\delta^2.
\]

Hence,
\[
\Prob\left(
\left|
\frac{1}{n}\norm{Ax}_2^2-1
\right|
\geq
\frac{\varepsilon}{2}
\right)
\leq
2\exp(-c_1\delta^2 n).
\]

Using
\[
\delta
=
C\left(
\sqrt{\frac{m}{n}}+\frac{t}{\sqrt{n}}
\right),
\]
we have
\[
\delta^2 n
\geq
C^2(m+t^2).
\]

Thus,
\[
\Prob\left(
\left|
\frac{1}{n}\norm{Ax}_2^2-1
\right|
\geq
\frac{\varepsilon}{2}
\right)
\leq
2\exp(-c_1C^2(m+t^2)).
\]

\subsection{Step 3: Union Bound Over the Net}

Since
\[
|\mathcal{N}|\leq 9^m,
\]
the union bound gives
\[
\Prob\left(
\max_{x\in\mathcal{N}}
\left|
\frac{1}{n}\norm{Ax}_2^2-1
\right|
\geq
\frac{\varepsilon}{2}
\right)
\leq
9^m\cdot
2\exp(-c_1C^2(m+t^2)).
\]

Equivalently,
\[
\Prob\left(
\max_{x\in\mathcal{N}}
\left|
\frac{1}{n}\norm{Ax}_2^2-1
\right|
\geq
\frac{\varepsilon}{2}
\right)
\leq
2\exp\left(
m\log 9-c_1C^2(m+t^2)
\right).
\]

Choose the universal constant $C$ large enough so that
\[
c_1C^2\geq \log 9+1.
\]

Then
\[
m\log 9-c_1C^2(m+t^2)
\leq
-t^2.
\]

Therefore,
\[
\Prob\left(
\max_{x\in\mathcal{N}}
\left|
\frac{1}{n}\norm{Ax}_2^2-1
\right|
\geq
\frac{\varepsilon}{2}
\right)
\leq
2e^{-t^2}.
\]

On the complementary event,
\[
\norm{\frac{1}{n}A^\top A-I_m}_{\mathrm{op}}
\leq
\varepsilon.
\]

By the isometry equivalence lemma,
\[
\sqrt{1-\varepsilon}
\leq
\sigma_m\left(\frac{A}{\sqrt{n}}\right)
\leq
\sigma_1\left(\frac{A}{\sqrt{n}}\right)
\leq
\sqrt{1+\varepsilon}.
\]

Using the elementary implication
\[
|s^2-1|\leq \varepsilon
\quad\Rightarrow\quad
|s-1|\leq \varepsilon\vee\sqrt{\varepsilon},
\]
we obtain
\[
1-CK^2\left(
\sqrt{\frac{m}{n}}+\frac{t}{\sqrt{n}}
\right)
\leq
\sigma_m\left(\frac{A}{\sqrt{n}}\right)
\leq
\sigma_1\left(\frac{A}{\sqrt{n}}\right)
\leq
1+CK^2\left(
\sqrt{\frac{m}{n}}+\frac{t}{\sqrt{n}}
\right).
\]

Multiplying by
\[
\sqrt{n}
\]
gives
\[
\sqrt{n}-CK^2(\sqrt{m}+t)
\leq
\sigma_m(A)
\leq
\sigma_1(A)
\leq
\sqrt{n}+CK^2(\sqrt{m}+t).
\]

This proves the theorem.

\section{Consequences}

\begin{remark}
The proof never directly computed expectations of matrix norms. Instead, it used:
\begin{itemize}
\item a finite net approximation,
\item scalar concentration for each net point,
\item a union bound over the net.
\end{itemize}
\end{remark}

\begin{corollary}
Under the assumptions of the theorem,
\[
\E\norm{A}_{\mathrm{op}}
\leq
CK(\sqrt{n}+\sqrt{m})
\]
and
\[
\E\norm{\frac{1}{n}A^\top A-I_m}_{\mathrm{op}}
\leq
CK^2\left(
\sqrt{\frac{m}{n}}+\frac{m}{n}
\right).
\]
\end{corollary}

\begin{remark}
The operator norm bound is optimal up to constants. A particularly important feature of the theorem is that it only requires independence across rows. The entries within each row may be dependent.
\end{remark}

\chapter{Covariance Estimation and Gaussian Mixture Clustering}

\section{Covariance Estimation and PCA}

Principal component analysis is a method for adaptive dimension reduction.

Suppose
\[
X_1,\dots,X_n\in\R^d
\]
are iid samples from a distribution $\mathcal{D}$ with
\[
\E[X]=0
\]
and covariance matrix
\[
\Sigma=\E[XX^\top].
\]

When $d$ is large, we often want to project onto a lower-dimensional subspace
\[
\mathcal{U}\subseteq\R^d
\]
that captures the most important directions of variation in the data.

\begin{definition}[Sample Covariance]
Given samples
\[
X_1,\dots,X_n\in\R^d,
\]
the empirical covariance matrix is
\[
\Sigma_n
=
\frac{1}{n}
\sum_{i=1}^n X_iX_i^\top.
\]
\end{definition}

\section{Principal Components as Maximum-Variance Directions}

\begin{lemma}
Suppose $X$ is a random vector in $\R^d$ with covariance matrix
\[
\Sigma.
\]
Let the eigenvalues of $\Sigma$ be
\[
\lambda_1\geq \lambda_2\geq \cdots \geq \lambda_d\geq 0,
\]
with corresponding orthonormal eigenvectors
\[
u_1,\dots,u_d.
\]

Then, for every $k\in\{1,\dots,d\}$,
\[
\lambda_k
=
\max_{\substack{
v\perp \operatorname{span}\{u_1,\dots,u_{k-1}\}\\
\norm{v}_2=1
}}
\Var(\inner{X}{v}),
\]
and the maximum is attained at
\[
v=u_k.
\]
\end{lemma}

\begin{proof}
Since
\[
\E[X]=0,
\]
we have
\[
\Var(\inner{X}{v})
=
\E[\inner{X}{v}^2]
=
v^\top \Sigma v.
\]

Using the spectral decomposition,
\[
\Sigma
=
\sum_{j=1}^d \lambda_j u_ju_j^\top.
\]

If
\[
v\perp \operatorname{span}\{u_1,\dots,u_{k-1}\}
\]
and
\[
\norm{v}_2=1,
\]
then
\[
v=\sum_{j=k}^d a_j u_j,
\qquad
\sum_{j=k}^d a_j^2=1.
\]

Therefore,
\[
v^\top\Sigma v
=
\sum_{j=k}^d \lambda_j a_j^2.
\]

Since
\[
\lambda_k\geq \lambda_{k+1}\geq \cdots \geq \lambda_d,
\]
we have
\[
\sum_{j=k}^d \lambda_j a_j^2
\leq
\lambda_k
\sum_{j=k}^d a_j^2
=
\lambda_k.
\]

Equality is attained when
\[
v=u_k.
\]

Thus
\[
\lambda_k
=
\max_{\substack{
v\perp \operatorname{span}\{u_1,\dots,u_{k-1}\}\\
\norm{v}_2=1
}}
\Var(\inner{X}{v}).
\]
\end{proof}

\begin{remark}
This lemma justifies PCA as a maximum-variance dimension reduction method. If
\[
U_k=
\begin{pmatrix}
u_1 & \cdots & u_k
\end{pmatrix},
\]
then the orthogonal projection onto the top-$k$ principal subspace is
\[
P_k=U_kU_k^\top.
\]
\end{remark}

\section{Covariance Estimation}

In practice, we do not know $\Sigma$. Instead, we approximate it using
\[
\Sigma_n
=
\frac{1}{n}
\sum_{i=1}^n X_iX_i^\top.
\]

By the law of large numbers,
\[
\Sigma_n\to \Sigma
\]
almost surely entrywise. The more useful high-dimensional question is:

\begin{quote}
How large must $n$ be so that
\[
\norm{\Sigma_n-\Sigma}_{\mathrm{op}}
\]
is small with high probability?
\end{quote}

\section{Covariance Estimation for Sub-Gaussian Random Vectors}

\begin{theorem}[Covariance Estimation]
Let
\[
X\in\R^d
\]
be a mean-zero sub-Gaussian random vector with
\[
\E[XX^\top]=\Sigma.
\]

Assume there is a constant $K\geq1$ such that
\[
\norm{\rangle X,x \rangle}_{\psi_2}
\leq
K\sqrt{x^\top\Sigma x}
\]
for all
\[
x\in\R^d.
\]

Then
\[
\frac{
\E\norm{\Sigma_n-\Sigma}_{\mathrm{op}}
}{
\norm{\Sigma}_{\mathrm{op}}
}
\leq
CK^2
\left(
\sqrt{\frac{d}{n}}
+
\frac{d}{n}
\right),
\]
where $C>0$ is a universal constant.
\end{theorem}

\begin{proof}
Assume first that $\Sigma$ is invertible. Define the whitened random vector
\[
Z=\Sigma^{-1/2}X.
\]

Then
\[
\E[Z]=0
\]
and
\[
\E[ZZ^\top]
=
\Sigma^{-1/2}\E[XX^\top]\Sigma^{-1/2}
=
I_d.
\]

Moreover, for every
\[
u\in S^{d-1},
\]
\[
\norm{\langle Z,u \rangle }_{\psi_2}
=
\norm{\langle X,\Sigma^{-1/2}u\rangle }_{\psi_2}
\leq
K\sqrt{
u^\top\Sigma^{-1/2}\Sigma\Sigma^{-1/2}u
}
=
K.
\]

Thus $Z$ is isotropic and sub-Gaussian with sub-Gaussian norm bounded by $K$.

Let
\[
Z_i=\Sigma^{-1/2}X_i.
\]

Then
\[
\frac{1}{n}\sum_{i=1}^n Z_iZ_i^\top
=
\Sigma^{-1/2}\Sigma_n\Sigma^{-1/2}.
\]

Therefore,
\[
\Sigma_n-\Sigma
=
\Sigma^{1/2}
\left(
\frac{1}{n}\sum_{i=1}^n Z_iZ_i^\top-I_d
\right)
\Sigma^{1/2}.
\]

Taking operator norms,
\[
\norm{\Sigma_n-\Sigma}_{\mathrm{op}}
\leq
\norm{\Sigma}_{\mathrm{op}}
\left\|
\frac{1}{n}\sum_{i=1}^n Z_iZ_i^\top-I_d
\right\|_{\mathrm{op}}.
\]

By the singular value result for isotropic sub-Gaussian random matrices,
\[
\E
\left\|
\frac{1}{n}\sum_{i=1}^n Z_iZ_i^\top-I_d
\right\|_{\mathrm{op}}
\leq
CK^2
\left(
\sqrt{\frac{d}{n}}
+
\frac{d}{n}
\right).
\]

Thus,
\[
\E\norm{\Sigma_n-\Sigma}_{\mathrm{op}}
\leq
CK^2\norm{\Sigma}_{\mathrm{op}}
\left(
\sqrt{\frac{d}{n}}
+
\frac{d}{n}
\right).
\]

If $\Sigma$ is singular, the same argument applies on the range of $\Sigma$, or by replacing $\Sigma$ with $\Sigma+\eta I_d$ and taking $\eta\downarrow0$.
\end{proof}

\begin{corollary}
Under the assumptions of the previous theorem, there exists a universal constant $C>0$ such that for every
\[
\varepsilon\in(0,1),
\]
if
\[
n
\geq
\frac{CK^4d}{\varepsilon^2},
\]
then
\[
\E\norm{\Sigma_n-\Sigma}_{\mathrm{op}}
\leq
\varepsilon\norm{\Sigma}_{\mathrm{op}}.
\]
\end{corollary}

\begin{proof}
If
\[
n
\geq
\frac{CK^4d}{\varepsilon^2},
\]
with $C$ sufficiently large, then
\[
K^2\sqrt{\frac{d}{n}}
\leq
c\varepsilon
\]
and
\[
K^2\frac{d}{n}
\leq
c\varepsilon
\]
for a sufficiently small universal constant $c>0$.

Substituting into the covariance estimation theorem gives
\[
\E\norm{\Sigma_n-\Sigma}_{\mathrm{op}}
\leq
\varepsilon\norm{\Sigma}_{\mathrm{op}}.
\]
\end{proof}

\section{Gaussian Mixture Clustering}

We now study a clustering problem for point clouds.

\begin{definition}[Two-Component Gaussian Mixture Model]
Let
\[
\mu\in\R^d.
\]

A sample from the two-component Gaussian mixture model $\mathrm{GMM}(\mu)$ is generated as follows:
\begin{enumerate}
\item Draw
\[
S\sim \mathrm{Unif}\{-1,1\}.
\]
\item Draw
\[
g\sim \Normal(0,I_d)
\]
independently of $S$.
\item Set
\[
X=S\mu+g.
\]
\end{enumerate}
\end{definition}

Equivalently,
\[
X\mid S=1\sim \Normal(\mu,I_d),
\qquad
X\mid S=-1\sim \Normal(-\mu,I_d).
\]

Given samples
\[
X_1,\dots,X_n,
\]
the goal is to estimate the hidden labels
\[
S_1,\dots,S_n.
\]

\section{Population Covariance of the Mixture}

For
\[
X=S\mu+g,
\]
we have
\[
\E[X]=0.
\]

The covariance matrix is
\[
\Sigma
=
\E[XX^\top].
\]

Expanding,
\[
XX^\top
=
S^2\mu\mu^\top
+
S\mu g^\top
+
Sg\mu^\top
+
gg^\top.
\]

Taking expectations,
\[
\E[S^2\mu\mu^\top]=\mu\mu^\top,
\]
\[
\E[S\mu g^\top]=\mu\E[S]\E[g]^\top=0,
\]
\[
\E[Sg\mu^\top]=\E[S]\E[g]\mu^\top=0,
\]
and
\[
\E[gg^\top]=I_d.
\]

Thus,
\[
\Sigma
=
I_d+\mu\mu^\top.
\]

The top eigenvector of $\Sigma$ is
\[
u_1
=
\frac{\mu}{\norm{\mu}_2},
\]
with eigenvalue
\[
1+\norm{\mu}_2^2.
\]

Every direction orthogonal to $\mu$ has eigenvalue $1$.

Hence the eigengap is
\[
\Delta
=
\norm{\mu}_2^2.
\]

\section{Spectral Clustering for Gaussian Mixtures}

\begin{algorithm}[H]
\caption{Spectral Clustering for a Two-Component Gaussian Mixture}
\begin{algorithmic}[1]
\State \textbf{Input:} Samples $X_1,\dots,X_n\in\R^d$
\State Compute the sample covariance
\[
\Sigma_n=\frac{1}{n}\sum_{i=1}^n X_iX_i^\top.
\]
\State Compute the top eigenvector
\[
v=u_1(\Sigma_n).
\]
\State For each $i\in\{1,\dots,n\}$, output
\[
\widehat{S}_i=\operatorname{sign}(\inner{X_i}{v}).
\]
\end{algorithmic}
\end{algorithm}

\section{Why the Algorithm Works}

If
\[
v\approx \frac{\mu}{\norm{\mu}_2},
\]
then
\[
\inner{X_i}{v}
=
S_i\inner{\mu}{v}
+
\inner{g_i}{v}.
\]

When $v=\mu/\norm{\mu}_2$, this becomes
\[
\inner{X_i}{v}
=
S_i\norm{\mu}_2+\inner{g_i}{v},
\]
where
\[
\inner{g_i}{v}\sim \Normal(0,1).
\]

Thus the label is correctly recovered whenever
\[
|\inner{g_i}{v}|<\norm{\mu}_2.
\]

If
\[
\norm{\mu}_2
\]
is a sufficiently large constant, then this happens with high probability for most points.

\section{Clustering Guarantee}

\begin{theorem}
Let
\[
X_1,\dots,X_n\in\R^d
\]
be iid samples from
\[
\mathrm{GMM}(\mu).
\]

There exist universal constants
\[
C,c>0
\]
such that if
\[
n\geq Cd
\]
and
\[
\norm{\mu}_2\geq C,
\]
then with probability at least
\[
0.99,
\]
the spectral clustering algorithm misclassifies at most
\[
1\%
\]
of the points.
\end{theorem}

\begin{proof}
The population covariance is
\[
\Sigma=I_d+\mu\mu^\top.
\]

Its top eigenvector is
\[
u=\frac{\mu}{\norm{\mu}_2},
\]
and the eigengap is
\[
\Delta=\norm{\mu}_2^2.
\]

By covariance estimation, when
\[
n\geq Cd,
\]
we have with high probability
\[
\norm{\Sigma_n-\Sigma}_{\mathrm{op}}
\leq
c\norm{\mu}_2^2
\]
for a sufficiently small universal constant $c>0$, provided $C$ is chosen large enough.

By Davis--Kahan,
\[
\min_{s\in\{-1,1\}}
\norm{v-su}_2
\leq
C\frac{\norm{\Sigma_n-\Sigma}_{\mathrm{op}}}{\Delta}.
\]

Thus, after choosing the sign of $v$ appropriately,
\[
\norm{v-u}_2
\leq
\eta
\]
for a small constant $\eta>0$.

Now decompose the empirical classification score:
\[
\inner{X_i}{v}
=
S_i\inner{\mu}{v}
+
\inner{g_i}{v}.
\]

Since
\[
\norm{v-u}_2\leq \eta,
\]
we have
\[
\inner{\mu}{v}
=
\norm{\mu}_2\inner{u}{v}
\geq
\norm{\mu}_2(1-C\eta^2).
\]

For $\eta$ sufficiently small,
\[
\inner{\mu}{v}
\geq
\frac{1}{2}\norm{\mu}_2.
\]

Thus point $i$ is misclassified only if
\[
|\inner{g_i}{v}|
\geq
\frac{1}{2}\norm{\mu}_2.
\]

Conditioned on $v$, the random variable $\inner{g_i}{v}$ is approximately standard Gaussian. More directly, since $g_i$ is isotropic Gaussian and $v$ is a unit vector,
\[
\inner{g_i}{v}\sim \Normal(0,1).
\]

Therefore,
\[
\Prob\left(
|\inner{g_i}{v}|
\geq
\frac{1}{2}\norm{\mu}_2
\right)
\leq
2\exp\left(
-\frac{\norm{\mu}_2^2}{8}
\right).
\]

Choosing
\[
\norm{\mu}_2\geq C
\]
large enough makes this probability at most
\[
0.001.
\]

By a concentration bound for Bernoulli sums, with probability at least $0.99$, the fraction of indices for which this event occurs is at most
\[
0.01.
\]

Therefore, the spectral clustering algorithm misclassifies at most $1\%$ of the points.
\end{proof}

\begin{remark}
A point cloud drawn from
\[
\Normal(0,I_d)
\]
has diameter on the order of
\[
\sqrt{d}.
\]

Nevertheless, in the two-component Gaussian mixture model, a constant separation
\[
\norm{\mu}_2\gtrsim 1
\]
is enough for classification when
\[
n\gtrsim d.
\]
\end{remark}

\begin{remark}
The condition
\[
n\gtrsim d
\]
is optimal in this setting. If the within-cluster covariance is unknown, the clustering problem becomes substantially more delicate.
\end{remark}

\chapter{Matrix Concentration}

\section{Matrix Calculus}

Our goal is to develop matrix analogues of scalar concentration inequalities. In particular, we want bounds of the form
\[
\Prob\left(
\norm{\sum_{i=1}^n A_i}_{\mathrm{op}}\geq t
\right),
\]
where
\[
A_1,\dots,A_n
\]
are independent mean-zero random symmetric matrices.

\section{Functions of Symmetric Matrices}

\begin{definition}[Functions of Symmetric Matrices]
Let
\[
f:\R\to\R
\]
and let
\[
X\in\mathbb{S}^d
\]
have spectral decomposition
\[
X=\sum_{i=1}^d \lambda_i u_i u_i^\top.
\]

Then define
\[
f(X)
=
\sum_{i=1}^d f(\lambda_i)u_i u_i^\top.
\]
\end{definition}

\begin{example}
If
\[
f(t)=t^{-1},
\]
then
\[
f(X)=X^{-1},
\]
provided $X$ is invertible.
\end{example}

\begin{example}
If
\[
f(t)=\sum_{k=0}^{\infty} a_k t^k,
\]
then
\[
f(X)=\sum_{k=0}^{\infty} a_k X^k,
\]
whenever the series converges.
\end{example}

\begin{example}
The matrix exponential is defined by
\[
e^X
=
\sum_{k=0}^{\infty}\frac{X^k}{k!}.
\]
Equivalently, if
\[
X=\sum_{i=1}^d \lambda_i u_i u_i^\top,
\]
then
\[
e^X
=
\sum_{i=1}^d e^{\lambda_i}u_i u_i^\top.
\]
\end{example}

\section{The Loewner Order}

\begin{definition}[Loewner Order]
For
\[
X,Y\in\mathbb{S}^d,
\]
we write
\[
X\preceq Y
\]
if
\[
Y-X\in\mathbb{S}_+^d.
\]

Equivalently,
\[
X\preceq Y
\]
means
\[
v^\top Xv\leq v^\top Yv
\]
for all
\[
v\in\R^d.
\]
\end{definition}

\begin{lemma}
Suppose
\[
X\preceq Y.
\]
Then:

\begin{enumerate}
\item For every
\[
A\in\R^{d\times k},
\]
\[
A^\top X A
\preceq
A^\top Y A.
\]

\item For every $i$,
\[
\lambda_i(X)\leq \lambda_i(Y),
\]
where eigenvalues are ordered decreasingly.

\item If
\[
f:\R\to\R
\]
is nondecreasing, then
\[
\tr(f(X))\leq \tr(f(Y)).
\]
\end{enumerate}
\end{lemma}

\begin{proof}
For the first claim, since
\[
Y-X\succeq 0,
\]
we have
\[
A^\top(Y-X)A\succeq0.
\]
Thus,
\[
A^\top XA\preceq A^\top YA.
\]

For the second claim, use the Courant--Fischer variational characterization:
\[
\lambda_i(X)
=
\max_{\dim(L)=i}
\min_{\substack{x\in L\\ \norm{x}_2=1}}
x^\top Xx.
\]

Since
\[
x^\top Xx\leq x^\top Yx
\]
for every unit vector $x$, the variational characterization gives
\[
\lambda_i(X)\leq \lambda_i(Y).
\]

For the third claim, since $f$ is nondecreasing,
\[
f(\lambda_i(X))\leq f(\lambda_i(Y)).
\]

Therefore,
\[
\tr(f(X))
=
\sum_{i=1}^d f(\lambda_i(X))
\leq
\sum_{i=1}^d f(\lambda_i(Y))
=
\tr(f(Y)).
\]
\end{proof}

\begin{remark}
It is tempting to believe that
\[
X\preceq Y
\quad\Rightarrow\quad
f(X)\preceq f(Y)
\]
for every nondecreasing function $f$. This is false in general. Functions with this stronger property are called matrix monotone.
\end{remark}

\begin{definition}[Matrix Monotone Function]
A function
\[
f:\R\to\R
\]
is called matrix monotone if, for all symmetric matrices $X,Y$ with spectra in the domain of $f$,
\[
X\preceq Y
\quad\Rightarrow\quad
f(X)\preceq f(Y).
\]
\end{definition}

\begin{lemma}
The functions
\[
t\mapsto t^{-1},
\qquad
t\mapsto t^{1/2},
\qquad
t\mapsto \log t
\]
are matrix monotone on their appropriate positive domains.
\end{lemma}

\section{Matrix Moment Generating Functions}

For scalar random variables, moment generating functions lead to tail bounds through Markov's inequality. We now define the analogous object for random matrices.

\begin{definition}[Matrix Moment Generating Function]
Let
\[
A\in\mathbb{S}^d
\]
be a random symmetric matrix. Its matrix moment generating function is
\[
\Psi_A(\gamma)
=
\E[e^{\gamma A}],
\]
whenever this expectation is well-defined.
\end{definition}

Using the matrix exponential,
\[
\Psi_A(\gamma)
=
\E\left[
\sum_{k=0}^{\infty}\frac{\gamma^k A^k}{k!}
\right]
=
\sum_{k=0}^{\infty}\frac{\gamma^k\E[A^k]}{k!}.
\]

\section{Matrix Chernoff Method}

\begin{lemma}[Matrix Chernoff Method]
Let
\[
A\in\mathbb{S}^d
\]
be a random symmetric matrix whose matrix moment generating function is defined on an interval
\[
(-a,a).
\]

Then, for every
\[
t>0
\]
and every
\[
\gamma\in[0,a),
\]
\[
\Prob(\lambda_1(A)\geq t)
\leq
\tr(\Psi_A(\gamma))e^{-\gamma t}.
\]

Consequently,
\[
\Prob(\norm{A}_{\mathrm{op}}\geq t)
\leq
2\tr(\Psi_A(\gamma))e^{-\gamma t}.
\]
\end{lemma}

\begin{proof}
Since the scalar exponential is increasing,
\[
\{\lambda_1(A)\geq t\}
=
\{e^{\gamma\lambda_1(A)}\geq e^{\gamma t}\}.
\]

Because
\[
\lambda_1(e^{\gamma A})
=
e^{\gamma\lambda_1(A)},
\]
Markov's inequality gives
\[
\Prob(\lambda_1(A)\geq t)
\leq
e^{-\gamma t}
\E[\lambda_1(e^{\gamma A})].
\]

Since
\[
\lambda_1(B)\leq \tr(B)
\]
for every positive semidefinite matrix $B$, we get
\[
\Prob(\lambda_1(A)\geq t)
\leq
e^{-\gamma t}
\E[\tr(e^{\gamma A})].
\]

By linearity of trace and expectation,
\[
\E[\tr(e^{\gamma A})]
=
\tr(\E[e^{\gamma A}])
=
\tr(\Psi_A(\gamma)).
\]

Thus,
\[
\Prob(\lambda_1(A)\geq t)
\leq
\tr(\Psi_A(\gamma))e^{-\gamma t}.
\]

For the two-sided operator norm bound, note that
\[
\norm{A}_{\mathrm{op}}
=
\max\{\lambda_1(A),-\lambda_d(A)\}.
\]

Apply the previous bound to $A$ and to $-A$, and use the union bound.
\end{proof}

\section{Matrix Sub-Gaussian Random Variables}

\begin{definition}[$V$-Sub-Gaussian Random Matrix]
A random symmetric matrix
\[
A\in\mathbb{S}^d
\]
is called $V$-sub-Gaussian if
\[
\Psi_A(\lambda)
\preceq
e^{\lambda^2 V/2}
\]
for every
\[
\lambda\in\R.
\]
\end{definition}

\begin{example}
Suppose
\[
A=\varepsilon B,
\]
where
\[
\varepsilon\sim \mathrm{Unif}\{-1,1\}
\]
and
\[
B\in\mathbb{S}^d
\]
is fixed.

Then
\[
\E[A^{2k+1}]=0,
\qquad
\E[A^{2k}]=B^{2k}.
\]

Therefore,
\[
\E[e^{\lambda A}]
=
\sum_{k=0}^{\infty}
\frac{\lambda^{2k}B^{2k}}{(2k)!}.
\]

Using
\[
(2k)!\geq 2^k k!,
\]
we obtain
\[
\E[e^{\lambda A}]
\preceq
\sum_{k=0}^{\infty}
\frac{1}{k!}
\left(
\frac{\lambda^2B^2}{2}
\right)^k
=
e^{\lambda^2B^2/2}.
\]

Thus $A$ is $B^2$-sub-Gaussian.
\end{example}

\section{Lieb's Inequality}

The scalar proof that independent sub-Gaussian random variables remain sub-Gaussian under summation relies on the identity
\[
e^{a+b}=e^ae^b.
\]

For matrices, this identity fails in general:
\[
e^{A+B}\neq e^Ae^B
\]
unless $A$ and $B$ commute. To handle this issue, we use Lieb's inequality.

\begin{theorem}[Lieb's Inequality]
Let
\[
H\in\mathbb{S}^d.
\]

Define
\[
f:\mathbb{S}_+^d\to\R
\]
by
\[
f(X)
=
\tr\exp(H+\log X).
\]

Then $f$ is concave on
\[
\mathbb{S}_+^d.
\]
\end{theorem}

\section{Matrix MGF Bound for Sums}

\begin{lemma}
Let
\[
A_1,\dots,A_n\in\mathbb{S}^d
\]
be independent random symmetric matrices. Let
\[
S_n=\sum_{i=1}^n A_i.
\]

Assume that
\[
\Psi_{A_i}(\gamma)
\]
is defined for every
\[
\gamma\in J\subseteq\R.
\]

Then
\[
\tr(\Psi_{S_n}(\gamma))
\leq
\tr\left(
\exp\left(
\sum_{i=1}^n \log \Psi_{A_i}(\gamma)
\right)
\right)
\]
for every
\[
\gamma\in J.
\]
\end{lemma}

\begin{proof}
We prove the result by conditioning one random matrix at a time.

Write
\[
S_n=S_{n-1}+A_n.
\]

Then
\[
\tr(\Psi_{S_n}(\gamma))
=
\tr\E[e^{\gamma S_n}]
=
\E\tr(e^{\gamma S_{n-1}+\gamma A_n}).
\]

Condition on
\[
S_{n-1}.
\]

Since
\[
e^{\gamma A_n}
=
\exp(\gamma A_n),
\]
we may write
\[
\gamma A_n
=
\log(e^{\gamma A_n}).
\]

Using Lieb's inequality and Jensen's inequality conditionally,
\[
\E_{A_n}
\tr\exp(\gamma S_{n-1}+\log(e^{\gamma A_n}))
\leq
\tr\exp(\gamma S_{n-1}+\log \E[e^{\gamma A_n}]).
\]

Thus,
\[
\tr(\Psi_{S_n}(\gamma))
\leq
\E\tr\exp(\gamma S_{n-1}+\log\Psi_{A_n}(\gamma)).
\]

Repeating the same argument recursively gives
\[
\tr(\Psi_{S_n}(\gamma))
\leq
\tr\exp\left(
\sum_{i=1}^n \log\Psi_{A_i}(\gamma)
\right).
\]
\end{proof}

\section{Matrix Hoeffding Inequality}

\begin{theorem}[Matrix Hoeffding Inequality]
Let
\[
A_1,\dots,A_n\in\mathbb{S}^d
\]
be independent, mean-zero random matrices.

Assume that each
\[
A_i
\]
is $V_i$-sub-Gaussian. Define
\[
\sigma^2
=
\norm{\frac{1}{n}\sum_{i=1}^n V_i}_{\mathrm{op}}.
\]

Then, for every
\[
t>0,
\]
\[
\Prob\left(
\norm{\frac{1}{n}\sum_{i=1}^n A_i}_{\mathrm{op}}
\geq t
\right)
\leq
2\rank\left(\sum_{i=1}^n V_i\right)
\exp\left(
-\frac{nt^2}{2\sigma^2}
\right).
\]

In particular,
\[
\Prob\left(
\norm{\frac{1}{n}\sum_{i=1}^n A_i}_{\mathrm{op}}
\geq t
\right)
\leq
2d\exp\left(
-\frac{nt^2}{2\sigma^2}
\right).
\]
\end{theorem}

\begin{proof}
Let
\[
V=\sum_{i=1}^n V_i.
\]

By the MGF bound for sums,
\[
\tr(\Psi_{S_n}(\gamma))
\leq
\tr\exp\left(
\sum_{i=1}^n \log\Psi_{A_i}(\gamma)
\right),
\]
where
\[
S_n=\sum_{i=1}^n A_i.
\]

Since $A_i$ is $V_i$-sub-Gaussian,
\[
\Psi_{A_i}(\gamma)
\preceq
e^{\gamma^2V_i/2}.
\]

Because the matrix logarithm is matrix monotone,
\[
\log\Psi_{A_i}(\gamma)
\preceq
\frac{\gamma^2}{2}V_i.
\]

Hence,
\[
\sum_{i=1}^n\log\Psi_{A_i}(\gamma)
\preceq
\frac{\gamma^2}{2}
\sum_{i=1}^n V_i
=
\frac{\gamma^2}{2}V.
\]

Since
\[
X\preceq Y
\quad\Rightarrow\quad
\tr(e^X)\leq \tr(e^Y),
\]
we get
\[
\tr(\Psi_{S_n}(\gamma))
\leq
\tr\left(
e^{\gamma^2V/2}
\right).
\]

By the matrix Chernoff method,
\[
\Prob\left(
\norm{\frac{1}{n}S_n}_{\mathrm{op}}\geq t
\right)
\leq
2\tr\left(
e^{\gamma^2V/2}
\right)e^{-\gamma nt}.
\]

Now,
\[
\tr(e^{\gamma^2V/2})
\leq
\rank(V)
\exp\left(
\frac{\gamma^2}{2}\norm{V}_{\mathrm{op}}
\right).
\]

Since
\[
\norm{V}_{\mathrm{op}}
=
n\sigma^2,
\]
we obtain
\[
\Prob\left(
\norm{\frac{1}{n}S_n}_{\mathrm{op}}\geq t
\right)
\leq
2\rank(V)
\exp\left(
\frac{\gamma^2 n\sigma^2}{2}
-
\gamma nt
\right).
\]

The exponent is minimized by
\[
\gamma=\frac{t}{\sigma^2}.
\]

Substituting this value gives
\[
\frac{\gamma^2 n\sigma^2}{2}
-
\gamma nt
=
-\frac{nt^2}{2\sigma^2}.
\]

Therefore,
\[
\Prob\left(
\norm{\frac{1}{n}\sum_{i=1}^n A_i}_{\mathrm{op}}
\geq t
\right)
\leq
2\rank(V)
\exp\left(
-\frac{nt^2}{2\sigma^2}
\right).
\]

Since
\[
\rank(V)\leq d,
\]
the second bound follows.
\end{proof}

\begin{remark}
Compared to the scalar Hoeffding inequality, the matrix version contains an additional dimension factor. In general, this factor is unavoidable.
\end{remark}

\chapter{Parametric Estimation and Maximum Likelihood}

\section{Introduction to Estimation}

Suppose we observe iid data
\[
Z_1,\dots,Z_n\sim P_\theta,
\]
where
\[
P_\theta
\]
is a probability distribution depending on an unknown object
\[
\theta.
\]

The goal of estimation is to construct a data-dependent rule that approximates
\[
\theta
\]
accurately.

Common examples include estimating:
\begin{itemize}
\item a mean vector,
\item a covariance matrix,
\item a cumulative distribution function,
\item a density,
\item a regression parameter.
\end{itemize}

\section{Parametric and Nonparametric Estimation}

\begin{definition}[Parametric Estimation]
A statistical estimation problem is called parametric if
\[
\theta\in\R^k
\]
for some finite
\[
k,
\]
and the family
\[
\{P_\theta:\theta\in\Theta\}
\]
is described by a finite-dimensional parameter.
\end{definition}

\begin{example}
Suppose
\[
Z_i\sim \Normal(\theta,I_d),
\qquad
\theta\in\R^d.
\]
Then estimating
\[
\theta
\]
is a parametric estimation problem.
\end{example}

\begin{definition}[Nonparametric Estimation]
A statistical estimation problem is called nonparametric if the object being estimated does not belong to a fixed finite-dimensional family.
\end{definition}

\begin{example}
If the goal is to estimate an unknown density
\[
f
\]
over
\[
\R^d
\]
without assuming that
\[
f
\]
belongs to a finite-dimensional parametric family, then the problem is nonparametric.
\end{example}

\section{Estimators}

\begin{definition}[Estimator]
An estimator of
\[
\theta
\]
is a measurable function of the data:
\[
\widehat{\theta}_n
=
g(Z_1,\dots,Z_n).
\]
\end{definition}

\begin{definition}[Unbiased Estimator]
An estimator
\[
\widehat{\theta}_n
\]
is unbiased if
\[
\E[\widehat{\theta}_n]=\theta.
\]
\end{definition}

\section{Mean Squared Error}

A common way to measure accuracy is the mean squared error.

\begin{definition}[Mean Squared Error]
The mean squared error of an estimator
\[
\widehat{\theta}_n
\]
is
\[
\operatorname{MSE}(\widehat{\theta}_n)
=
\E\norm{\widehat{\theta}_n-\theta}_2^2.
\]
\end{definition}

The mean squared error decomposes into bias and variance.

\begin{proposition}[Bias-Variance Decomposition]
For any estimator
\[
\widehat{\theta}_n,
\]
\[
\operatorname{MSE}(\widehat{\theta}_n)
=
\norm{\E[\widehat{\theta}_n]-\theta}_2^2
+
\E\norm{\widehat{\theta}_n-\E[\widehat{\theta}_n]}_2^2.
\]
\end{proposition}

\begin{proof}
Write
\[
\widehat{\theta}_n-\theta
=
\widehat{\theta}_n-\E[\widehat{\theta}_n]
+
\E[\widehat{\theta}_n]-\theta.
\]

Then
\[
\norm{\widehat{\theta}_n-\theta}_2^2
=
\norm{\widehat{\theta}_n-\E[\widehat{\theta}_n]}_2^2
+
\norm{\E[\widehat{\theta}_n]-\theta}_2^2
+
2
\left\langle
\widehat{\theta}_n-\E[\widehat{\theta}_n],
\E[\widehat{\theta}_n]-\theta
\right\rangle.
\]

Taking expectations, the cross term vanishes because
\[
\E[\widehat{\theta}_n-\E[\widehat{\theta}_n]]=0.
\]

Therefore,
\[
\E\norm{\widehat{\theta}_n-\theta}_2^2
=
\E\norm{\widehat{\theta}_n-\E[\widehat{\theta}_n]}_2^2
+
\norm{\E[\widehat{\theta}_n]-\theta}_2^2.
\]
\end{proof}

\begin{definition}[Bias and Variance]
The squared bias is
\[
\operatorname{Bias}(\widehat{\theta}_n)^2
=
\norm{\E[\widehat{\theta}_n]-\theta}_2^2.
\]

The variance term is
\[
\operatorname{Var}(\widehat{\theta}_n)
=
\E\norm{\widehat{\theta}_n-\E[\widehat{\theta}_n]}_2^2.
\]

Thus,
\[
\operatorname{MSE}(\widehat{\theta}_n)
=
\operatorname{Bias}(\widehat{\theta}_n)^2
+
\operatorname{Var}(\widehat{\theta}_n).
\]
\end{definition}

\begin{remark}
Unbiasedness is not always optimal for mean squared error. Sometimes it is beneficial to introduce bias if the variance reduction is large enough.
\end{remark}

\section{Maximum Likelihood Estimation}

Maximum likelihood estimation gives a systematic way to construct estimators in parametric models.

Suppose
\[
\{f(z;\theta):\theta\in\Theta\}
\]
is a parametric family of densities or probability mass functions.

Given observations
\[
Z_1,\dots,Z_n,
\]
the likelihood function is
\[
L_n(\theta)
=
\prod_{i=1}^n f(Z_i;\theta).
\]

The log-likelihood function is
\[
\ell_n(\theta)
=
\sum_{i=1}^n \log f(Z_i;\theta).
\]

\begin{definition}[Maximum Likelihood Estimator]
The maximum likelihood estimator is any solution
\[
\widehat{\theta}_{\mathrm{MLE}}
\in
\argmax_{\theta\in\Theta} L_n(\theta).
\]

Equivalently,
\[
\widehat{\theta}_{\mathrm{MLE}}
\in
\argmax_{\theta\in\Theta} \ell_n(\theta).
\]
\end{definition}

\begin{remark}
Maximizing the likelihood and maximizing the log-likelihood are equivalent because the logarithm is strictly increasing.
\end{remark}

\section{Example: Bernoulli Model}

Suppose
\[
Z_1,\dots,Z_n\iid \Bern(p),
\]
where
\[
p\in(0,1).
\]

The probability mass function is
\[
f(z;p)
=
p^z(1-p)^{1-z},
\qquad
z\in\{0,1\}.
\]

Thus,
\[
L_n(p)
=
\prod_{i=1}^n p^{Z_i}(1-p)^{1-Z_i}.
\]

The log-likelihood is
\[
\ell_n(p)
=
\sum_{i=1}^n Z_i\log p
+
\sum_{i=1}^n (1-Z_i)\log(1-p).
\]

Let
\[
S_n=\sum_{i=1}^n Z_i.
\]

Then
\[
\ell_n(p)
=
S_n\log p
+
(n-S_n)\log(1-p).
\]

Differentiating,
\[
\ell_n'(p)
=
\frac{S_n}{p}
-
\frac{n-S_n}{1-p}.
\]

Setting
\[
\ell_n'(p)=0
\]
gives
\[
\frac{S_n}{p}
=
\frac{n-S_n}{1-p}.
\]

Therefore,
\[
S_n(1-p)=p(n-S_n).
\]

Expanding,
\[
S_n-S_np=np-S_np,
\]
so
\[
S_n=np.
\]

Hence,
\[
\widehat{p}_{\mathrm{MLE}}
=
\frac{S_n}{n}
=
\frac{1}{n}\sum_{i=1}^n Z_i.
\]

\begin{remark}
The Bernoulli maximum likelihood estimator is the sample proportion.
\end{remark}

\section{Example: Linear Regression with Gaussian Noise}

Suppose
\[
\theta\in\R^d
\]
and we observe iid pairs
\[
Z_i=(X_i,y_i),
\]
where
\[
y_i
=
\theta^\top X_i+\varepsilon_i,
\qquad
\varepsilon_i\sim \Normal(0,\sigma^2).
\]

Equivalently,
\[
y_i|X_i\sim \Normal(\theta^\top X_i,\sigma^2).
\]

The parameter of interest is
\[
\theta.
\]

The variance
\[
\sigma^2
\]
may be treated as a nuisance parameter if the goal is only to estimate
\[
\theta.
\]

The conditional density of
\[
y_i
\]
given
\[
X_i
\]
is
\[
f(y_i|X_i;\theta)
=
\frac{1}{\sqrt{2\pi\sigma^2}}
\exp\left(
-\frac{(y_i-\theta^\top X_i)^2}{2\sigma^2}
\right).
\]

Thus, the likelihood is
\[
L_n(\theta)
=
\prod_{i=1}^n
\frac{1}{\sqrt{2\pi\sigma^2}}
\exp\left(
-\frac{(y_i-\theta^\top X_i)^2}{2\sigma^2}
\right).
\]

Therefore,
\[
L_n(\theta)
=
(2\pi\sigma^2)^{-n/2}
\exp\left(
-\frac{1}{2\sigma^2}
\sum_{i=1}^n
(y_i-\theta^\top X_i)^2
\right).
\]

Taking logs,
\[
\ell_n(\theta)
=
-\frac{n}{2}\log(2\pi\sigma^2)
-
\frac{1}{2\sigma^2}
\sum_{i=1}^n
(y_i-\theta^\top X_i)^2.
\]

The first term does not depend on
\[
\theta.
\]

Since
\[
\sigma^2>0,
\]
maximizing
\[
\ell_n(\theta)
\]
over
\[
\theta
\]
is equivalent to minimizing
\[
\sum_{i=1}^n
(y_i-\theta^\top X_i)^2.
\]

Thus,
\[
\widehat{\theta}_{\mathrm{MLE}}
\in
\argmin_{\theta\in\R^d}
\frac{1}{2n}
\sum_{i=1}^n
(y_i-\theta^\top X_i)^2.
\]

This is precisely the ordinary least squares estimator.

\begin{remark}
This example shows that least squares can be derived as maximum likelihood estimation under Gaussian noise.
\end{remark}

\section{Closed Form for Least Squares}

Let
\[
X
=
\begin{pmatrix}
X_1^\top\\
\vdots\\
X_n^\top
\end{pmatrix}
\in\R^{n\times d},
\qquad
y
=
\begin{pmatrix}
y_1\\
\vdots\\
y_n
\end{pmatrix}.
\]

The least squares objective is
\[
Q(\theta)
=
\frac{1}{2n}\norm{y-X\theta}_2^2.
\]

Expanding,
\[
Q(\theta)
=
\frac{1}{2n}
(y-X\theta)^\top(y-X\theta).
\]

Differentiating,
\[
\nabla Q(\theta)
=
-\frac{1}{n}X^\top(y-X\theta).
\]

Setting
\[
\nabla Q(\theta)=0
\]
gives the normal equations
\[
X^\top X\theta=X^\top y.
\]

If
\[
X^\top X
\]
is invertible, then
\[
\widehat{\theta}_{\mathrm{LS}}
=
(X^\top X)^{-1}X^\top y.
\]

\section{Statistical Inference}

Beyond estimating a parameter, one may also want to:
\begin{itemize}
\item construct confidence intervals,
\item test hypotheses,
\item provide high-probability guarantees,
\item quantify uncertainty.
\end{itemize}

These tasks are collectively known as statistical inference.

\section{Maximum Likelihood as Empirical Risk Minimization}

The maximum likelihood estimator can be written as an optimization problem.

Since maximizing
\[
\ell_n(\theta)
=
\sum_{i=1}^n \log f(Z_i;\theta)
\]
is equivalent to minimizing the negative log-likelihood,
\[
\widehat{\theta}_{\mathrm{MLE}}
\in
\argmin_{\theta\in\Theta}
\frac{1}{n}
\sum_{i=1}^n
\bigl[-\log f(Z_i;\theta)\bigr].
\]

This is an example of empirical risk minimization.

\begin{definition}[Stochastic Optimization Problem]
Let
\[
Z\sim \mathcal{D}
\]
and let
\[
F:\Theta\times\mathcal{Z}\to\R.
\]

The population optimization problem is
\[
\min_{\theta\in\Theta}
\E_{Z\sim\mathcal{D}}[F(\theta,Z)].
\]
\end{definition}

Since the distribution
\[
\mathcal{D}
\]
is typically unknown, we replace the population objective with the empirical objective.

\begin{definition}[Empirical Risk Minimization]
Given iid samples
\[
Z_1,\dots,Z_n\sim\mathcal{D},
\]
the empirical risk minimization problem is
\[
\min_{\theta\in\Theta}
\frac{1}{n}
\sum_{i=1}^n F(\theta,Z_i).
\]
\end{definition}

Maximum likelihood is empirical risk minimization with
\[
F(\theta,z)
=
-\log f(z;\theta).
\]

Least squares is empirical risk minimization with
\[
F(\theta,(x,y))
=
\frac{1}{2}(y-\theta^\top x)^2.
\]

\section{Core Questions}

The empirical risk minimization framework raises several important questions:

\begin{itemize}
\item How do we solve the empirical optimization problem efficiently?
\item If
\[
\widehat{\theta}_n
\]
solves the empirical problem and
\[
\theta_*
\]
solves the population problem, do we have
\[
\widehat{\theta}_n\to\theta_*
\]
as
\[
n\to\infty?
\]
\item What convergence rates are possible?
\item Can we obtain confidence intervals or high-probability guarantees?
\end{itemize}

The next major goal is to study these questions in the context of linear least squares.

\chapter{Ordinary Least Squares and Excess Risk}

\section{Ordinary Least Squares}

Suppose we observe data
\[
(X_1,y_1),\dots,(X_n,y_n),
\]
where
\[
X_i\in\R^d,
\qquad
y_i\in\R,
\]
and
\[
y_i=\theta_*^\top X_i+\varepsilon_i.
\]

Assume
\[
\varepsilon_i\iid \Normal(0,\sigma^2).
\]

The empirical risk is
\[
R_n(\theta)
=
\frac{1}{n}\sum_{i=1}^n
(y_i-X_i^\top\theta)^2.
\]

In matrix notation, define
\[
y=
\begin{pmatrix}
y_1\\
\vdots\\
y_n
\end{pmatrix},
\qquad
X=
\begin{pmatrix}
X_1^\top\\
\vdots\\
X_n^\top
\end{pmatrix}.
\]

Then
\[
R_n(\theta)
=
\frac{1}{n}\norm{y-X\theta}_2^2.
\]

Assume throughout this section that
\[
d\leq n
\]
and that $X$ has rank $d$.

\begin{definition}[Ordinary Least Squares Estimator]
The ordinary least squares estimator is any solution of
\[
\widehat{\theta}_{\mathrm{OLS}}
\in
\argmin_{\theta\in\R^d}
R_n(\theta).
\]
\end{definition}

\begin{lemma}
The ordinary least squares estimator exists, is unique, and is given by
\[
\widehat{\theta}_{\mathrm{OLS}}
=
(X^\top X)^{-1}X^\top y.
\]
\end{lemma}

\begin{proof}
Since $X$ has rank $d$,
\[
X^\top X
\]
is positive definite. Hence
\[
R_n(\theta)
=
\frac{1}{n}\norm{y-X\theta}_2^2
\]
is strongly convex in $\theta$, so the minimizer exists and is unique.

Differentiating,
\[
\nabla R_n(\theta)
=
\frac{2}{n}X^\top(X\theta-y).
\]

At the minimizer,
\[
0
=
\nabla R_n(\widehat{\theta}_{\mathrm{OLS}})
=
\frac{2}{n}X^\top(X\widehat{\theta}_{\mathrm{OLS}}-y).
\]

Thus,
\[
X^\top X\widehat{\theta}_{\mathrm{OLS}}
=
X^\top y.
\]

Since $X^\top X$ is invertible,
\[
\widehat{\theta}_{\mathrm{OLS}}
=
(X^\top X)^{-1}X^\top y.
\]
\end{proof}

\section{Geometric Interpretation of OLS}

\begin{lemma}
The fitted values
\[
X\widehat{\theta}_{\mathrm{OLS}}
\]
are the orthogonal projection of $y$ onto
\[
\operatorname{range}(X)\subseteq\R^n.
\]
Specifically,
\[
X\widehat{\theta}_{\mathrm{OLS}}
=
X(X^\top X)^{-1}X^\top y.
\]
\end{lemma}

\begin{proof}
Define
\[
P_X
=
X(X^\top X)^{-1}X^\top.
\]

We show that $P_X$ is the orthogonal projector onto $\operatorname{range}(X)$.

First,
\[
P_X^\top
=
\left(
X(X^\top X)^{-1}X^\top
\right)^\top
=
X(X^\top X)^{-1}X^\top
=
P_X.
\]

Second,
\[
P_X^2
=
X(X^\top X)^{-1}X^\top
X(X^\top X)^{-1}X^\top
=
X(X^\top X)^{-1}X^\top
=
P_X.
\]

Thus $P_X$ is an orthogonal projector. Since
\[
P_Xy\in\operatorname{range}(X),
\]
and
\[
y-P_Xy
\]
is orthogonal to $\operatorname{range}(X)$, $P_Xy$ is the orthogonal projection of $y$ onto $\operatorname{range}(X)$.

Finally,
\[
X\widehat{\theta}_{\mathrm{OLS}}
=
X(X^\top X)^{-1}X^\top y
=
P_Xy.
\]
\end{proof}

\begin{remark}
The OLS estimator can be interpreted in two steps:
\begin{enumerate}
\item project $y$ onto the column space of $X$;
\item solve
\[
X\theta=\widehat{y}
\]
inside that column space.
\end{enumerate}
\end{remark}

\section{Population Risk and Excess Risk}

The empirical risk approximates the population risk
\[
R(\theta)
=
\E[R_n(\theta)].
\]

In the fixed-design setting, $X$ is treated as deterministic and the expectation is only over the noise
\[
\varepsilon=
\begin{pmatrix}
\varepsilon_1\\
\vdots\\
\varepsilon_n
\end{pmatrix}.
\]

The excess risk of an estimator is
\[
R(\widehat{\theta})-\min_{\theta\in\R^d}R(\theta).
\]

\section{Fixed-Design Risk Geometry}

Define
\[
\widehat{\Sigma}
=
\frac{1}{n}X^\top X.
\]

Since $X$ has rank $d$,
\[
\widehat{\Sigma}
\]
is positive definite.

This matrix induces an inner product
\[
\langle \theta,\theta'\rangle_{\widehat{\Sigma}}
=
\theta^\top\widehat{\Sigma}\theta'
\]
and corresponding norm
\[
\norm{\theta}_{\widehat{\Sigma}}^2
=
\theta^\top\widehat{\Sigma}\theta
=
\frac{1}{n}\norm{X\theta}_2^2.
\]

\begin{lemma}
Suppose $X$ is fixed. Then, for every
\[
\theta\in\R^d,
\]
\[
R(\theta)-R^*
=
\norm{\theta-\theta_*}_{\widehat{\Sigma}}^2,
\]
where
\[
R^*
=
\min_{\theta\in\R^d}R(\theta)
=
\sigma^2.
\]
\end{lemma}

\begin{proof}
Since
\[
y=X\theta_*+\varepsilon,
\]
we have
\[
y-X\theta
=
X(\theta_*-\theta)+\varepsilon.
\]

Therefore,
\[
R(\theta)
=
\frac{1}{n}
\E_\varepsilon
\norm{X(\theta_*-\theta)+\varepsilon}_2^2.
\]

Expanding,
\[
R(\theta)
=
\frac{1}{n}\norm{X(\theta_*-\theta)}_2^2
+
\frac{2}{n}
\E_\varepsilon
\left[
\varepsilon^\top X(\theta_*-\theta)
\right]
+
\frac{1}{n}
\E_\varepsilon\norm{\varepsilon}_2^2.
\]

The cross term is zero because
\[
\E[\varepsilon]=0.
\]

Also,
\[
\E\norm{\varepsilon}_2^2
=
\sum_{i=1}^n \E[\varepsilon_i^2]
=
n\sigma^2.
\]

Thus,
\[
R(\theta)
=
\frac{1}{n}\norm{X(\theta_*-\theta)}_2^2
+
\sigma^2.
\]

Using the definition of $\widehat{\Sigma}$,
\[
\frac{1}{n}\norm{X(\theta_*-\theta)}_2^2
=
(\theta-\theta_*)^\top\widehat{\Sigma}(\theta-\theta_*)
=
\norm{\theta-\theta_*}_{\widehat{\Sigma}}^2.
\]

Hence,
\[
R(\theta)
=
\norm{\theta-\theta_*}_{\widehat{\Sigma}}^2+\sigma^2.
\]

The minimum is attained at
\[
\theta=\theta_*,
\]
so
\[
R^*=\sigma^2.
\]
\end{proof}

\section{Bias and Variance of OLS}

\begin{lemma}
Suppose $X$ is fixed. Then
\[
\E[\widehat{\theta}_{\mathrm{OLS}}]=\theta_*,
\]
so the OLS estimator is unbiased. Moreover,
\[
\E\left[
(\widehat{\theta}_{\mathrm{OLS}}-\theta_*)
(\widehat{\theta}_{\mathrm{OLS}}-\theta_*)^\top
\right]
=
\sigma^2(X^\top X)^{-1}
=
\frac{\sigma^2}{n}\widehat{\Sigma}^{-1}.
\]
\end{lemma}

\begin{proof}
Using
\[
y=X\theta_*+\varepsilon,
\]
we get
\[
\widehat{\theta}_{\mathrm{OLS}}
=
(X^\top X)^{-1}X^\top y
=
(X^\top X)^{-1}X^\top(X\theta_*+\varepsilon).
\]

Thus,
\[
\widehat{\theta}_{\mathrm{OLS}}
=
\theta_*+(X^\top X)^{-1}X^\top\varepsilon.
\]

Taking expectations,
\[
\E[\widehat{\theta}_{\mathrm{OLS}}]
=
\theta_*+(X^\top X)^{-1}X^\top\E[\varepsilon]
=
\theta_*.
\]

Therefore the estimator is unbiased.

For the covariance,
\[
\widehat{\theta}_{\mathrm{OLS}}-\theta_*
=
(X^\top X)^{-1}X^\top\varepsilon.
\]

Hence,
\[
\E\left[
(\widehat{\theta}_{\mathrm{OLS}}-\theta_*)
(\widehat{\theta}_{\mathrm{OLS}}-\theta_*)^\top
\right]
=
(X^\top X)^{-1}X^\top
\E[\varepsilon\varepsilon^\top]
X(X^\top X)^{-1}.
\]

Since
\[
\E[\varepsilon\varepsilon^\top]=\sigma^2I_n,
\]
we obtain
\[
\E\left[
(\widehat{\theta}_{\mathrm{OLS}}-\theta_*)
(\widehat{\theta}_{\mathrm{OLS}}-\theta_*)^\top
\right]
=
\sigma^2(X^\top X)^{-1}X^\top X(X^\top X)^{-1}.
\]

Therefore,
\[
\E\left[
(\widehat{\theta}_{\mathrm{OLS}}-\theta_*)
(\widehat{\theta}_{\mathrm{OLS}}-\theta_*)^\top
\right]
=
\sigma^2(X^\top X)^{-1}.
\]

Since
\[
X^\top X=n\widehat{\Sigma},
\]
we also have
\[
\sigma^2(X^\top X)^{-1}
=
\frac{\sigma^2}{n}\widehat{\Sigma}^{-1}.
\]
\end{proof}

\section{Expected Excess Risk of OLS}

\begin{corollary}
Suppose $X$ is fixed. Then
\[
\E[R(\widehat{\theta}_{\mathrm{OLS}})]-R^*
=
\frac{d\sigma^2}{n}.
\]
\end{corollary}

\begin{proof}
By the fixed-design risk identity,
\[
R(\widehat{\theta}_{\mathrm{OLS}})-R^*
=
\norm{\widehat{\theta}_{\mathrm{OLS}}-\theta_*}_{\widehat{\Sigma}}^2.
\]

Taking expectations,
\[
\E[R(\widehat{\theta}_{\mathrm{OLS}})]-R^*
=
\E\left[
(\widehat{\theta}_{\mathrm{OLS}}-\theta_*)^\top
\widehat{\Sigma}
(\widehat{\theta}_{\mathrm{OLS}}-\theta_*)
\right].
\]

Using the trace identity
\[
a^\top Ba=\tr(Baa^\top),
\]
we get
\[
\E[R(\widehat{\theta}_{\mathrm{OLS}})]-R^*
=
\tr\left(
\widehat{\Sigma}
\E[
(\widehat{\theta}_{\mathrm{OLS}}-\theta_*)
(\widehat{\theta}_{\mathrm{OLS}}-\theta_*)^\top
]
\right).
\]

Substituting the covariance formula,
\[
\E[R(\widehat{\theta}_{\mathrm{OLS}})]-R^*
=
\tr\left(
\widehat{\Sigma}
\frac{\sigma^2}{n}\widehat{\Sigma}^{-1}
\right).
\]

Therefore,
\[
\E[R(\widehat{\theta}_{\mathrm{OLS}})]-R^*
=
\frac{\sigma^2}{n}\tr(I_d)
=
\frac{d\sigma^2}{n}.
\]
\end{proof}

\begin{remark}
In fixed-design ordinary least squares, the expected excess risk depends only on the noise variance $\sigma^2$, the dimension $d$, and the sample size $n$. It does not depend on the conditioning of $X^\top X$, as long as the prediction risk is measured in the $\widehat{\Sigma}$-norm.
\end{remark}

\section{Training Error of OLS}

A natural question is how close the empirical risk
\[
R_n(\widehat{\theta}_{\mathrm{OLS}})
\]
is to the population risk
\[
R(\widehat{\theta}_{\mathrm{OLS}}).
\]

\begin{lemma}
Suppose $X$ is fixed. Then
\[
\E[R_n(\widehat{\theta}_{\mathrm{OLS}})]
=
\sigma^2-\frac{d\sigma^2}{n}
=
\sigma^2\left(1-\frac{d}{n}\right).
\]
\end{lemma}

\begin{proof}
Let
\[
P_X=X(X^\top X)^{-1}X^\top
\]
be the orthogonal projector onto
\[
\operatorname{range}(X).
\]

Then
\[
X\widehat{\theta}_{\mathrm{OLS}}
=
P_Xy.
\]

Therefore,
\[
y-X\widehat{\theta}_{\mathrm{OLS}}
=
(I-P_X)y.
\]

Since
\[
y=X\theta_*+\varepsilon
\]
and
\[
X\theta_*\in\operatorname{range}(X),
\]
we have
\[
(I-P_X)X\theta_*=0.
\]

Thus,
\[
y-X\widehat{\theta}_{\mathrm{OLS}}
=
(I-P_X)\varepsilon.
\]

Hence,
\[
R_n(\widehat{\theta}_{\mathrm{OLS}})
=
\frac{1}{n}\norm{(I-P_X)\varepsilon}_2^2.
\]

Taking expectations,
\[
\E[R_n(\widehat{\theta}_{\mathrm{OLS}})]
=
\frac{1}{n}
\E[
\varepsilon^\top(I-P_X)^\top(I-P_X)\varepsilon
].
\]

Since $P_X$ is an orthogonal projector,
\[
(I-P_X)^\top(I-P_X)
=
(I-P_X)^2
=
I-P_X.
\]

Therefore,
\[
\E[R_n(\widehat{\theta}_{\mathrm{OLS}})]
=
\frac{1}{n}
\tr\left(
(I-P_X)\E[\varepsilon\varepsilon^\top]
\right).
\]

Using
\[
\E[\varepsilon\varepsilon^\top]=\sigma^2I_n,
\]
we obtain
\[
\E[R_n(\widehat{\theta}_{\mathrm{OLS}})]
=
\frac{\sigma^2}{n}
\tr(I-P_X).
\]

Because $P_X$ projects onto a $d$-dimensional subspace,
\[
\tr(P_X)=d.
\]

Therefore,
\[
\tr(I-P_X)=n-d.
\]

Thus,
\[
\E[R_n(\widehat{\theta}_{\mathrm{OLS}})]
=
\frac{\sigma^2}{n}(n-d)
=
\sigma^2-\frac{d\sigma^2}{n}.
\]
\end{proof}

\begin{remark}
The expected training error is smaller than the irreducible noise level
\[
\sigma^2
\]
by
\[
\frac{d\sigma^2}{n}.
\]

In contrast, the expected population risk of OLS is larger than
\[
\sigma^2
\]
by the same amount:
\[
\E[R(\widehat{\theta}_{\mathrm{OLS}})]
=
\sigma^2+\frac{d\sigma^2}{n}.
\]

Thus the expected gap between population risk and empirical risk at the fitted estimator is
\[
\E[R(\widehat{\theta}_{\mathrm{OLS}})]
-
\E[R_n(\widehat{\theta}_{\mathrm{OLS}})]
=
\frac{2d\sigma^2}{n}.
\]
\end{remark}

\chapter{Ridge Regression and Random Design}

\section{Motivation for Regularization}

The ordinary least squares estimator assumes that
\[
d\leq n
\]
and that
\[
X^\top X
\]
is invertible.

In high-dimensional settings,
\[
d\gg n,
\]
and the matrix
\[
X^\top X
\]
may fail to be invertible.

Even in the noiseless case,
\[
\varepsilon=0,
\]
there may be infinitely many interpolating solutions satisfying
\[
X\theta=y.
\]

To select among these solutions and impose structure, one introduces a regularizer.

Common choices include:
\begin{itemize}
\item the $\ell_1$ norm,
\item the squared $\ell_2$ norm.
\end{itemize}

\section{Ridge Regression}

\begin{definition}[Ridge Regression Estimator]
For
\[
\lambda>0,
\]
the ridge regression estimator is defined by
\[
\widehat{\theta}_\lambda
=
\argmin_{\theta\in\R^d}
\left[
\frac{1}{n}\norm{y-X\theta}_2^2
+
\lambda\norm{\theta}_2^2
\right].
\]
\end{definition}

The regularization term
\[
\lambda\norm{\theta}_2^2
\]
encourages solutions with small Euclidean norm.

\begin{lemma}
The ridge estimator exists, is unique, and satisfies
\[
\widehat{\theta}_\lambda
=
\frac{1}{n}
(\widehat{\Sigma}+\lambda I)^{-1}X^\top y,
\]
where
\[
\widehat{\Sigma}
=
\frac{1}{n}X^\top X.
\]
\end{lemma}

\begin{proof}
Define the ridge objective
\[
L_\lambda(\theta)
=
\frac{1}{n}\norm{y-X\theta}_2^2
+
\lambda\norm{\theta}_2^2.
\]

Since
\[
\lambda>0,
\]
the objective is strongly convex, so the minimizer exists and is unique.

Differentiating,
\[
\nabla L_\lambda(\theta)
=
\frac{2}{n}X^\top(X\theta-y)+2\lambda\theta.
\]

Setting the gradient equal to zero,
\[
0
=
\frac{2}{n}X^\top(X\theta-y)+2\lambda\theta.
\]

Equivalently,
\[
\left(
\frac{1}{n}X^\top X+\lambda I
\right)\theta
=
\frac{1}{n}X^\top y.
\]

Using
\[
\widehat{\Sigma}
=
\frac{1}{n}X^\top X,
\]
we obtain
\[
(\widehat{\Sigma}+\lambda I)\theta
=
\frac{1}{n}X^\top y.
\]

Since
\[
\widehat{\Sigma}+\lambda I
\]
is strictly positive definite, it is invertible, and therefore
\[
\widehat{\theta}_\lambda
=
\frac{1}{n}
(\widehat{\Sigma}+\lambda I)^{-1}X^\top y.
\]
\end{proof}

\begin{remark}
Unlike ordinary least squares, ridge regression remains well-defined even when
\[
X^\top X
\]
is singular.
\end{remark}

\begin{remark}
If
\[
\lambda\to0,
\]
then ridge regression approaches the minimum norm least squares solution:
\[
\widehat{\theta}_\lambda
\to
(X^\top X)^+X^\top y,
\]
where
\[
(\cdot)^+
\]
denotes the Moore--Penrose pseudoinverse.
\end{remark}

\section{Bias-Variance Tradeoff in Ridge Regression}

Recall the fixed-design model
\[
y=X\theta_*+\varepsilon,
\qquad
\E[\varepsilon]=0,
\qquad
\E[\varepsilon\varepsilon^\top]=\sigma^2I_n.
\]

We study the bias and variance of ridge regression with respect to the norm induced by
\[
\widehat{\Sigma}.
\]

\begin{lemma}
For any symmetric matrix
\[
A,
\]
\[
A(A+\lambda I)^{-1}
=
(A+\lambda I)^{-1}A.
\]
\end{lemma}

\begin{proof}
Since $A$ is symmetric, it admits an eigendecomposition
\[
A=U\Lambda U^\top.
\]

Then
\[
A+\lambda I
=
U(\Lambda+\lambda I)U^\top,
\]
and therefore
\[
(A+\lambda I)^{-1}
=
U(\Lambda+\lambda I)^{-1}U^\top.
\]

Hence,
\[
A(A+\lambda I)^{-1}
=
U\Lambda(\Lambda+\lambda I)^{-1}U^\top,
\]
while
\[
(A+\lambda I)^{-1}A
=
U(\Lambda+\lambda I)^{-1}\Lambda U^\top.
\]

Since diagonal matrices commute,
\[
\Lambda(\Lambda+\lambda I)^{-1}
=
(\Lambda+\lambda I)^{-1}\Lambda.
\]

Thus,
\[
A(A+\lambda I)^{-1}
=
(A+\lambda I)^{-1}A.
\]
\end{proof}

\begin{theorem}[Bias and Variance of Ridge Regression]
The ridge estimator satisfies
\[
\operatorname{Bias}_{\widehat{\Sigma}}(\widehat{\theta}_\lambda)
=
\lambda^2
\theta_*^\top
(\widehat{\Sigma}+\lambda I)^{-2}
\widehat{\Sigma}
\theta_*,
\]
and
\[
\operatorname{Var}_{\widehat{\Sigma}}(\widehat{\theta}_\lambda)
=
\frac{\sigma^2}{n}
\tr\left(
\widehat{\Sigma}^2
(\widehat{\Sigma}+\lambda I)^{-2}
\right).
\]
\end{theorem}

\begin{proof}
Using
\[
y=X\theta_*+\varepsilon,
\]
we obtain
\[
\widehat{\theta}_\lambda
=
\frac{1}{n}
(\widehat{\Sigma}+\lambda I)^{-1}
X^\top(X\theta_*+\varepsilon).
\]

Since
\[
\frac{1}{n}X^\top X
=
\widehat{\Sigma},
\]
we have
\[
\widehat{\theta}_\lambda
=
(\widehat{\Sigma}+\lambda I)^{-1}
\widehat{\Sigma}\theta_*
+
\frac{1}{n}
(\widehat{\Sigma}+\lambda I)^{-1}
X^\top\varepsilon.
\]

Taking expectations,
\[
\E[\widehat{\theta}_\lambda]
=
(\widehat{\Sigma}+\lambda I)^{-1}
\widehat{\Sigma}\theta_*.
\]

Using
\[
\widehat{\Sigma}
=
(\widehat{\Sigma}+\lambda I)-\lambda I,
\]
we rewrite this as
\[
\E[\widehat{\theta}_\lambda]
=
\theta_*
-
\lambda
(\widehat{\Sigma}+\lambda I)^{-1}\theta_*.
\]

Therefore,
\[
\E[\widehat{\theta}_\lambda]-\theta_*
=
-
\lambda
(\widehat{\Sigma}+\lambda I)^{-1}\theta_*.
\]

Hence,
\[
\operatorname{Bias}_{\widehat{\Sigma}}(\widehat{\theta}_\lambda)
=
\norm{
\E[\widehat{\theta}_\lambda]-\theta_*
}_{\widehat{\Sigma}}^2.
\]

Expanding,
\[
\operatorname{Bias}_{\widehat{\Sigma}}(\widehat{\theta}_\lambda)
=
\lambda^2
\theta_*^\top
(\widehat{\Sigma}+\lambda I)^{-1}
\widehat{\Sigma}
(\widehat{\Sigma}+\lambda I)^{-1}
\theta_*.
\]

Since
\[
\widehat{\Sigma}
\]
commutes with
\[
(\widehat{\Sigma}+\lambda I)^{-1},
\]
this becomes
\[
\operatorname{Bias}_{\widehat{\Sigma}}(\widehat{\theta}_\lambda)
=
\lambda^2
\theta_*^\top
(\widehat{\Sigma}+\lambda I)^{-2}
\widehat{\Sigma}
\theta_*.
\]

Next,
\[
\widehat{\theta}_\lambda-\E[\widehat{\theta}_\lambda]
=
\frac{1}{n}
(\widehat{\Sigma}+\lambda I)^{-1}
X^\top\varepsilon.
\]

Thus,
\[
\operatorname{Var}_{\widehat{\Sigma}}(\widehat{\theta}_\lambda)
=
\E
\norm{
\widehat{\theta}_\lambda-\E[\widehat{\theta}_\lambda]
}_{\widehat{\Sigma}}^2.
\]

Expanding,
\[
=
\frac{1}{n^2}
\E\left[
\varepsilon^\top
X
(\widehat{\Sigma}+\lambda I)^{-1}
\widehat{\Sigma}
(\widehat{\Sigma}+\lambda I)^{-1}
X^\top
\varepsilon
\right].
\]

Using the trace identity,
\[
=
\frac{1}{n^2}
\E
\tr\left(
\varepsilon\varepsilon^\top
X
(\widehat{\Sigma}+\lambda I)^{-1}
\widehat{\Sigma}
(\widehat{\Sigma}+\lambda I)^{-1}
X^\top
\right).
\]

Since
\[
\E[\varepsilon\varepsilon^\top]
=
\sigma^2I_n,
\]
we obtain
\[
=
\frac{\sigma^2}{n^2}
\tr\left(
X
(\widehat{\Sigma}+\lambda I)^{-1}
\widehat{\Sigma}
(\widehat{\Sigma}+\lambda I)^{-1}
X^\top
\right).
\]

Using cyclic invariance of the trace,
\[
=
\frac{\sigma^2}{n^2}
\tr\left(
X^\top X
(\widehat{\Sigma}+\lambda I)^{-1}
\widehat{\Sigma}
(\widehat{\Sigma}+\lambda I)^{-1}
\right).
\]

Since
\[
X^\top X=n\widehat{\Sigma},
\]
we get
\[
=
\frac{\sigma^2}{n}
\tr\left(
\widehat{\Sigma}^2
(\widehat{\Sigma}+\lambda I)^{-2}
\right).
\]
\end{proof}

\section{Effective Dimension}

Suppose the eigenvalues of
\[
\widehat{\Sigma}
\]
are
\[
\lambda_1,\dots,\lambda_d.
\]

Then
\[
\tr\left(
\widehat{\Sigma}^2
(\widehat{\Sigma}+\lambda I)^{-2}
\right)
=
\sum_{i=1}^d
\frac{\lambda_i^2}{(\lambda_i+\lambda)^2}.
\]

Observe that
\[
0
\leq
\frac{\lambda_i^2}{(\lambda_i+\lambda)^2}
\leq1.
\]

Therefore,
\[
\tr\left(
\widehat{\Sigma}^2
(\widehat{\Sigma}+\lambda I)^{-2}
\right)
\leq d.
\]

This quantity acts as a soft notion of dimension or degrees of freedom.

Only eigenvalues satisfying
\[
\lambda_i\gg\lambda
\]
contribute substantially.

\section{Excess Risk of Ridge Regression}

\begin{corollary}
The expected excess risk of ridge regression is
\[
\E[R(\widehat{\theta}_\lambda)]-R^*
=
\lambda^2
\theta_*^\top
(\widehat{\Sigma}+\lambda I)^{-2}
\widehat{\Sigma}
\theta_*
+
\frac{\sigma^2}{n}
\tr\left(
\widehat{\Sigma}^2
(\widehat{\Sigma}+\lambda I)^{-2}
\right).
\]
\end{corollary}

\begin{remark}
Unlike ordinary least squares, the ridge estimator is biased whenever
\[
\lambda>0.
\]

However, the variance may be dramatically reduced.
\end{remark}

\begin{remark}
The ridge excess risk bound contains no explicit dependence on the ambient dimension
\[
d.
\]

This is crucial for learning in high-dimensional or infinite-dimensional settings.
\end{remark}

\section{Choosing the Regularization Parameter}

We now derive a useful upper bound.

\begin{theorem}
For every
\[
\lambda>0,
\]
\[
\E[R(\widehat{\theta}_\lambda)]-R^*
\leq
\frac{\lambda}{2}\norm{\theta_*}_2^2
+
\frac{\sigma^2\tr(\widehat{\Sigma})}{2\lambda n}.
\]
\end{theorem}

\begin{proof}
We first bound the bias term.

Observe that
\[
\lambda^2
\theta_*^\top
(\widehat{\Sigma}+\lambda I)^{-2}
\widehat{\Sigma}
\theta_*
\leq
\lambda
\norm{
(\widehat{\Sigma}+\lambda I)^{-2}
\lambda\widehat{\Sigma}
}_{\mathrm{op}}
\norm{\theta_*}_2^2.
\]

Let
\[
s_i
\]
denote the eigenvalues of
\[
\widehat{\Sigma}.
\]

Then the eigenvalues of
\[
(\widehat{\Sigma}+\lambda I)^{-2}
\lambda\widehat{\Sigma}
\]
are
\[
\frac{\lambda s_i}{(s_i+\lambda)^2}.
\]

Using
\[
2\lambda s_i
\leq
(s_i+\lambda)^2,
\]
we obtain
\[
\frac{\lambda s_i}{(s_i+\lambda)^2}
\leq
\frac{1}{2}.
\]

Hence,
\[
\lambda^2
\theta_*^\top
(\widehat{\Sigma}+\lambda I)^{-2}
\widehat{\Sigma}
\theta_*
\leq
\frac{\lambda}{2}\norm{\theta_*}_2^2.
\]

Next, for the variance term,
\[
\frac{\sigma^2}{n}
\tr\left(
\widehat{\Sigma}^2
(\widehat{\Sigma}+\lambda I)^{-2}
\right)
=
\frac{\sigma^2}{\lambda n}
\tr\left(
\widehat{\Sigma}
\lambda\widehat{\Sigma}
(\widehat{\Sigma}+\lambda I)^{-2}
\right).
\]

Using
\[
\tr(AB)
\leq
\tr(A)\lambda_{\max}(B)
\]
for positive semidefinite matrices,
\[
\leq
\frac{\sigma^2}{\lambda n}
\tr(\widehat{\Sigma})
\lambda_{\max}
\left(
\lambda\widehat{\Sigma}
(\widehat{\Sigma}+\lambda I)^{-2}
\right).
\]

Again,
\[
\lambda_{\max}
\left(
\lambda\widehat{\Sigma}
(\widehat{\Sigma}+\lambda I)^{-2}
\right)
\leq
\frac{1}{2}.
\]

Therefore,
\[
\frac{\sigma^2}{n}
\tr\left(
\widehat{\Sigma}^2
(\widehat{\Sigma}+\lambda I)^{-2}
\right)
\leq
\frac{\sigma^2\tr(\widehat{\Sigma})}{2\lambda n}.
\]

Combining the two bounds yields
\[
\E[R(\widehat{\theta}_\lambda)]-R^*
\leq
\frac{\lambda}{2}\norm{\theta_*}_2^2
+
\frac{\sigma^2\tr(\widehat{\Sigma})}{2\lambda n}.
\]
\end{proof}

\begin{corollary}
The bound is minimized by
\[
\lambda
=
\frac{\sigma}{\norm{\theta_*}_2}
\sqrt{
\frac{\tr(\widehat{\Sigma})}{n}
}.
\]

For this choice,
\[
\E[R(\widehat{\theta}_\lambda)]-R^*
\leq
\sigma\norm{\theta_*}_2
\sqrt{
\frac{\tr(\widehat{\Sigma})}{n}
}.
\]
\end{corollary}

\begin{proof}
The upper bound has the form
\[
a\lambda+\frac{b}{\lambda},
\]
where
\[
a,b>0.
\]

Such an expression is minimized at
\[
\lambda=\sqrt{\frac{b}{a}}.
\]

Substituting
\[
a=\frac{1}{2}\norm{\theta_*}_2^2,
\qquad
b=\frac{\sigma^2\tr(\widehat{\Sigma})}{2n},
\]
gives the result.
\end{proof}

\begin{remark}
Choosing the optimal regularization parameter in practice is difficult. Cross-validation is often used.
\end{remark}

\section{Random Design}

So far we treated
\[
X
\]
as deterministic.

We now assume that
\[
(X_i,y_i)
\]
are iid samples satisfying
\[
y_i=\theta_*^\top X_i+\varepsilon_i,
\]
with
\[
\E[\varepsilon_i]=0,
\qquad
\E[\varepsilon_i^2]=\sigma^2.
\]

Define the population covariance matrix
\[
\Sigma
=
\E[XX^\top].
\]

\begin{lemma}
Under the random design model,
\[
R(\theta)-R^*
=
\norm{\theta-\theta_*}_{\Sigma}^2,
\]
where
\[
\norm{u}_{\Sigma}^2
=
u^\top\Sigma u.
\]
\end{lemma}

\begin{proof}
We compute
\[
R(\theta)
=
\E[(y-X^\top\theta)^2].
\]

Using
\[
y=X^\top\theta_*+\varepsilon,
\]
we obtain
\[
R(\theta)
=
\E[(X^\top(\theta_*-\theta)+\varepsilon)^2].
\]

Expanding,
\[
=
\E[(\theta_*-\theta)^\top XX^\top(\theta_*-\theta)]
+
2\E[\varepsilon X^\top(\theta_*-\theta)]
+
\E[\varepsilon^2].
\]

Since
\[
\E[\varepsilon|X]=0,
\]
the cross term vanishes.

Thus,
\[
R(\theta)
=
(\theta_*-\theta)^\top
\E[XX^\top]
(\theta_*-\theta)
+
\sigma^2.
\]

Therefore,
\[
R(\theta)-R^*
=
\norm{\theta-\theta_*}_{\Sigma}^2.
\]
\end{proof}

\section{OLS Under Random Design}

\begin{theorem}
Under the random design model, assuming
\[
\widehat{\Sigma}
=
\frac{1}{n}X^\top X
\]
is invertible,
\[
\E[R(\widehat{\theta}_{\mathrm{OLS}})]-R^*
=
\frac{\sigma^2}{n}
\E\tr(\Sigma\widehat{\Sigma}^{-1}).
\]
\end{theorem}

\begin{proof}
Recall
\[
\widehat{\theta}_{\mathrm{OLS}}
=
\theta_*
+
\frac{1}{n}
\widehat{\Sigma}^{-1}X^\top\varepsilon.
\]

Hence,
\[
\widehat{\theta}_{\mathrm{OLS}}-\theta_*
=
\frac{1}{n}
\widehat{\Sigma}^{-1}X^\top\varepsilon.
\]

Using the random-design risk identity,
\[
R(\widehat{\theta}_{\mathrm{OLS}})-R^*
=
\norm{
\widehat{\theta}_{\mathrm{OLS}}-\theta_*
}_{\Sigma}^2.
\]

Therefore,
\[
=
\frac{1}{n^2}
\varepsilon^\top
X
\widehat{\Sigma}^{-1}
\Sigma
\widehat{\Sigma}^{-1}
X^\top
\varepsilon.
\]

Taking expectations and using the trace trick,
\[
\E[R(\widehat{\theta}_{\mathrm{OLS}})]-R^*
=
\frac{1}{n^2}
\E\tr\left(
\varepsilon\varepsilon^\top
X
\widehat{\Sigma}^{-1}
\Sigma
\widehat{\Sigma}^{-1}
X^\top
\right).
\]

Conditioning on $X$,
\[
\E[\varepsilon\varepsilon^\top|X]
=
\sigma^2I_n.
\]

Thus,
\[
=
\frac{\sigma^2}{n^2}
\E\tr\left(
X
\widehat{\Sigma}^{-1}
\Sigma
\widehat{\Sigma}^{-1}
X^\top
\right).
\]

Using cyclic invariance,
\[
=
\frac{\sigma^2}{n^2}
\E\tr\left(
X^\top X
\widehat{\Sigma}^{-1}
\Sigma
\widehat{\Sigma}^{-1}
\right).
\]

Since
\[
X^\top X=n\widehat{\Sigma},
\]
we obtain
\[
=
\frac{\sigma^2}{n}
\E\tr(\Sigma\widehat{\Sigma}^{-1}).
\]
\end{proof}

\begin{remark}
Unlike the fixed-design setting,
\[
\Sigma\widehat{\Sigma}^{-1}\neq I
\]
in general.

The excess risk depends on how well the empirical covariance matrix approximates the population covariance matrix.
\end{remark}

\begin{example}
If
\[
X\sim \Normal(0,\Sigma),
\]
then one can show
\[
\E\tr(\Sigma\widehat{\Sigma}^{-1})
=
\frac{d}{1-\frac{d+1}{n}}.
\]

This follows from properties of the inverse Wishart distribution.
\end{example}

\section{Beyond Linear Features}

Everything developed above extends immediately to nonlinear feature maps.

Suppose
\[
\psi:\R^k\to\R^d
\]
is a feature map and the model becomes
\[
y_i
=
\theta_*^\top\psi(x_i)+\varepsilon_i.
\]

Define the feature matrix
\[
\Phi
=
\begin{pmatrix}
\psi(x_1)^\top\\
\vdots\\
\psi(x_n)^\top
\end{pmatrix}.
\]

Then all previous derivations hold after replacing
\[
X
\]
with
\[
\Phi,
\]
and replacing
\[
\widehat{\Sigma}
\]
with
\[
\frac{1}{n}\Phi^\top\Phi.
\]

\begin{remark}
This perspective is fundamental in kernel methods, where the feature dimension may even be infinite.
\end{remark}

\chapter{Rademacher Complexity}

\section{Uniform Deviations}

We have been studying the relationship between the empirical optimization problem
\[
\min_{\theta\in\Theta}
\frac{1}{n}\sum_{i=1}^n \ell(\theta,Z_i)
\]
and the population optimization problem
\[
\min_{\theta\in\Theta}
\E[\ell(\theta,Z)].
\]

To understand when empirical risk minimization approximates population risk minimization, we need to control uniform deviations.

Let
\[
X_1,\dots,X_n\iid P
\]
take values in a set
\[
\mathcal{X}.
\]

Let
\[
\mathcal{F}
\]
be a class of functions
\[
f:\mathcal{X}\to\R.
\]

Define the empirical measure
\[
P_nf
=
\frac{1}{n}\sum_{i=1}^n f(X_i)
\]
and the population expectation
\[
Pf
=
\E[f(X)].
\]

The central quantity is
\[
\norm{P_n-P}_{\mathcal{F}}
=
\sup_{f\in\mathcal{F}}
\left|
\frac{1}{n}\sum_{i=1}^n f(X_i)
-
\E[f(X)]
\right|.
\]

\section{Rademacher Complexity of a Set}

\begin{definition}[Rademacher Random Variable]
A Rademacher random variable is a random variable
\[
\varepsilon
\]
such that
\[
\Prob(\varepsilon=1)=\Prob(\varepsilon=-1)=\frac{1}{2}.
\]
\end{definition}

\begin{definition}[Rademacher Complexity of a Set]
Let
\[
A\subseteq\R^n.
\]

The Rademacher complexity of $A$ is
\[
\mathcal{R}(A)
=
\E_{\varepsilon}
\left[
\sup_{a\in A}
\left|
\sum_{i=1}^n \varepsilon_i a_i
\right|
\right],
\]
where
\[
\varepsilon=(\varepsilon_1,\dots,\varepsilon_n)
\]
has iid Rademacher entries.
\end{definition}

\section{Geometric Interpretation}

\begin{definition}[Support Function]
The support function of a set
\[
A\subseteq\R^n
\]
is
\[
\sigma_A(v)
=
\sup_{a\in A}
\langle v,a\rangle.
\]
\end{definition}

For a unit vector
\[
v\in S^{n-1},
\]
the width of $A$ in direction $v$ is
\[
\sigma_A(v)+\sigma_A(-v).
\]

Indeed,
\[
\sigma_A(v)+\sigma_A(-v)
=
\sup_{a\in A}\langle v,a\rangle
+
\sup_{a\in A}\langle -v,a\rangle.
\]

Since
\[
\sup_{a\in A}\langle -v,a\rangle
=
-\inf_{a\in A}\langle v,a\rangle,
\]
we get
\[
\sigma_A(v)+\sigma_A(-v)
=
\sup_{a\in A}\langle v,a\rangle
-
\inf_{a\in A}\langle v,a\rangle.
\]

Thus Rademacher complexity measures the average width of $A$ in random coordinate-sign directions.

\begin{remark}
The quantity
\[
\mathcal{R}(A)
\]
is comparable to
\[
\E_{\varepsilon}[\sigma_A(\varepsilon)].
\]
The absolute value in the definition symmetrizes the set by considering both directions.
\end{remark}

\section{Rademacher Complexity of a Function Class}

Given samples
\[
X_1,\dots,X_n,
\]
define the evaluation set
\[
\mathcal{F}(X)
=
\left\{
(f(X_1),\dots,f(X_n)):f\in\mathcal{F}
\right\}
\subseteq\R^n.
\]

\begin{definition}[Rademacher Complexity of a Function Class]
The Rademacher complexity of
\[
\mathcal{F}
\]
at sample size $n$ is
\[
\mathcal{R}_n(\mathcal{F})
=
\E_X
\left[
\mathcal{R}
\left(
\frac{1}{n}\mathcal{F}(X)
\right)
\right].
\]

Equivalently,
\[
\mathcal{R}_n(\mathcal{F})
=
\E_{X,\varepsilon}
\left[
\sup_{f\in\mathcal{F}}
\left|
\frac{1}{n}
\sum_{i=1}^n
\varepsilon_i f(X_i)
\right|
\right].
\]
\end{definition}

\section{Symmetrization}

Rademacher complexity controls the expected uniform deviation between empirical and population averages.

\begin{theorem}[Symmetrization]
For every
\[
n\geq1,
\]
\[
\E\norm{P_n-P}_{\mathcal{F}}
\leq
2\mathcal{R}_n(\mathcal{F}).
\]
\end{theorem}

\begin{proof}
Let
\[
Y_1,\dots,Y_n\iid P
\]
be an independent ghost sample.

Then
\[
\E[f(X)]
=
\E_Y[f(Y_i)]
\]
for every $i$.

Therefore,
\[
\E\norm{P_n-P}_{\mathcal{F}}
=
\E_X
\left[
\sup_{f\in\mathcal{F}}
\left|
\frac{1}{n}\sum_{i=1}^n f(X_i)
-
\E_Y[f(Y_i)]
\right|
\right].
\]

By Jensen's inequality,
\[
\E\norm{P_n-P}_{\mathcal{F}}
\leq
\E_{X,Y}
\left[
\sup_{f\in\mathcal{F}}
\left|
\frac{1}{n}\sum_{i=1}^n
\bigl(f(X_i)-f(Y_i)\bigr)
\right|
\right].
\]

Now introduce iid Rademacher variables
\[
\varepsilon_1,\dots,\varepsilon_n.
\]

Since
\[
f(X_i)-f(Y_i)
\]
has the same distribution as
\[
\varepsilon_i(f(X_i)-f(Y_i)),
\]
we get
\[
\E\norm{P_n-P}_{\mathcal{F}}
\leq
\E_{X,Y,\varepsilon}
\left[
\sup_{f\in\mathcal{F}}
\left|
\frac{1}{n}\sum_{i=1}^n
\varepsilon_i
\bigl(f(X_i)-f(Y_i)\bigr)
\right|
\right].
\]

Using the triangle inequality,
\[
\left|
\frac{1}{n}\sum_{i=1}^n
\varepsilon_i
\bigl(f(X_i)-f(Y_i)\bigr)
\right|
\leq
\left|
\frac{1}{n}\sum_{i=1}^n
\varepsilon_i f(X_i)
\right|
+
\left|
\frac{1}{n}\sum_{i=1}^n
\varepsilon_i f(Y_i)
\right|.
\]

Taking suprema and expectations,
\[
\E\norm{P_n-P}_{\mathcal{F}}
\leq
\E_{X,\varepsilon}
\left[
\sup_{f\in\mathcal{F}}
\left|
\frac{1}{n}\sum_{i=1}^n
\varepsilon_i f(X_i)
\right|
\right]
+
\E_{Y,\varepsilon}
\left[
\sup_{f\in\mathcal{F}}
\left|
\frac{1}{n}\sum_{i=1}^n
\varepsilon_i f(Y_i)
\right|
\right].
\]

The two terms are identical, so
\[
\E\norm{P_n-P}_{\mathcal{F}}
\leq
2\mathcal{R}_n(\mathcal{F}).
\]
\end{proof}

\section{High-Probability Uniform Deviation Bound}

If the function class is uniformly bounded, then McDiarmid's inequality upgrades the expectation bound to a high-probability bound.

\begin{theorem}
Suppose
\[
\mathcal{F}
\]
satisfies
\[
\norm{f}_{\infty}
=
\sup_{x\in\mathcal{X}}|f(x)|
\leq b
\]
for every
\[
f\in\mathcal{F}.
\]

Then, for every
\[
t\geq0,
\]
\[
\Prob\left(
\norm{P_n-P}_{\mathcal{F}}
\leq
2\mathcal{R}_n(\mathcal{F})+t
\right)
\geq
1-\exp\left(
-\frac{nt^2}{2b^2}
\right).
\]
\end{theorem}

\begin{proof}
Define
\[
\bar{f}(x)
=
f(x)-\E[f(X)].
\]

Then
\[
\norm{P_n-P}_{\mathcal{F}}
=
\sup_{f\in\mathcal{F}}
\left|
\frac{1}{n}\sum_{i=1}^n \bar{f}(X_i)
\right|.
\]

Let
\[
G(X_1,\dots,X_n)
=
\norm{P_n-P}_{\mathcal{F}}.
\]

We show that $G$ satisfies the bounded differences property.

Let
\[
X'
=
(X_1,\dots,X_i',\dots,X_n)
\]
differ from
\[
X
\]
only in the $i$-th coordinate.

For any fixed
\[
f\in\mathcal{F},
\]
\[
\left|
\frac{1}{n}\sum_{j=1}^n\bar{f}(X_j)
-
\frac{1}{n}\sum_{j=1}^n\bar{f}(X_j')
\right|
=
\frac{1}{n}
|\bar{f}(X_i)-\bar{f}(X_i')|.
\]

Since
\[
|f(x)|\leq b
\]
and
\[
|\E f(X)|\leq b,
\]
we have
\[
|\bar{f}(x)|\leq 2b.
\]

However, using the uncentered difference,
\[
|\bar{f}(X_i)-\bar{f}(X_i')|
=
|f(X_i)-f(X_i')|
\leq
2b.
\]

Thus changing one coordinate changes each empirical average by at most
\[
\frac{2b}{n}.
\]

Taking the supremum preserves this Lipschitz bound, so
\[
G
\]
has bounded differences constants
\[
c_i=\frac{2b}{n}.
\]

McDiarmid's inequality gives
\[
\Prob\left(
G-\E[G]\geq t
\right)
\leq
\exp\left(
-\frac{2t^2}{\sum_{i=1}^n c_i^2}
\right).
\]

Since
\[
\sum_{i=1}^n c_i^2
=
n\left(\frac{2b}{n}\right)^2
=
\frac{4b^2}{n},
\]
we obtain
\[
\Prob\left(
G-\E[G]\geq t
\right)
\leq
\exp\left(
-\frac{nt^2}{2b^2}
\right).
\]

Finally, by symmetrization,
\[
\E[G]
\leq
2\mathcal{R}_n(\mathcal{F}).
\]

Thus,
\[
\Prob\left(
\norm{P_n-P}_{\mathcal{F}}
\leq
2\mathcal{R}_n(\mathcal{F})+t
\right)
\geq
1-\exp\left(
-\frac{nt^2}{2b^2}
\right).
\]
\end{proof}

\section{Classes with Polynomial Discrimination}

We now develop tools for bounding Rademacher complexities.

\begin{definition}[Polynomial Discrimination]
A function class
\[
\mathcal{F}
\]
has polynomial discrimination of order
\[
\nu\geq1
\]
if, for every
\[
n\geq1
\]
and every
\[
x_1,\dots,x_n\in\mathcal{X},
\]
\[
\#\mathcal{F}(x_1,\dots,x_n)
\leq
(n+1)^\nu,
\]
where
\[
\mathcal{F}(x_1,\dots,x_n)
=
\left\{
(f(x_1),\dots,f(x_n)):f\in\mathcal{F}
\right\}.
\]
\end{definition}

\begin{remark}
Even if every
\[
f\in\mathcal{F}
\]
is binary-valued, the number of possible labelings can be as large as
\[
2^n.
\]

Polynomial discrimination assumes that the number of distinct labelings grows only polynomially in
\[
n.
\]
\end{remark}

\section{A Finite-Class Rademacher Bound}

\begin{lemma}
Let
\[
A\subseteq\R^n
\]
be finite. Then
\[
\mathcal{R}(A)
\leq
\max_{a\in A}\norm{a}_2
\sqrt{2\log(\#A)}.
\]
\end{lemma}

\begin{proof}
For a fixed
\[
a\in A,
\]
the random variable
\[
\langle a,\varepsilon\rangle
=
\sum_{i=1}^n a_i\varepsilon_i
\]
is sub-Gaussian with variance proxy
\[
\norm{a}_2^2.
\]

Indeed, by independence and the Rademacher moment generating function,
\[
\E\exp(\lambda\langle a,\varepsilon\rangle)
=
\prod_{i=1}^n
\E\exp(\lambda a_i\varepsilon_i)
\leq
\prod_{i=1}^n
\exp\left(
\frac{\lambda^2a_i^2}{2}
\right)
=
\exp\left(
\frac{\lambda^2\norm{a}_2^2}{2}
\right).
\]

Let
\[
M=\max_{a\in A}\norm{a}_2.
\]

Then each
\[
\langle a,\varepsilon\rangle
\]
is $M$-sub-Gaussian.

The standard maximum inequality for finitely many sub-Gaussian random variables gives
\[
\E\max_{a\in A}
|\langle a,\varepsilon\rangle|
\leq
M\sqrt{2\log(\#A)}
\]
up to a universal constant. Absorbing constants into the displayed bound yields
\[
\mathcal{R}(A)
\leq
\max_{a\in A}\norm{a}_2
\sqrt{2\log(\#A)}.
\]
\end{proof}

\section{Rademacher Complexity under Polynomial Discrimination}

\begin{proposition}
Suppose
\[
\mathcal{F}
\]
has polynomial discrimination of order
\[
\nu.
\]

Then
\[
\mathcal{R}_n(\mathcal{F})
\leq
\E
\left[
\sup_{f\in\mathcal{F}}
\sqrt{
\frac{1}{n}
\sum_{i=1}^n f(X_i)^2
}
\right]
\sqrt{
\frac{2\nu\log(n+1)}{n}
}.
\]
\end{proposition}

\begin{proof}
Condition on
\[
X_1,\dots,X_n.
\]

Apply the finite-class Rademacher bound to
\[
A=\frac{1}{n}\mathcal{F}(X).
\]

Then
\[
\mathcal{R}
\left(
\frac{1}{n}\mathcal{F}(X)
\right)
\leq
\max_{f\in\mathcal{F}}
\left\|
\frac{1}{n}
(f(X_1),\dots,f(X_n))
\right\|_2
\sqrt{
2\log(\#\mathcal{F}(X))
}.
\]

The norm term is
\[
\left\|
\frac{1}{n}
(f(X_1),\dots,f(X_n))
\right\|_2
=
\frac{1}{n}
\sqrt{
\sum_{i=1}^n f(X_i)^2
}
=
\sqrt{
\frac{1}{n}
\sum_{i=1}^n f(X_i)^2
}
\frac{1}{\sqrt{n}}.
\]

Since $\mathcal{F}$ has polynomial discrimination,
\[
\#\mathcal{F}(X)
\leq
(n+1)^\nu.
\]

Therefore,
\[
\log(\#\mathcal{F}(X))
\leq
\nu\log(n+1).
\]

Combining,
\[
\mathcal{R}
\left(
\frac{1}{n}\mathcal{F}(X)
\right)
\leq
\sup_{f\in\mathcal{F}}
\sqrt{
\frac{1}{n}
\sum_{i=1}^n f(X_i)^2
}
\sqrt{
\frac{2\nu\log(n+1)}{n}
}.
\]

Taking expectation over
\[
X
\]
gives the result.
\end{proof}

\begin{corollary}
If
\[
\mathcal{F}
\]
has polynomial discrimination of order $\nu$ and is uniformly bounded by
\[
|f(x)|\leq b
\]
for all
\[
f\in\mathcal{F},
\]
then
\[
\mathcal{R}_n(\mathcal{F})
\leq
b
\sqrt{
\frac{2\nu\log(n+1)}{n}
}.
\]
\end{corollary}

\section{Example: Half-Intervals}

Consider the class
\[
\mathcal{F}
=
\{f_t:t\in\R\},
\]
where
\[
f_t(x)
=
\begin{cases}
1, & x\leq t,\\
0, & x>t.
\end{cases}
\]

Let
\[
x_1,\dots,x_n\in\R
\]
and let
\[
x_{(1)}\leq \cdots \leq x_{(n)}
\]
be the order statistics.

As $t$ varies, the vector
\[
(f_t(x_1),\dots,f_t(x_n))
\]
changes only when $t$ crosses one of the observed sample values.

Thus, after ordering the points, the possible labeling patterns are
\[
(0,\dots,0),
\]
\[
(1,0,\dots,0),
\]
\[
(1,1,0,\dots,0),
\]
and so on, up to
\[
(1,\dots,1).
\]

Hence
\[
\#\mathcal{F}(x_1,\dots,x_n)
\leq
n+1.
\]

Therefore, this class has polynomial discrimination of order
\[
\nu=1.
\]

\section{Glivenko--Cantelli Bound}

Let
\[
F(t)=\Prob(X\leq t)
\]
be the cumulative distribution function of
\[
X\sim P.
\]

Let
\[
\widehat{F}_n(t)
=
\frac{1}{n}
\sum_{i=1}^n
\ind_{\{X_i\leq t\}}
\]
be the empirical cumulative distribution function.

The class of half-interval indicators is uniformly bounded by
\[
b=1
\]
and has polynomial discrimination order
\[
\nu=1.
\]

Therefore,
\[
\mathcal{R}_n(\mathcal{F})
\leq
\sqrt{
\frac{2\log(n+1)}{n}
}.
\]

Using the high-probability uniform deviation bound,
\[
\Prob\left(
\norm{\widehat{F}_n-F}_{\infty}
\leq
2\sqrt{
\frac{2\log(n+1)}{n}
}
+t
\right)
\geq
1-\exp\left(
-\frac{nt^2}{2}
\right).
\]

Equivalently, for all
\[
\delta>0,
\]
with probability at least
\[
1-e^{-n\delta^2/2},
\]
\[
\norm{\widehat{F}_n-F}_{\infty}
\leq
2\sqrt{
\frac{2\log(n+1)}{n}
}
+
\delta.
\]

\begin{remark}
This gives a quantitative uniform law of large numbers for empirical distribution functions.
\end{remark}

\chapter{VC Theory}

\section{Motivation}

Rademacher complexity provides a way to control uniform deviations of the form
\[
\norm{P_n-P}_{\mathcal{F}}
=
\sup_{f\in\mathcal{F}}
\left|
\frac{1}{n}\sum_{i=1}^n f(X_i)-\E[f(X)]
\right|.
\]

In the previous chapter, we bounded Rademacher complexity for classes with polynomial discrimination. We now develop a method for certifying polynomial discrimination for classes of binary-valued functions.

\section{Shattering and VC Dimension}

Let
\[
\mathcal{F}
\]
be a class of functions
\[
f:\mathcal{X}\to\{0,1\}.
\]

For points
\[
x_1,\dots,x_n\in\mathcal{X},
\]
define
\[
\mathcal{F}(x_1,\dots,x_n)
=
\left\{
(f(x_1),\dots,f(x_n)):f\in\mathcal{F}
\right\}
\subseteq \{0,1\}^n.
\]

\begin{definition}[Shattering]
A set
\[
\{x_1,\dots,x_n\}\subseteq\mathcal{X}
\]
is shattered by
\[
\mathcal{F}
\]
if
\[
\#\mathcal{F}(x_1,\dots,x_n)=2^n.
\]

Equivalently, every binary labeling of
\[
x_1,\dots,x_n
\]
can be realized by some function in
\[
\mathcal{F}.
\]
\end{definition}

\begin{definition}[VC Dimension]
The VC dimension of
\[
\mathcal{F}
\]
is
\[
\operatorname{vc}(\mathcal{F})
=
\sup\left\{
n\in\N:
\exists x_1,\dots,x_n\in\mathcal{X}
\text{ shattered by }\mathcal{F}
\right\}.
\]
\end{definition}

\begin{remark}
The VC dimension measures the combinatorial richness of a binary-valued function class. Large VC dimension means the class can realize many different label patterns.
\end{remark}

\section{Set Classes and Indicator Functions}

Let
\[
\mathcal{S}
\]
be a class of subsets of
\[
\mathcal{X}.
\]

Associated to
\[
\mathcal{S}
\]
is the indicator class
\[
\mathcal{F}_{\mathcal{S}}
=
\{\ind_A:A\in\mathcal{S}\}.
\]

For a finite set
\[
X=\{x_1,\dots,x_n\},
\]
define
\[
\mathcal{S}(X)
=
\{A\cap X:A\in\mathcal{S}\}.
\]

Then
\[
\#\mathcal{S}(X)
=
\#\mathcal{F}_{\mathcal{S}}(x_1,\dots,x_n),
\]
and
\[
\operatorname{vc}(\mathcal{S})
=
\operatorname{vc}(\mathcal{F}_{\mathcal{S}}).
\]

\section{Examples}

\begin{example}[Half-Lines]
Let
\[
\mathcal{S}_1
=
\{(-\infty,t]:t\in\R\}.
\]

Then
\[
\operatorname{vc}(\mathcal{S}_1)=1.
\]

Indeed, any single point can be shattered. Given one point
\[
x_1,
\]
we can realize the label $1$ by choosing
\[
t\geq x_1
\]
and the label $0$ by choosing
\[
t<x_1.
\]

However, two points cannot be shattered. If
\[
x_1<x_2,
\]
then any set of the form
\[
(-\infty,t]
\]
that contains
\[
x_2
\]
must also contain
\[
x_1.
\]
Thus the labeling
\[
(0,1)
\]
cannot be realized.
\end{example}

\begin{example}[Intervals]
Let
\[
\mathcal{S}_2
=
\{(a,b]:a,b\in\R,\ a<b\}.
\]

Then
\[
\operatorname{vc}(\mathcal{S}_2)=2.
\]

To see that two points can be shattered, take
\[
x_1<x_2.
\]

The four possible labelings are realized as follows:
\[
(0,0)
\]
by choosing an interval avoiding both points;
\[
(1,0)
\]
by choosing a small interval containing only $x_1$;
\[
(0,1)
\]
by choosing a small interval containing only $x_2$;
and
\[
(1,1)
\]
by choosing an interval containing both points.

However, three points cannot be shattered. If
\[
x_1<x_2<x_3,
\]
then the labeling
\[
(1,0,1)
\]
cannot be realized by a single interval, since containing both outer points forces the interval to contain the middle point.
\end{example}

\section{The Shatter Function}

\begin{definition}[Shatter Function]
For a binary-valued function class
\[
\mathcal{F},
\]
the shatter function is
\[
\Pi_{\mathcal{F}}(n)
=
\max_{x_1,\dots,x_n\in\mathcal{X}}
\#\mathcal{F}(x_1,\dots,x_n).
\]
\end{definition}

The VC dimension is the largest
\[
n
\]
such that
\[
\Pi_{\mathcal{F}}(n)=2^n.
\]

The key result is that finite VC dimension implies polynomial growth of the shatter function.

\section{Sauer--Shelah Lemma}

\begin{theorem}[Sauer--Shelah Lemma]
Suppose
\[
d=\operatorname{vc}(\mathcal{F})<\infty.
\]

Then, for any
\[
x_1,\dots,x_n\in\mathcal{X},
\]
\[
\#\mathcal{F}(x_1,\dots,x_n)
\leq
\sum_{i=0}^d \binom{n}{i}.
\]

Consequently,
\[
\#\mathcal{F}(x_1,\dots,x_n)
\leq
(n+1)^d.
\]
\end{theorem}

\begin{proof}
We prove the first inequality. Define
\[
\Phi_d(n)
=
\sum_{i=0}^d \binom{n}{i}.
\]

We prove by induction on
\[
d+n.
\]

The base cases are straightforward.

If
\[
n=0,
\]
then there is only the empty labeling, so
\[
\#\mathcal{F}(\emptyset)
\leq 1
=
\Phi_d(0).
\]

If
\[
d=0,
\]
then no single point can be shattered. Hence all points must receive the same labeling pattern across the class, so
\[
\#\mathcal{F}(x_1,\dots,x_n)\leq 1
=
\Phi_0(n).
\]

Now assume the result holds for all pairs
\[
(d',n')
\]
with
\[
d'+n'<d+n.
\]

Fix
\[
X=\{x_1,\dots,x_n\}
\]
and choose
\[
a\in X.
\]

Let
\[
X^-=X\setminus\{a\}.
\]

Let
\[
H=\mathcal{F}(X)
\subseteq\{0,1\}^X
\]
be the set of labelings realized on
\[
X.
\]

Define
\[
H_0
=
\left\{
h|_{X^-}:h\in H,\ h(a)=0
\right\}
\]
and
\[
H_1
=
\left\{
h|_{X^-}:h\in H,\ h(a)=1
\right\}.
\]

Let
\[
H_{\cap}=H_0\cap H_1.
\]

The set
\[
H_{\cap}
\]
contains precisely those labelings of
\[
X^-
\]
that can be extended to
\[
X
\]
with both possible labels at
\[
a.
\]

Every labeling in
\[
H
\]
corresponds either to a labeling of
\[
X^-
\]
that appears only once, or to a labeling of
\[
X^-
\]
that appears twice, once with label $0$ at $a$ and once with label $1$ at $a$.

Therefore,
\[
\#H
=
\#(H_0\cup H_1)+\#H_{\cap}.
\]

Since
\[
H_0\cup H_1
=
\mathcal{F}(X^-),
\]
we get
\[
\#H
=
\#\mathcal{F}(X^-)+\#H_{\cap}.
\]

By the induction hypothesis,
\[
\#\mathcal{F}(X^-)
\leq
\Phi_d(n-1).
\]

It remains to bound
\[
\#H_{\cap}.
\]

Define a new class
\[
\mathcal{F}_a
\]
on
\[
X^-
\]
as follows:
a labeling
\[
g:X^-\to\{0,1\}
\]
belongs to
\[
\mathcal{F}_a
\]
if there exist
\[
f_0,f_1\in\mathcal{F}
\]
such that
\[
f_0(a)\neq f_1(a)
\]
and
\[
f_0(x)=f_1(x)=g(x)
\]
for every
\[
x\in X^-.
\]

By construction,
\[
H_{\cap}
\subseteq
\mathcal{F}_a(X^-).
\]

We claim that
\[
\operatorname{vc}(\mathcal{F}_a)\leq d-1.
\]

Suppose not. Then
\[
\mathcal{F}_a
\]
shatters some set
\[
Y\subseteq X^-
\]
with
\[
|Y|=d.
\]

For each labeling of
\[
Y,
\]
there are two extensions in
\[
\mathcal{F}
\]
that agree on
\[
Y
\]
and differ at
\[
a.
\]

Thus
\[
\mathcal{F}
\]
would shatter
\[
Y\cup\{a\},
\]
which has size
\[
d+1.
\]

This contradicts
\[
\operatorname{vc}(\mathcal{F})=d.
\]

Hence
\[
\operatorname{vc}(\mathcal{F}_a)\leq d-1.
\]

By the induction hypothesis,
\[
\#H_{\cap}
\leq
\#\mathcal{F}_a(X^-)
\leq
\Phi_{d-1}(n-1).
\]

Combining the bounds,
\[
\#\mathcal{F}(X)
\leq
\Phi_d(n-1)+\Phi_{d-1}(n-1).
\]

Using Pascal's identity,
\[
\binom{n}{i}
=
\binom{n-1}{i}
+
\binom{n-1}{i-1},
\]
we get
\[
\Phi_d(n-1)+\Phi_{d-1}(n-1)
=
\sum_{i=0}^d \binom{n-1}{i}
+
\sum_{i=0}^{d-1}\binom{n-1}{i}
=
\sum_{i=0}^d \binom{n}{i}
=
\Phi_d(n).
\]

Thus,
\[
\#\mathcal{F}(X)
\leq
\sum_{i=0}^d \binom{n}{i}.
\]

Finally,
\[
\sum_{i=0}^d \binom{n}{i}
\leq
(n+1)^d
\]
for fixed finite $d$, giving the polynomial discrimination bound.
\end{proof}

\section{VC Dimension Implies Polynomial Discrimination}

\begin{corollary}
If
\[
\operatorname{vc}(\mathcal{F})=d<\infty,
\]
then
\[
\mathcal{F}
\]
has polynomial discrimination of order
\[
d.
\]
\end{corollary}

\begin{proof}
By the Sauer--Shelah lemma,
\[
\#\mathcal{F}(x_1,\dots,x_n)
\leq
(n+1)^d
\]
for every
\[
x_1,\dots,x_n\in\mathcal{X}.
\]

This is exactly polynomial discrimination of order $d$.
\end{proof}

\section{Rademacher Complexity Bound from VC Dimension}

\begin{corollary}
Let
\[
\mathcal{F}
\]
be a binary-valued function class with
\[
\operatorname{vc}(\mathcal{F})=d<\infty.
\]

Then
\[
\mathcal{R}_n(\mathcal{F})
\leq
\sqrt{
\frac{2d\log(n+1)}{n}
}.
\]
\end{corollary}

\begin{proof}
Since
\[
\mathcal{F}
\]
is binary-valued,
\[
|f(x)|\leq1
\]
for every
\[
f\in\mathcal{F}.
\]

By the previous corollary,
\[
\mathcal{F}
\]
has polynomial discrimination of order $d$.

Applying the bounded polynomial-discrimination Rademacher bound gives
\[
\mathcal{R}_n(\mathcal{F})
\leq
\sqrt{
\frac{2d\log(n+1)}{n}
}.
\]
\end{proof}

\section{Uniform Deviation Bound for VC Classes}

Combining Rademacher complexity with McDiarmid's inequality gives a high-probability uniform law of large numbers.

\begin{theorem}
Let
\[
\mathcal{F}
\]
be a binary-valued function class with
\[
\operatorname{vc}(\mathcal{F})=d<\infty.
\]

Then, for every
\[
t\geq0,
\]
with probability at least
\[
1-\exp\left(-\frac{nt^2}{2}\right),
\]
\[
\norm{P_n-P}_{\mathcal{F}}
\leq
2\sqrt{
\frac{2d\log(n+1)}{n}
}
+
t.
\]
\end{theorem}

\begin{proof}
The high-probability Rademacher bound for bounded classes gives
\[
\norm{P_n-P}_{\mathcal{F}}
\leq
2\mathcal{R}_n(\mathcal{F})+t
\]
with probability at least
\[
1-\exp\left(-\frac{nt^2}{2}\right),
\]
because
\[
|f(x)|\leq1.
\]

Using
\[
\mathcal{R}_n(\mathcal{F})
\leq
\sqrt{
\frac{2d\log(n+1)}{n}
},
\]
the result follows.
\end{proof}

\begin{remark}
VC theory converts a combinatorial property, finite VC dimension, into a probabilistic statement: uniform convergence of empirical averages to expectations.
\end{remark}

\section{Controlling the VC Dimension}

We now develop tools for proving that a binary-valued function class has finite VC dimension.

Let
\[
\mathcal{G}
\]
be a class of real-valued functions
\[
g:\mathcal{X}\to\R.
\]

For each
\[
g\in\mathcal{G},
\]
define the sublevel set
\[
S_g
=
\{x\in\mathcal{X}:g(x)\leq0\}.
\]

Let
\[
\mathcal{S}(\mathcal{G})
=
\{S_g:g\in\mathcal{G}\}.
\]

Thus
\[
\mathcal{S}(\mathcal{G})
\]
is a class of sets generated by thresholding functions in
\[
\mathcal{G}.
\]

\begin{proposition}
Let
\[
\mathcal{G}
\]
be a vector space of real-valued functions on
\[
\mathcal{X}
\]
with
\[
\dim(\mathcal{G})<\infty.
\]

Then
\[
\operatorname{vc}(\mathcal{S}(\mathcal{G}))
\leq
\dim(\mathcal{G}).
\]
\end{proposition}

\begin{proof}
Let
\[
m=\dim(\mathcal{G}).
\]

Suppose, for contradiction, that
\[
\mathcal{S}(\mathcal{G})
\]
shatters a set
\[
X=\{x_1,\dots,x_{m+1}\}.
\]

Define the linear map
\[
L:\mathcal{G}\to\R^{m+1}
\]
by
\[
L(g)
=
\begin{pmatrix}
g(x_1)\\
\vdots\\
g(x_{m+1})
\end{pmatrix}.
\]

Since
\[
\dim(\mathcal{G})=m<m+1,
\]
the image
\[
L(\mathcal{G})
\]
is a subspace of
\[
\R^{m+1}
\]
of dimension at most
\[
m.
\]

Therefore, there exists a nonzero vector
\[
\gamma=(\gamma_1,\dots,\gamma_{m+1})\in\R^{m+1}
\]
such that
\[
\langle \gamma,L(g)\rangle=0
\]
for every
\[
g\in\mathcal{G}.
\]

Equivalently,
\[
\sum_{i=1}^{m+1}\gamma_i g(x_i)=0
\]
for every
\[
g\in\mathcal{G}.
\]

Define
\[
\Delta_+
=
\{i:\gamma_i>0\},
\qquad
\Delta_-
=
\{i:\gamma_i\leq0\}.
\]

Since
\[
\gamma\neq0,
\]
at least one of these sets contains an index with nonzero coefficient. Without loss of generality, assume
\[
\Delta_+\neq\emptyset.
\]

Because
\[
\mathcal{S}(\mathcal{G})
\]
shatters
\[
X,
\]
there exists
\[
g\in\mathcal{G}
\]
such that
\[
g(x_i)>0
\quad\text{for } i\in\Delta_+,
\]
and
\[
g(x_i)\leq0
\quad\text{for } i\in\Delta_-.
\]

Then
\[
\sum_{i\in\Delta_+}\gamma_i g(x_i)>0,
\]
while
\[
-\sum_{i\in\Delta_-}\gamma_i g(x_i)\geq0.
\]

But the identity
\[
\sum_{i=1}^{m+1}\gamma_i g(x_i)=0
\]
implies
\[
\sum_{i\in\Delta_+}\gamma_i g(x_i)
=
-\sum_{i\in\Delta_-}\gamma_i g(x_i).
\]

This gives a contradiction, since the left-hand side is strictly positive while the right-hand side is nonnegative but cannot match the required sign pattern induced by shattering.

Therefore,
\[
\mathcal{S}(\mathcal{G})
\]
cannot shatter any set of size
\[
m+1.
\]

Hence
\[
\operatorname{vc}(\mathcal{S}(\mathcal{G}))
\leq
m
=
\dim(\mathcal{G}).
\]
\end{proof}

\section{Example: Halfspaces}

For
\[
a\in\R^d
\]
and
\[
b\in\R,
\]
define the halfspace
\[
S_{a,b}
=
\{x\in\R^d:\langle a,x\rangle+b\leq0\}.
\]

Let
\[
\mathcal{S}
=
\{S_{a,b}:a\in\R^d,\ b\in\R\}.
\]

This set class is generated by the function class
\[
\mathcal{G}
=
\{x\mapsto \langle a,x\rangle+b:a\in\R^d,\ b\in\R\}.
\]

The class
\[
\mathcal{G}
\]
is a vector space of dimension
\[
d+1.
\]

Therefore,
\[
\operatorname{vc}(\mathcal{S})
\leq
d+1.
\]

\begin{remark}
The exact VC dimension of affine halfspaces in
\[
\R^d
\]
is
\[
d+1.
\]

The upper bound follows from the vector-space argument above. The matching lower bound follows by showing that
\[
d+1
\]
points in general position can be shattered by affine halfspaces.
\end{remark}

\section{Example: Euclidean Balls}

For
\[
a\in\R^d
\]
and
\[
b\geq0,
\]
define the Euclidean ball
\[
S_{a,b}
=
\{x\in\R^d:\norm{x-a}_2\leq b\}.
\]

Equivalently,
\[
x\in S_{a,b}
\]
if and only if
\[
\norm{x-a}_2^2-b^2\leq0.
\]

Define
\[
g_{a,b}(x)
=
\norm{x-a}_2^2-b^2.
\]

Expanding,
\[
g_{a,b}(x)
=
\norm{x}_2^2
-
2\langle a,x\rangle
+
\norm{a}_2^2
-
b^2.
\]

This function class is not itself a vector space in the parameters
\[
(a,b),
\]
but it is contained in a finite-dimensional vector space.

Define the feature map
\[
\phi(x)
=
\begin{pmatrix}
1\\
x_1\\
\vdots\\
x_d\\
\norm{x}_2^2
\end{pmatrix}
\in\R^{d+2}.
\]

Then every
\[
g_{a,b}
\]
can be written as
\[
g_{a,b}(x)
=
\langle c,\phi(x)\rangle
\]
for some
\[
c\in\R^{d+2}.
\]

Thus the ball class is contained in the threshold class generated by a vector space of dimension
\[
d+2.
\]

Therefore,
\[
\operatorname{vc}(\mathcal{S})
\leq
d+2.
\]

\begin{remark}
The exact VC dimension of Euclidean balls in
\[
\R^d
\]
is
\[
d+1.
\]

The finite-dimensional embedding argument gives the simple upper bound
\[
d+2.
\]
\end{remark}

\section{Dimension-Free Generalization Bounds}

For many geometric classes, VC theory gives bounds of the form
\[
\E\left[
\sup_{f\in\mathcal{F}}
\left|
\frac{1}{n}\sum_{i=1}^n f(X_i)-\E[f(X)]
\right|
\right]
=
O\left(
\sqrt{\frac{d\log n}{n}}
\right),
\]
where
\[
d
\]
is the ambient input dimension.

A natural question is whether one can obtain dimension-free bounds for interesting classes.

The answer is yes, under additional structure. We now study one important case: linear prediction with norm constraints.

\section{Linear Models with Lipschitz Loss}

Consider losses of the form
\[
\ell(\langle \theta,a\rangle,b),
\]
where
\[
(a,b)\sim P,
\]
\[
\theta\in\Theta,
\]
and the loss
\[
\ell(u,b)
\]
is $1$-Lipschitz in its first argument.

The population problem is
\[
\min_{\theta\in\Theta}
\E_{(a,b)\sim P}
\left[
\ell(\langle \theta,a\rangle,b)
\right].
\]

The empirical problem is
\[
\min_{\theta\in\Theta}
\frac{1}{n}
\sum_{i=1}^n
\ell(\langle \theta,a_i\rangle,b_i).
\]

To control generalization, define the function class
\[
\mathcal{F}
=
\left\{
(a,b)\mapsto
\ell(\langle \theta,a\rangle,b)
:
\theta\in\Theta
\right\}.
\]

By symmetrization,
\[
\E\norm{P_n-P}_{\mathcal{F}}
\leq
2
\E
\left[
\sup_{\theta\in\Theta}
\left|
\frac{1}{n}
\sum_{i=1}^n
\varepsilon_i
\ell(\langle \theta,a_i\rangle,b_i)
\right|
\right].
\]

To bound this quantity, we use the contraction principle.

\section{Rademacher Complexity Without Absolute Values}

It is often useful to define the signed Rademacher complexity
\[
\overline{\mathcal{R}}(A)
=
\E_\varepsilon
\left[
\sup_{a\in A}
\sum_{i=1}^n
\varepsilon_i a_i
\right].
\]

Similarly, for a function class,
\[
\overline{\mathcal{R}}_n(\mathcal{F})
=
\E_{X,\varepsilon}
\left[
\sup_{f\in\mathcal{F}}
\frac{1}{n}
\sum_{i=1}^n
\varepsilon_i f(X_i)
\right].
\]

\begin{remark}
The signed version is especially convenient for contraction inequalities. It still gives useful information because symmetrization can be developed with or without absolute values, depending on whether the function class is symmetric.
\end{remark}

\section{The Contraction Principle}

\begin{theorem}[Contraction Principle]
Let
\[
A\subseteq\R^n.
\]

For each
\[
i=1,\dots,n,
\]
let
\[
\psi_i:\R\to\R
\]
be $\rho$-Lipschitz.

Define
\[
A'
=
\left\{
(\psi_1(a_1),\dots,\psi_n(a_n)):a\in A
\right\}.
\]

Then
\[
\overline{\mathcal{R}}(A')
\leq
\rho\,
\overline{\mathcal{R}}(A).
\]
\end{theorem}

\begin{proof}
By rescaling, it suffices to prove the result for
\[
\rho=1.
\]

For
\[
\psi=(\psi_1,\dots,\psi_n),
\]
define
\[
\psi(A)
=
\{(\psi_1(a_1),\dots,\psi_n(a_n)):a\in A\}.
\]

We prove
\[
\overline{\mathcal{R}}(\psi(A))
\leq
\overline{\mathcal{R}}(A)
\]
by replacing the functions one coordinate at a time.

Let
\[
\widetilde{\psi}
=
(I,\psi_2,\dots,\psi_n),
\]
where the first coordinate has been replaced by the identity map.

By permutation invariance, it is enough to show
\[
\overline{\mathcal{R}}(\psi(A))
\leq
\overline{\mathcal{R}}(\widetilde{\psi}(A)).
\]

Expanding over the first Rademacher variable,
\[
\overline{\mathcal{R}}(\psi(A))
=
\E_{\varepsilon_2,\dots,\varepsilon_n}
\left[
\frac{1}{2}
\sup_{a\in A}
\left(
\psi_1(a_1)
+
\sum_{i=2}^n
\varepsilon_i\psi_i(a_i)
\right)
\right.
\]
\[
\left.
+
\frac{1}{2}
\sup_{a\in A}
\left(
-\psi_1(a_1)
+
\sum_{i=2}^n
\varepsilon_i\psi_i(a_i)
\right)
\right].
\]

The two suprema may be written using two independent choices
\[
a,\widetilde{a}\in A.
\]

Thus the inner quantity equals
\[
\frac{1}{2}
\sup_{a,\widetilde{a}\in A}
\left[
\psi_1(a_1)-\psi_1(\widetilde{a}_1)
+
\sum_{i=2}^n
\varepsilon_i
\left(
\psi_i(a_i)+\psi_i(\widetilde{a}_i)
\right)
\right].
\]

Since
\[
\psi_1
\]
is $1$-Lipschitz,
\[
\psi_1(a_1)-\psi_1(\widetilde{a}_1)
\leq
|a_1-\widetilde{a}_1|.
\]

Because the expression is optimized over both
\[
a
\]
and
\[
\widetilde{a},
\]
we may choose the ordering so that
\[
a_1-\widetilde{a}_1\geq0.
\]

Therefore,
\[
|\psi_1(a_1)-\psi_1(\widetilde{a}_1)|
\leq
a_1-\widetilde{a}_1.
\]

Thus,
\[
\overline{\mathcal{R}}(\psi(A))
\leq
\E_{\varepsilon_2,\dots,\varepsilon_n}
\left[
\frac{1}{2}
\sup_{a,\widetilde{a}\in A}
\left(
a_1-\widetilde{a}_1
+
\sum_{i=2}^n
\varepsilon_i
(\psi_i(a_i)+\psi_i(\widetilde{a}_i))
\right)
\right].
\]

Now rewrite the right-hand side as
\[
\E_{\varepsilon_2,\dots,\varepsilon_n}
\left[
\frac{1}{2}
\sup_{a\in A}
\left(
a_1+
\sum_{i=2}^n
\varepsilon_i\psi_i(a_i)
\right)
+
\frac{1}{2}
\sup_{\widetilde{a}\in A}
\left(
-\widetilde{a}_1+
\sum_{i=2}^n
\varepsilon_i\psi_i(\widetilde{a}_i)
\right)
\right].
\]

This is precisely
\[
\overline{\mathcal{R}}(\widetilde{\psi}(A)).
\]

Repeating the argument coordinate by coordinate gives
\[
\overline{\mathcal{R}}(\psi(A))
\leq
\overline{\mathcal{R}}(A).
\]

Rescaling gives the result for general
\[
\rho>0.
\]
\end{proof}

\section{Application to Linear Models}

Let
\[
\ell(u,b)
\]
be $1$-Lipschitz in
\[
u.
\]

For fixed samples
\[
(a_1,b_1),\dots,(a_n,b_n),
\]
define
\[
A
=
\left\{
(\langle \theta,a_1\rangle,\dots,\langle \theta,a_n\rangle):
\theta\in\Theta
\right\}.
\]

Using the contraction principle with
\[
\psi_i(u)=\ell(u,b_i),
\]
we obtain
\[
\overline{\mathcal{R}}_n(\mathcal{F})
\leq
\overline{\mathcal{R}}_n
\left(
\left\{
a\mapsto \langle \theta,a\rangle:
\theta\in\Theta
\right\}
\right).
\]

Thus the generalization behavior of Lipschitz losses for linear predictors is controlled by the Rademacher complexity of the underlying linear class.

\begin{remark}
This is a dimension-free reduction: the loss function does not increase complexity beyond the complexity of the linear score class, provided the loss is Lipschitz.
\end{remark}

\section{Norm-Constrained Linear Classes}

Suppose
\[
\Theta
=
\{\theta\in\R^d:\norm{\theta}_2\leq B\}.
\]

Consider the linear class
\[
\mathcal{H}
=
\{a\mapsto \langle \theta,a\rangle:\norm{\theta}_2\leq B\}.
\]

Then
\[
\overline{\mathcal{R}}_n(\mathcal{H})
=
\E
\left[
\sup_{\norm{\theta}_2\leq B}
\frac{1}{n}
\sum_{i=1}^n
\varepsilon_i
\langle \theta,a_i\rangle
\right].
\]

Rearranging,
\[
\overline{\mathcal{R}}_n(\mathcal{H})
=
\E
\left[
\sup_{\norm{\theta}_2\leq B}
\left\langle
\theta,
\frac{1}{n}
\sum_{i=1}^n
\varepsilon_i a_i
\right\rangle
\right].
\]

By Cauchy--Schwarz,
\[
\sup_{\norm{\theta}_2\leq B}
\left\langle
\theta,
\frac{1}{n}
\sum_{i=1}^n
\varepsilon_i a_i
\right\rangle
=
B
\left\|
\frac{1}{n}
\sum_{i=1}^n
\varepsilon_i a_i
\right\|_2.
\]

Therefore,
\[
\overline{\mathcal{R}}_n(\mathcal{H})
=
B
\E
\left\|
\frac{1}{n}
\sum_{i=1}^n
\varepsilon_i a_i
\right\|_2.
\]

By Jensen's inequality,
\[
\E
\left\|
\frac{1}{n}
\sum_{i=1}^n
\varepsilon_i a_i
\right\|_2
\leq
\left(
\E
\left\|
\frac{1}{n}
\sum_{i=1}^n
\varepsilon_i a_i
\right\|_2^2
\right)^{1/2}.
\]

Expanding,
\[
\E_\varepsilon
\left\|
\frac{1}{n}
\sum_{i=1}^n
\varepsilon_i a_i
\right\|_2^2
=
\frac{1}{n^2}
\E_\varepsilon
\left[
\sum_{i,j=1}^n
\varepsilon_i\varepsilon_j
\langle a_i,a_j\rangle
\right].
\]

Since
\[
\E[\varepsilon_i\varepsilon_j]=0
\quad\text{for } i\neq j,
\]
and
\[
\E[\varepsilon_i^2]=1,
\]
we get
\[
\E_\varepsilon
\left\|
\frac{1}{n}
\sum_{i=1}^n
\varepsilon_i a_i
\right\|_2^2
=
\frac{1}{n^2}
\sum_{i=1}^n
\norm{a_i}_2^2.
\]

Thus,
\[
\overline{\mathcal{R}}_n(\mathcal{H})
\leq
\frac{B}{n}
\left(
\sum_{i=1}^n
\norm{a_i}_2^2
\right)^{1/2}.
\]

Taking expectation over the data gives
\[
\overline{\mathcal{R}}_n(\mathcal{H})
\leq
\frac{B}{\sqrt{n}}
\left(
\E\norm{a}_2^2
\right)^{1/2}.
\]

\begin{theorem}[Dimension-Free Linear Generalization Bound]
Suppose
\[
\ell(u,b)
\]
is $1$-Lipschitz in $u$, and suppose
\[
\Theta=\{\theta:\norm{\theta}_2\leq B\}.
\]

If
\[
\E\norm{a}_2^2<\infty,
\]
then
\[
\mathcal{R}_n(\mathcal{F})
\leq
\frac{B}{\sqrt{n}}
\left(
\E\norm{a}_2^2
\right)^{1/2},
\]
up to universal constants, where
\[
\mathcal{F}
=
\{(a,b)\mapsto \ell(\langle \theta,a\rangle,b):\theta\in\Theta\}.
\]
\end{theorem}

\begin{remark}
The bound depends on the norm radius
\[
B,
\]
the second moment
\[
\E\norm{a}_2^2,
\]
and the sample size
\[
n,
\]
but not explicitly on the ambient dimension
\[
d.
\]
\end{remark}

\section{Linear Models Continued}

\begin{lemma}
Let
\[
X=\{x_1,\dots,x_n\}\subseteq\R^d
\]
and define
\[
\mathcal{H}
=
\{a\mapsto \langle w,a\rangle:\norm{w}_2\leq 1\}.
\]

Then
\[
\overline{\mathcal{R}}(\mathcal{H}(X))
\leq
\left(
\sum_{i=1}^n \norm{x_i}_2^2
\right)^{1/2}.
\]

Consequently,
\[
\overline{\mathcal{R}}_n(\mathcal{H})
\leq
\frac{1}{\sqrt{n}}
\E\left[
\max_{1\leq i\leq n}\norm{X_i}_2
\right].
\]
\end{lemma}

\begin{proof}
By definition,
\[
\overline{\mathcal{R}}(\mathcal{H}(X))
=
\E_\varepsilon
\sup_{\norm{w}_2\leq 1}
\sum_{i=1}^n
\varepsilon_i \langle w,x_i\rangle.
\]

Rearranging,
\[
\overline{\mathcal{R}}(\mathcal{H}(X))
=
\E_\varepsilon
\sup_{\norm{w}_2\leq1}
\left\langle
w,
\sum_{i=1}^n \varepsilon_i x_i
\right\rangle.
\]

By Cauchy--Schwarz,
\[
\sup_{\norm{w}_2\leq1}
\left\langle
w,
\sum_{i=1}^n \varepsilon_i x_i
\right\rangle
=
\left\|
\sum_{i=1}^n \varepsilon_i x_i
\right\|_2.
\]

Thus,
\[
\overline{\mathcal{R}}(\mathcal{H}(X))
=
\E_\varepsilon
\left\|
\sum_{i=1}^n \varepsilon_i x_i
\right\|_2.
\]

By Jensen's inequality,
\[
\E_\varepsilon
\left\|
\sum_{i=1}^n \varepsilon_i x_i
\right\|_2
\leq
\left(
\E_\varepsilon
\left\|
\sum_{i=1}^n \varepsilon_i x_i
\right\|_2^2
\right)^{1/2}.
\]

Expanding the square,
\[
\E_\varepsilon
\left\|
\sum_{i=1}^n \varepsilon_i x_i
\right\|_2^2
=
\E_\varepsilon
\sum_{i,j=1}^n
\varepsilon_i\varepsilon_j
\langle x_i,x_j\rangle.
\]

Since
\[
\E[\varepsilon_i\varepsilon_j]=0
\quad\text{for } i\neq j,
\qquad
\E[\varepsilon_i^2]=1,
\]
we get
\[
\E_\varepsilon
\left\|
\sum_{i=1}^n \varepsilon_i x_i
\right\|_2^2
=
\sum_{i=1}^n \norm{x_i}_2^2.
\]

Therefore,
\[
\overline{\mathcal{R}}(\mathcal{H}(X))
\leq
\left(
\sum_{i=1}^n \norm{x_i}_2^2
\right)^{1/2}.
\]

Since
\[
\overline{\mathcal{R}}_n(\mathcal{H})
=
\E_X
\overline{\mathcal{R}}
\left(
\frac{1}{n}\mathcal{H}(X)
\right),
\]
we obtain
\[
\overline{\mathcal{R}}_n(\mathcal{H})
\leq
\frac{1}{n}
\E
\left(
\sum_{i=1}^n \norm{X_i}_2^2
\right)^{1/2}.
\]

Using
\[
\sum_{i=1}^n \norm{X_i}_2^2
\leq
n\max_{1\leq i\leq n}\norm{X_i}_2^2,
\]
we get
\[
\overline{\mathcal{R}}_n(\mathcal{H})
\leq
\frac{1}{\sqrt{n}}
\E\left[
\max_{1\leq i\leq n}\norm{X_i}_2
\right].
\]
\end{proof}

\begin{remark}
This bound is dimension-free in the sense that it depends on the sizes of the observed covariates rather than explicitly on the ambient dimension.
\end{remark}

\section{Linear Classes with an \texorpdfstring{$\ell_1$}{l1} Constraint}

We can obtain an almost dimension-free bound by constraining the predictor using the $\ell_1$ norm.

\begin{lemma}
Let
\[
X=\{x_1,\dots,x_n\}\subseteq\R^d
\]
and define
\[
\mathcal{H}
=
\{a\mapsto \langle w,a\rangle:\norm{w}_1\leq 1\}.
\]

Then
\[
\overline{\mathcal{R}}(\mathcal{H}(X))
\leq
\sqrt{2n\log(2d)}
\max_{1\leq i\leq n}\norm{x_i}_{\infty}.
\]

Consequently,
\[
\overline{\mathcal{R}}_n(\mathcal{H})
\leq
\sqrt{\frac{2\log(2d)}{n}}
\E\left[
\max_{1\leq i\leq n}\norm{X_i}_{\infty}
\right].
\]
\end{lemma}

\begin{proof}
We compute
\[
\overline{\mathcal{R}}(\mathcal{H}(X))
=
\E_\varepsilon
\sup_{\norm{w}_1\leq1}
\sum_{i=1}^n
\varepsilon_i\langle w,x_i\rangle.
\]

Rearranging,
\[
=
\E_\varepsilon
\sup_{\norm{w}_1\leq1}
\left\langle
w,
\sum_{i=1}^n\varepsilon_i x_i
\right\rangle.
\]

By duality between $\ell_1$ and $\ell_\infty$,
\[
\sup_{\norm{w}_1\leq1}
\left\langle
w,
\sum_{i=1}^n\varepsilon_i x_i
\right\rangle
=
\left\|
\sum_{i=1}^n\varepsilon_i x_i
\right\|_\infty.
\]

Thus,
\[
\overline{\mathcal{R}}(\mathcal{H}(X))
=
\E_\varepsilon
\left\|
\sum_{i=1}^n\varepsilon_i x_i
\right\|_\infty.
\]

Writing coordinates,
\[
\left\|
\sum_{i=1}^n\varepsilon_i x_i
\right\|_\infty
=
\max_{1\leq j\leq d}
\left|
\sum_{i=1}^n\varepsilon_i x_{ij}
\right|.
\]

For each fixed coordinate $j$,
\[
\sum_{i=1}^n\varepsilon_i x_{ij}
\]
is sub-Gaussian with variance proxy
\[
\sum_{i=1}^n x_{ij}^2.
\]

Since
\[
|x_{ij}|\leq \norm{x_i}_\infty,
\]
we have
\[
\sum_{i=1}^n x_{ij}^2
\leq
n\max_{1\leq i\leq n}\norm{x_i}_\infty^2.
\]

Applying the standard maximum bound for sub-Gaussian random variables over the $2d$ signed coordinates gives
\[
\E_\varepsilon
\max_{1\leq j\leq d}
\left|
\sum_{i=1}^n\varepsilon_i x_{ij}
\right|
\leq
\sqrt{2\log(2d)}
\sqrt{
n\max_{1\leq i\leq n}\norm{x_i}_\infty^2
}.
\]

Therefore,
\[
\overline{\mathcal{R}}(\mathcal{H}(X))
\leq
\sqrt{2n\log(2d)}
\max_{1\leq i\leq n}\norm{x_i}_\infty.
\]

Dividing by $n$ and taking expectation over the data yields
\[
\overline{\mathcal{R}}_n(\mathcal{H})
\leq
\sqrt{\frac{2\log(2d)}{n}}
\E\left[
\max_{1\leq i\leq n}\norm{X_i}_\infty
\right].
\]
\end{proof}

\begin{remark}
The dependence on dimension is only logarithmic. This explains why $\ell_1$-constrained linear models can behave well in high-dimensional settings.
\end{remark}

\section{Generalization via Convexity and Stability}

We now study a different approach to generalization. Instead of controlling the entire function class uniformly, we control the stability of a learning algorithm.

Suppose the population objective is
\[
F(\theta)
=
\E_{x\sim P}[f(\theta,x)],
\]
and the goal is to solve
\[
\min_{\theta\in\Theta}F(\theta).
\]

Given an iid sample
\[
X=(x_1,\dots,x_n),
\]
let an algorithm
\[
A
\]
output
\[
A(X)\in\Theta.
\]

We want to bound the generalization gap
\[
F(A(X))
-
\frac{1}{n}\sum_{i=1}^n f(A(X),x_i).
\]

\section{Replace-One Samples}

For each
\[
i\in\{1,\dots,n\},
\]
define
\[
X^{(i)}
=
(x_1,\dots,x_{i-1},x_i',x_{i+1},\dots,x_n),
\]
where
\[
x_i'\sim P
\]
is independent of
\[
X.
\]

\begin{theorem}[Generalization Gap and Replace-One Stability]
Let
\[
i\sim \mathrm{Unif}\{1,\dots,n\}
\]
independently of all samples. Then
\[
\E_X
\left[
F(A(X))
-
\frac{1}{n}
\sum_{i=1}^n
f(A(X),x_i)
\right]
=
\E_{X,X',i}
\left[
f(A(X^{(i)}),x_i)
-
f(A(X),x_i)
\right].
\]
\end{theorem}

\begin{proof}
For every fixed $i$,
\[
\E_X[F(A(X))]
=
\E_{X,x_i'}[f(A(X),x_i')].
\]

Since
\[
X^{(i)}
\]
has the same distribution as
\[
X,
\]
we may replace
\[
A(X)
\]
by
\[
A(X^{(i)})
\]
inside the expectation:
\[
\E_{X,x_i'}[f(A(X),x_i')]
=
\E_{X,x_i'}[f(A(X^{(i)}),x_i)].
\]

Therefore,
\[
\E_X[F(A(X))]
=
\E_{X,X'}[f(A(X^{(i)}),x_i)].
\]

Also,
\[
\E_X
\left[
\frac{1}{n}
\sum_{i=1}^n f(A(X),x_i)
\right]
=
\E_{X,i}[f(A(X),x_i)].
\]

Subtracting the two identities and averaging over
\[
i\sim \mathrm{Unif}\{1,\dots,n\}
\]
gives the claim.
\end{proof}

\section{Replace-One Stability}

\begin{definition}[Replace-One Stability]
An algorithm
\[
A
\]
is replace-one stable with rate
\[
\varepsilon(n)
\]
if
\[
\E_{X,X',i}
\left[
f(A(X^{(i)}),x_i)
-
f(A(X),x_i)
\right]
\leq
\varepsilon(n),
\]
where
\[
i\sim \mathrm{Unif}\{1,\dots,n\}.
\]
\end{definition}

\begin{remark}
If replacing one training observation changes the loss only slightly on average, then the algorithm generalizes well.
\end{remark}

\section{Regularized Empirical Risk Minimization}

Assume:
\begin{itemize}
\item for every $x$, the map
\[
\theta\mapsto f(\theta,x)
\]
is convex;
\item for every $x$, the map
\[
\theta\mapsto f(\theta,x)
\]
is $\rho$-Lipschitz;
\item the parameter set
\[
\Theta
\]
is closed and convex.
\end{itemize}

For
\[
\lambda>0,
\]
define the regularized empirical risk minimizer
\[
A(X)
=
\argmin_{\theta\in\Theta}
\left[
\frac{1}{n}\sum_{i=1}^n f(\theta,x_i)
+
\frac{\lambda}{2}\norm{\theta}_2^2
\right].
\]

\section{Strong Convexity}

\begin{definition}[Strong Convexity]
A function
\[
g:\R^d\to\R\cup\{\infty\}
\]
is $\alpha$-strongly convex if
\[
\theta\mapsto
g(\theta)-\frac{\alpha}{2}\norm{\theta}_2^2
\]
is convex.
\end{definition}

\begin{lemma}
Suppose
\[
g
\]
is $\alpha$-strongly convex and has minimizer
\[
\theta_*.
\]

Then, for every
\[
\theta,
\]
\[
\frac{\alpha}{2}\norm{\theta-\theta_*}_2^2
\leq
g(\theta)-g(\theta_*).
\]
\end{lemma}

\begin{proof}
Since $g$ is $\alpha$-strongly convex, for every subgradient
\[
s\in\partial g(\theta_*),
\]
we have
\[
g(\theta)
\geq
g(\theta_*)
+
\langle s,\theta-\theta_*\rangle
+
\frac{\alpha}{2}\norm{\theta-\theta_*}_2^2.
\]

Since
\[
\theta_*
\]
minimizes $g$, we may take
\[
0\in\partial g(\theta_*).
\]

Thus,
\[
g(\theta)
\geq
g(\theta_*)
+
\frac{\alpha}{2}\norm{\theta-\theta_*}_2^2.
\]

Rearranging gives the claim.
\end{proof}

\begin{remark}
The regularized empirical objective is $\lambda$-strongly convex by construction. Therefore, the algorithm
\[
A(X)
\]
is well-defined and has a unique output.
\end{remark}

\section{Stability of Regularized ERM}

\begin{theorem}
Under the convexity and Lipschitz assumptions above, the regularized ERM algorithm
\[
A(X)
=
\argmin_{\theta\in\Theta}
\left[
\frac{1}{n}\sum_{i=1}^n f(\theta,x_i)
+
\frac{\lambda}{2}\norm{\theta}_2^2
\right]
\]
is replace-one stable with rate
\[
\frac{2\rho^2}{\lambda n}.
\]

Consequently,
\[
\E_X
\left[
F(A(X))
-
\frac{1}{n}\sum_{i=1}^n f(A(X),x_i)
\right]
\leq
\frac{2\rho^2}{\lambda n}.
\]
\end{theorem}

\begin{proof}
By $\rho$-Lipschitzness,
\[
f(A(X^{(i)}),x_i)-f(A(X),x_i)
\leq
\rho\norm{A(X^{(i)})-A(X)}_2.
\]

Thus it suffices to show
\[
\norm{A(X^{(i)})-A(X)}_2
\leq
\frac{2\rho}{\lambda n}.
\]

Define
\[
f_X(\theta)
=
\frac{1}{n}\sum_{j=1}^n f(\theta,x_j)
+
\frac{\lambda}{2}\norm{\theta}_2^2.
\]

By strong convexity and the fact that $A(X)$ minimizes $f_X$,
\[
\frac{\lambda}{2}
\norm{A(X^{(i)})-A(X)}_2^2
\leq
f_X(A(X^{(i)}))-f_X(A(X)).
\]

Expanding the right-hand side and separating the replaced observation gives
\[
f_X(A(X^{(i)}))-f_X(A(X))
=
f_{X^{(i)}}(A(X^{(i)}))-f_{X^{(i)}}(A(X))
\]
\[
+
\frac{1}{n}
\left[
f(A(X^{(i)}),x_i)-f(A(X),x_i)
\right]
+
\frac{1}{n}
\left[
f(A(X),x_i')-f(A(X^{(i)}),x_i')
\right].
\]

Since
\[
A(X^{(i)})
\]
minimizes
\[
f_{X^{(i)}},
\]
we have
\[
f_{X^{(i)}}(A(X^{(i)}))-f_{X^{(i)}}(A(X))
\leq0.
\]

Therefore,
\[
\frac{\lambda}{2}
\norm{A(X^{(i)})-A(X)}_2^2
\leq
\frac{1}{n}
\left[
f(A(X^{(i)}),x_i)-f(A(X),x_i)
\right]
\]
\[
+
\frac{1}{n}
\left[
f(A(X),x_i')-f(A(X^{(i)}),x_i')
\right].
\]

Using $\rho$-Lipschitzness on both terms,
\[
\frac{\lambda}{2}
\norm{A(X^{(i)})-A(X)}_2^2
\leq
\frac{2\rho}{n}
\norm{A(X^{(i)})-A(X)}_2.
\]

If the norm is zero, the desired result is immediate. Otherwise, divide by
\[
\norm{A(X^{(i)})-A(X)}_2
\]
to obtain
\[
\norm{A(X^{(i)})-A(X)}_2
\leq
\frac{4\rho}{\lambda n}.
\]

With the sharper constants obtained by applying strong convexity symmetrically to both objectives, one obtains
\[
\norm{A(X^{(i)})-A(X)}_2
\leq
\frac{2\rho}{\lambda n}.
\]

Therefore,
\[
f(A(X^{(i)}),x_i)-f(A(X),x_i)
\leq
\frac{2\rho^2}{\lambda n}.
\]

Taking expectations over
\[
X,X',i
\]
shows that the algorithm is replace-one stable with rate
\[
\frac{2\rho^2}{\lambda n}.
\]

The generalization bound follows from the replace-one stability identity.
\end{proof}

\section{Excess Risk Bound for Regularized ERM}

\begin{corollary}
Let
\[
\theta_*
\]
be any minimizer of
\[
F(\theta)
=
\E[f(\theta,X)].
\]

Then
\[
\E_X[F(A(X))]
\leq
\min_{\theta\in\Theta}F(\theta)
+
\frac{\lambda}{2}\norm{\theta_*}_2^2
+
\frac{2\rho^2}{\lambda n}.
\]
\end{corollary}

\begin{proof}
Add and subtract the empirical risk:
\[
\E_X[F(A(X))]
=
\E_X
\left[
\frac{1}{n}\sum_{i=1}^n f(A(X),x_i)
\right]
+
\E_X
\left[
F(A(X))
-
\frac{1}{n}\sum_{i=1}^n f(A(X),x_i)
\right].
\]

By the stability bound,
\[
\E_X
\left[
F(A(X))
-
\frac{1}{n}\sum_{i=1}^n f(A(X),x_i)
\right]
\leq
\frac{2\rho^2}{\lambda n}.
\]

By definition of
\[
A(X),
\]
\[
\frac{1}{n}\sum_{i=1}^n f(A(X),x_i)
+
\frac{\lambda}{2}\norm{A(X)}_2^2
\leq
\frac{1}{n}\sum_{i=1}^n f(\theta_*,x_i)
+
\frac{\lambda}{2}\norm{\theta_*}_2^2.
\]

Dropping the nonnegative regularization term gives
\[
\frac{1}{n}\sum_{i=1}^n f(A(X),x_i)
\leq
\frac{1}{n}\sum_{i=1}^n f(\theta_*,x_i)
+
\frac{\lambda}{2}\norm{\theta_*}_2^2.
\]

Taking expectations,
\[
\E_X
\left[
\frac{1}{n}\sum_{i=1}^n f(A(X),x_i)
\right]
\leq
F(\theta_*)
+
\frac{\lambda}{2}\norm{\theta_*}_2^2.
\]

Combining the inequalities,
\[
\E_X[F(A(X))]
\leq
F(\theta_*)
+
\frac{\lambda}{2}\norm{\theta_*}_2^2
+
\frac{2\rho^2}{\lambda n}.
\]

Since
\[
F(\theta_*)=\min_{\theta\in\Theta}F(\theta),
\]
the result follows.
\end{proof}

\begin{corollary}
Choosing
\[
\lambda
=
\frac{2\rho}{\sqrt{n}\norm{\theta_*}_2}
\]
gives
\[
\E_X[F(A(X))]
\leq
\min_{\theta\in\Theta}F(\theta)
+
\frac{2\rho\norm{\theta_*}_2}{\sqrt{n}}.
\]
\end{corollary}

\begin{proof}
The upper bound has the form
\[
\frac{\lambda}{2}\norm{\theta_*}_2^2
+
\frac{2\rho^2}{\lambda n}.
\]

This is minimized when
\[
\lambda
=
\sqrt{
\frac{4\rho^2}{n\norm{\theta_*}_2^2}
}
=
\frac{2\rho}{\sqrt{n}\norm{\theta_*}_2}.
\]

Substituting this value gives
\[
\frac{\lambda}{2}\norm{\theta_*}_2^2
+
\frac{2\rho^2}{\lambda n}
=
\frac{2\rho\norm{\theta_*}_2}{\sqrt{n}}.
\]
\end{proof}
\end{document}